\section{Introduction}

This paper is concerned with the study of the following system of PDEs, hereafter referred to as
the \emph{Schrödinger-wave equation}
\begin{subequations}
\begin{alignat}{11}
\label{schro}&\ds \Bigl(i\partial_{t}u+\frac{1}{2}\Delta_{x}u\Bigr)(t,x)=\left(\ds\int_{\mathbb R^d\times\mathbb R^n}
\sigma_{1}(x-y) \sigma_{2}(z)\psi(t,y,z)\ud y\ud z\right)u(t,x),
\\
& \nonumber \hspace{8cm}t\in\R,\;x\in\R^{d},\\
\label{wave_schro}&(\partial_{tt}^{2}\psi-c^{2}\Delta_{z}\psi)(t,x,z)=-c^{2}\sigma_{2}(z)\biggl(\ds\int_{\mathbb R^d} \sigma_{1}(x-y)|u(t,y)|^{2}
\ud y\biggr),
\\
& \nonumber\hspace{8cm}t\in\R,\;x\in\R^{d},\;z\in\R^{n},
\end{alignat}
\end{subequations}
endowed with the initial data
\begin{equation}\label{CI_schro}
u(0,x)=u_{0}(x), \qquad (\psi(0,x,z),\partial_{t}\psi(0,x,z))=(\psi_{0}(x,z),\psi_{1}(x,z)).
\end{equation}
Here $u$ represents the wave function of a quantum particle, which interacts with the vibrational field $\psi$, and $c>0$ is a fixed parameter.
A key feature of the model is the fact that the particle motion holds in the space $\mathbb R^d$, but the
vibrations hold in a \emph{transverse direction} $\mathbb R^n$.
We are mainly interested in finding particular \emph{solitary wave} solutions of the system, with the specific form
\begin{equation}
\label{SolW}
u(t,x)=e^{i\omega t} Q(x),\qquad \psi(t,x,z)=\Psi(x,z),
\end{equation}
where $\omega\in \mathbb R$, and $Q,\Psi$ are real valued, and to investigate the stability of such solutions.

\subsection{Motivation}\label{motiv}

This work is motivated by the modeling of dissipative systems.
As suggested by A.\,Caldeira and A.\,Legget \cite{CL} the dissipation arising on a physical system might come from a coupling with a complex environment.
In this approach, dissipation is interpreted as the transfer of energy from the single
degree of freedom characterizing the system to the more
complex set of degrees of freedom describing the
environment; the energy is then evacuated into the environment and does not come back to the system.
There are many possible descriptions of the environment:
the case in which the environmental variables
are vibrational degrees of freedom is particularly appealing.
The system \eqref{schro}--\eqref{wave_schro} belongs to this class of models.
More generally, the questions we address are reminiscent to the
analysis of ``open systems'', which typically take the form of Hamiltonian systems
where the momentum and energy exchanges between a subsystem (describing a particle, say)
and the environment (an~energy reservoir, say) are expected to lead
to equilibration phenomena for the subsystem.
The issue is therefore to understand how the interactions produce an energy dissipation for the subsystem. These questions have been addressed for a large variety of classical and quantum couplings, appealing to a large panel of mathematical arguments,
ranging from ergodic theory to spectral analysis, PDE analysis, probability, and numerics,
\cite{bach,JP1,JP2,KKS,KKSb,bDPL,SW}.

The system \eqref{schro}--\eqref{wave_schro} is nothing but a quantum version of a model introduced by
L.\,Bruneau and S.\,de\,Bièvre in \cite{BdB} for describing
a classical particle interacting with its environment seen as a bath of oscillators. Roughly speaking in each space position $x\in\R^{d}$
there is a membrane oscillating on a transverse direction $z\in\R^{n}$. When the particle hits a membrane, its kinetic energy activates
vibrations and the energy is evacuated at infinity in the $\mathbb R^n$ directions.
In particular, the coordinates $(z_1,\dots,z_n)\in \mathbb R^n$ need not have the specific
dimension of a length (but adopting this language might definitely help the intuition).
These energy transfer mechanisms eventually act as a sort of friction force
on the particle, an intuition rigorously justified in
\cite[Th.\,2 \& Th.\,4]{BdB}.
The system for the position of the particle $t\mto q(t)$ and the state of the vibrational environment $(t,z)\mto \psi(t,z)$ reads
\begin{subequations}
\begin{alignat}{11}
\label{part}&\ds \overbigddot{q}(t)=-
\ds\int \nabla \sigma_1(q(t)-y) \sigma_2(z) \psi(t,y,z)\ud z\ud y,
\qquad&
t\in\R,\\
\label{wave_part}&\ds (\partial_{tt}^{2}\psi-c^{2}\Delta_{z}\psi)(t,z)=-\sigma_{2}(z)\,\sigma_{1}(x-q(t)),\qquad&
t\in\R,\;x\in\R^{d},\;z\in\R^{n},
\end{alignat}
\end{subequations}
completed by the initial data
\begin{equation}\label{CI_part}
(q(0),\overbigdot{q}(0))=(q_{0},p_{0}),\qquad (\psi(0,x,z),\partial_{t}\psi(0,x,z))=(\psi_{0}(x,z),\psi_{1}(x,z)).
\end{equation}
The functions $\sigma_{1}:\mathbb R^d\to [0,\infty)$ and $\sigma_{2}:\mathbb R^n\to [0,\infty)$ are form functions encoding the interaction domain between the particle and the environment.
The model can be extended by considering $P$-interacting particles,
and the mean-field regime $P\to \infty$ leads to the following Vlasov-wave system
\cite{GV}
\begin{subequations}
\begin{alignat}{11}
\label{kin}&\ds \partial_{t}f+v\cdot\nabla_{x}f-\nabla_{x}\left(\sigma_{1}\star_{x}\int\sigma_{2}\psi\ud z\right)\cdot\nabla_{v}f=0, & \hspace{-2.7cm}t\in\R,\;x\in\R^{d},\;v\in\R^{d},\\
\label{wave_kin}&\ds \partial_{tt}^{2}\psi-c^{2}\Delta_{z}\psi=-\sigma_{2}(z)\biggl(\sigma_{1}\star_{x}\int f\ud v\biggr), & \hspace{-2.7cm}t\in\R,\;x\in\R^{d},\;z\in\R^{n},\\
\label{CI_kin}
&f(0,x,v)=f_{0}(x,v),\qquad (\psi(0,x,z),\partial_{t}\psi(0,x,z))=(\psi_{0}(x,z),\psi_{1}(x,z)),
\end{alignat}
\end{subequations}
where $f$ stands for the particle distribution function in phase space.
This system is thoroughly
investigated in \cite{AGV,BGV,AVPhD}.
In \cite{dBGV}, it is proposed to rescale the wave equation \eqref{wave_kin} as follows
\begin{equation}\label{wave_kin2}
\partial_{tt}^{2}\psi-c^{2}\Delta_{z}\psi=-c^{2}\sigma_{2}\biggl(\sigma_{1}\star_{x}\int f\ud v\biggr).
\end{equation}
As $c$ goes to $+\infty$, the solutions of the rescaled system \eqref{kin}, \eqref{wave_kin2} tend to solutions of
\begin{subequations}
\begin{alignat}{11}
\label{Vlasov}&\ds \partial_{t}\tilde f+v\cdot\nabla_{x}\tilde f-\nabla_{x}\left(\sigma_{1}\star_{x}\int\sigma_{2}\tilde \psi \ud z\right)\cdot\nabla_{v}\tilde f=0,&\hspace{.5cm}t\in\R,\;x\in\R^{d},\;v\in\R^{d},\\
\label{Poissonz_Vlasov}&\ds -\Delta_{z}\tilde \psi=-\sigma_{2}\biggl(\sigma_{1}\star_{x}\int \tilde f\ud v\biggr),&\hspace{.5cm}t\in\R,\;x\in\R^{d},\;z\in\R^{n},
\end{alignat}
\end{subequations}
(Without the rescaling the regime $c\to \infty$ would simply lead to the free transport equation for the particle distribution function $\tilde f$.)
We can write
\[
\tilde \psi(t,x,z)=\Gamma(z)\,\biggl(\sigma_{1}\star\int \tilde f\ud v\biggr)(x),
\]
where $\Gamma$ denotes the unique solution of
\begin{equation}\label{eqGamma}
-\Delta_{z}\Gamma=-\sigma_{2}, \qquad \Gamma \in H^{1}(\R_{z}^{n}).
\end{equation}
This observation allows us to
express \eqref{Vlasov}--\eqref{Poissonz_Vlasov} as a standard Vlasov equation\vspace*{-3pt}
\begin{equation}\label{Vlasov_bis}
\partial_{t}\tilde f +v\cdot\nabla_{x}\tilde f+\kappa\nabla_{x}\left(\Sigma\star_{x}\int \tilde f\ud v\right)\cdot\nabla_{v}\tilde f=0,\hspace{1cm}t\in\R,\;x\in\R^{d},\;v\in\R^{d},
\end{equation}
where the potential is defined by a convolution with the macroscopic density, with
\begin{equation}\label{defkapp}\kappa=\|\nabla_{z}\Gamma\|_{L_{z}^{2}}^{2},\qquad \Sigma=\sigma_{1}\star\sigma_{1}.
\end{equation}
Quite surprisingly --- mind the sign $\kappa>0$ --- this corresponds to an attractive dynamics.
This unexpected connection guides the intuition to establish further features
of the solutions of the Vlasov-wave system; it particular, they
exhibit Landau damping phenomena \cite{LVNum1,GV_LD}.
The analysis of these models, either for a single particle or the kinetic description, brings out the critical role of the wave speed $c>0$ and the
dimension $n$ of the space for the wave equation.

The system \eqref{schro}--\eqref{wave_schro} then appears as the quantum version of the L.\,Bruneau and S.\,de\,Bièvre model.
This intuition can be justified by the semi-classical analysis {\it à la } P.-L.\,Lions-T.\,Paul \cite{Lions_Paul_Wigner}, which
makes a natural connection between the Vlasov-wave system and \eqref{schro}--\eqref{wave_schro}, see Appendix~\ref{SClassicalAnalysis} and \cite{LV_PhD}.
Note that here we have adopted from the beginning the rescaling where the coupling
term
in the wave equation \eqref{wave_schro} is of the order of $c^2$. We will motivate this
choice below.
According to the framework
introduced in \cite{BdB}, throughout this article we assume:
\begin{enumerate}\renewcommand{\theenumi}{H\arabic{enumi}}
\item\label{h1}
$n\geq 3$,
\item\label{h2}
The form functions $\sigma_{1}$ and $\sigma_{2}$ are non-negative, smooth, compactly supported and radially symmetric.
\end{enumerate}
As said above the role of the dimension $n$ for the wave equation is critical in these models.
Indeed, the evacuation of energy in the environment relies on the dispersion properties of the wave equation, which are strong enough when $n$ is sufficiently large~\cite{GV_LD}.
By the way, notice that the definition of $\kappa$ in \eqref{defkapp} makes sense when assuming $n\geq 3$.
The case $n=3$ also plays a specific role in the theory presented in \cite{BdB}.
The assumptions \eqref{h1} and \eqref{h2} on the form functions are very natural in the modeling framework of \cite{BdB}.
In what follows, we use the abuse of notation
to mix up a radially symmetric function of $x\in \mathbb R^d$ with the
underlying function of the scalar quantity $|x|$, and we will equally refer to the monotonicity
of this function.
It is likely that the compactness assumption on the support of the
functions $\sigma_1, \sigma_2$ can be relaxed into a fast enough decay property; radial symmetry is however crucial in the present work.
Following the observations made for classical particles, it is instructive to consider the regime
where $c$ goes to $+\infty$ in \eqref{schro}--\eqref{wave_schro}.
We are led to\vspace*{-3pt}
\begin{subequations}
\begin{alignat}{11}
\label{Hartree}&\ds i\partial_{t}\tilde u+\frac{1}{2}\Delta_{x}\tilde u =\left(\sigma_{1}\star_{x}\int\sigma_{2}\tilde \psi\ud z\right)\tilde u,&\hspace{1cm}t\in\R,\;x\in\R^{d},\\
\label{Poissonz_Hartree}&\ds -\Delta_{z}\tilde \psi=-\sigma_{2}(z)\left(\sigma_{1}\star_{x}|\tilde u|^{2}\right)(x),&\hspace{1cm}t\in\R,\;x\in\R^{d},\;z\in\R^{n},
\end{alignat}
\end{subequations}
which can be cast in the usual form of an Hartree type equation
\begin{equation}\label{hartree_bis}
i\partial_{t}\tilde u+\frac{1}{2}\Delta_{x}\tilde u=-\kappa\left(\Sigma\star_{x}|\tilde u|^{2}\right)\tilde u,\hspace{1cm}t\in\R,\;x\in\R^{d}.
\end{equation}
This remark will be helpful for the analysis.

The conservation of the total
energy
is a remarkable property of all these models.
For the particle equation \eqref{part}--\eqref{wave_part}, we set
\begin{align*}
\mathcal E_{\mathrm{part}}(t)=\ds\frac{|\overbigdot q(t)|^2}{2} &+\frac12\ds\int \big(
|\partial_t \psi | ^ 2+ c^2 |\nabla_z \psi|^2\big) (t,x,z) \ud z\ud x
\\
&+\int \sigma_1( q(t)-y)\sigma_2(z)\psi(t,y,z)\ud y \ud z
\end{align*}
and for the kinetic equation \eqref{kin}, with \eqref{wave_kin2}
(mind the rescaling for the wave equation), we set
\begin{align*}
\mathcal E_{\mathrm{kin}}(t)
=\frac{1}{2} \ds\int v^2 f(t,x,v)\ud v \ud x&+\frac{1}{2}\ds\int \Big(\frac{|\partial_t \psi | ^ 2}{c^2}+ |\nabla_z \psi|^2\Big) (t,x,z) \ud z\ud x
\\
&+\int \sigma_1( x-y)\sigma_2(z)\psi(t,y,z)f(t,x,v) \ud v \ud x \ud y \ud z.
\end{align*}
Then, we have
\[
\mathcal E_{\mathrm{part}}(t)=\mathcal E_{\mathrm{part}}(0),\qquad
\mathcal E_{\mathrm{kin}}(t)=\mathcal E_{\mathrm{kin}}(0).
\]
For the quantum model, \eqref{schro}--\eqref{wave_schro}, it becomes
\begin{equation}\label{energy_Sch}
\begin{array}{lll}
\mathcal E_{\mathrm {Schr}}(t)&=&
\ds\frac{1}{2}
\ds\int |\nabla_{x}u(t,x)|^2 \ud x+
\ds\frac{1}{2}\ds\int \Big(
\ds\frac{|\partial_t \psi | ^ 2}{c^2}+ |\nabla_z \psi|^2\Big) (t,x,z) \ud z\ud x
\\&&\hspace*{2.2cm} +
\ds \int
\sigma_{1} (x-y) \sigma_{2}(z)\psi(t,y,z) |u(t,x) |^{2}\ud z\ud y\ud x
\\&=&\mathcal E_{\mathrm {Schr}}(0).
\end{array}
\end{equation}
For
the asymptotic Hartree equation \eqref{hartree_bis}, we get
similarly
\begin{equation}\label{energy_Har}
\mathcal H(t)=\frac{1}{2}\ds\int |\nabla_{x}\tilde u(t,x)|^2 \ud x
-\frac{\kappa}{2}\ds\int \Sigma(x-y) |\tilde u(t,y) |^{2}|\tilde u(t,x) |^{2}\ud x\ud y
=\mathcal H(0).
\end{equation}
Moreover, both quantum equations
are invariant by translation and phase and they preserve the mass of the wave function:
\begin{equation}\label{mass_conservation}
\ds\int |u(t,x)|^2 \ud x=
\ds\int |u(0,x)|^2 \ud x,
\qquad
\ds\int |\tilde u(t,x)|^2 \ud x= \ds\int |\tilde u(0,x)|^2 \ud x,
\end{equation}
constant quantities that we denote $\mathscr{M}_0$ and $\tilde{\mathscr{M}}_0$.
However, there are fundamental
differences between the two equations.
Let
\[
p(t)=\Im \int\nabla_{x}u(t,x)\bar{u}(t,x)\ud x,\qquad
\tilde p(t)=\Im \int\nabla_{x}\tilde u (t,x)\overline{\tilde u}(t,x)\ud x
\]
be the momentum associated to \eqref{schro}--\eqref{wave_schro} and \eqref{hartree_bis}, respectively.
We have, for~\eqref{hartree_bis},
\[
\ds\frac{\ud}{\ud t }\, \tilde p=0,
\]
but
\[
\ds\frac{\ud}{\ud t }\,p(t)=-\int_{\R^{d}}\nabla_{x}\left(\sigma_{1}\star\int_{\R^{n}}\sigma_{2}(z)\psi(t,x,z)\ud z\right)|u(t,x)|^{2}\ud x
\]
for \eqref{schro}--\eqref{wave_schro}.
We also introduce the center of mass
\[
q(t)=\frac{\ds\int_{\R^{d}}x\,|u(t,x)|^{2}\ud x}{\ds\int_{\R^{d}}|u(t,x)|^{2}\ud x}=\frac{1}{\mathscr{M}_0}\int_{\R^{d}}x\,|u(t,x)|^{2}\ud x
\]
associated to \eqref{schro}--\eqref{wave_schro} and a similar definition $\tilde q(t)$ for \eqref{hartree_bis}.
We have
\[
\mathscr{M}_0\ds\frac{\ud}{\ud t }\, q(t)=p(t),\qquad \tilde{\mathscr{M}}_0\ds\frac{\ud}{\ud t }\,\tilde q(t)=\tilde p(t).
\]
Therefore,
the momentum conservation for \eqref{hartree_bis} implies that the center of mass follows a straight line at constant speed.
For \eqref{schro}--\eqref{wave_schro}, the analogy
with the case of a single classical particle would lead to conjecture that the center of mass will stop exponentially fast.
Numerical experiments shed some light on this issue \cite{LVNum2}.
Finally, we note that \eqref{hartree_bis} is also Galilean invariant: if $\tilde u$ is a solution of \eqref{hartree_bis}, then $v(t,x)=\tilde u(t,x-tp_{0})e^{ip_{0}\cdot(x-t\sfrac{p_{0}}{2})}$ still is a solution of \eqref{hartree_bis}.
This property
is not fulfilled by the
system \eqref{schro}--\eqref{wave_schro}, which leads to a specific behavior of the solutions, consistently with the previous remark.

Note that we find convenient to perform the analysis on a dimensionless version of the PDE system (even if we are using terminology like ``mass'', ``position'', ``speed'' that helps the intuition).
The identification
of the dimensionless parameters
is detailed in Appendix~\ref{adim}. In particular, in what follows the ``mass''
\eqref{mass_conservation}
is not necessarily normalized to 1.
The viewpoint has the advantage of making analogies
appear more clearly with the classical system \eqref{part}--\eqref{wave_part}.
It turns out
that the stability issue relies on the interaction between $c$, the mass of the initial data, and the amplitude of the perturbation, see \cite{LVNum2,LV_PhD}
for further numerical evidence.

\subsection{Scaling properties}\label{scaling}

It is well-known that scaling invariance plays a central role in the analysis of nonlinear Schrödinger equations.
Here, let $(u,\psi)$ be a solution of \eqref{schro}--\eqref{wave_schro}, and, for given $\lambda,\mu>0$, let us set
\[
\big(u_{\lambda,\mu}(t,x),\psi_{\lambda,\mu}(t,x,z)\big)=\big(\mu\,u(\lambda^{2}t,\lambda x),\mu\,\lambda^{n-1}\psi(\lambda^{2}t,\lambda x,\lambda^{2}z)\big).
\]
It turns out that $u_{\lambda,\mu}$ is a solution of \eqref{schro}--\eqref{wave_schro} but with the \emph{rescaled} form functions
\[
\sigma_{1}^{\lambda,\mu}(x)=\mu^{-1}\lambda^{d+1}\sigma_{1}(\lambda x) \hspace{0.2cm}\text{ and }\hspace{0.2cm} \sigma_{2}^{\lambda}(z)=\lambda^{n+2}\sigma_{2}(\lambda^{2}z).
\]
Since $\sigma_{1}$ and $\sigma_{2}$ are not homogeneous functions, $(u,\psi)$ and $(u_{\lambda,\mu},\psi_{\lambda,\mu})$ are solutions of the \emph{same} Schrödinger-wave system if and only if $\lambda=1=\mu$. The same conclusion applies to the limiting system: if $\tilde{u}$ is a solution of \eqref{hartree_bis} then $\tilde{u}_{\lambda,\mu}(t,x)=\mu\tilde{u}(\lambda^{2}t,\lambda x)$ is a solution of \eqref{hartree_bis} with the rescaled potential $\Sigma^{\lambda,\mu}(x)=\sigma_{1}^{\lambda,\mu}\star\nobreak\sigma_{1}^{\lambda,\mu}(x)=\mu^{-2}\lambda^{d+2}\Sigma(\lambda x)$.
Therefore, in contrast to the usual nonlinear Schrödinger or Hartree equations, we cannot find
a relation between $\lambda$ and~$\mu$ such that the $(u_{\lambda,\mu},\psi_{\lambda,\mu})$'s are solutions of the \emph{same} equation than $(u,\psi)$;
this lack of scale invariance will have an important role in the sequel of this paper.

Nevertheless, the scaling property implies that any result valid for the Hartree equation with a given potential $\Sigma$ equally applies
to the equations with the modified potentials $\Sigma^{\lambda,\mu}$. Considering the case where $\lambda=\mu=\epsilon^{-1}$ and letting
$\epsilon$ go to $ 0$, up to a suitable renormalization,
allows us to consider the regime $\Sigma\to\delta_{0}$ which formally leads to the standard cubic nonlinear Schrödinger equation
\begin{equation}\label{nls}
i\partial_{t} U+\frac{1}{2}\,\Delta_{x}U=-\kappa|U|^{2}U.
\end{equation}
This equation is $L^{2}$-subcritical in the case $d=1$, it is $L^{2}$-critical in the case $d=2$ and $L^{2}$-super-critical in the case $d\geq 3$. Hence, this formal limit suggests different behaviors for the Hartree equation (when a smooth potential is considered), depending on the dimension $d$.
Even if the continuity with respect to $\Sigma$ as $\Sigma\to\delta_{0}$ is certainly wrong when $d\geq 2$ --- \eqref{nls} admits solutions which blow up in finite time while solutions of \eqref{hartree_bis} are globally defined when $\Sigma$ is smooth --- our analysis shows several differences between the case $d=1$ and $d\geq 2$, which can be understood from the formal asymptotic to \eqref{nls}. It is thus not surprising that our main results, Theorem~\ref{orbital_stab_2} and Proposition~\ref{non_perturbative_L+}, require some additional assumptions on the form function $\sigma_{1}$. Namely, in the case $d=3$, we shall
consider $\Sigma=\sigma_{1}\star\sigma_{1}$ such that
the rescaled potentials $\Sigma^{\lambda,\mu}$, with $\lambda,\mu>0$, are \emph{close enough} to $|\cdot|^{-1}$ (note that when $d=3$ and $\Sigma=|\cdot|^{-1}$, the Hartree equation is $L^{2}$-subcritical). When $d=1$ we do not require any additional assumption on $\sigma_{1}$: see Proposition~\ref{non_perturbative_L+_d=1},
obtained precisely by using the $L^2$ subcritical feature of \eqref{nls} when $d=1$.

\subsection{Solitary waves}

The system \eqref{schro}--\eqref{wave_schro} can be shown to be well-posed,
in~natural functional spaces associated to the energy conservation.

\begin{theo}\label{TC}
Let \eqref{h1}--\eqref{h2} be fulfilled. For all initial data $u_{0}\in H^{1}(\R_{x}^{d})$, $\psi_{0}\in\nobreak L^{2}(\R_{x}^{d};\overbigdot{H}^{1}(\R_{z}^{n}))$ and $\psi_{1}\in L^{2}(\R_{x}^{d};L^{2}(\R_{z}^{n}))$, the system \eqref{schro}--\eqref{wave_schro} and \eqref{CI_schro} admits a unique global solution $(u,\psi)$ such that $u\in C^{0}([0,+\infty);H^{1}(\R_{x}^{d}))$ and
\[
\psi\in C^{0}\bigl([0,+\infty);L^{2}\bigl(\R_{x}^{d};\overbigdot{H}^{1}(\R_{z}^{n})\bigr)\bigr)\;\cap\;C^{1}\bigl([0,+\infty);L^{2}\left(\R_{x}^{d};L^{2}(\R_{z}^{n})\right)\bigr).
\]
\end{theo}

The proof is detailed in Appendix~\ref{CauchyTheory}. The local well-posedness is based on Stri\-chartz' estimates, which rely on the dispersive properties of
the Schrödinger and the wave equations in the coupling.
The difficulty comes from the fact that Strichartz' estimates for \eqref{schro} lead to estimates of $u$ in $L_{t}^{q}L_{x}^{r}$ norms whereas Strichartz' estimates for \eqref{wave_schro} lead to estimates on $\psi$ in $L_{x}^{r}L_{t}^{q}L_{z}^{p}$ norms. Then, in order to gather these estimates, it is necessary
to deal with
permutations of Lebesgue-norms in time and space. For this purpose, assumption \eqref{h2} allows us to apply Hölder and Young inequalities in order to always obtain estimates in $L_{t}^{q}L_{x}^{q}$-norms. Eventually, that solutions are globally defined comes from the Hamiltonian structure of the system.

The main purpose of this article is to show the \emph{existence and the orbital stability of solitary waves} for the Schrödinger-wave system.
Namely, we are going to study solutions of \eqref{schro}--\eqref{wave_schro} with the form \eqref{SolW}.
The existence of such non dispersive solutions is the translation of the presence of some attractive dynamics induced by the model.
The rescaling \eqref{wave_kin2} is important in the discussion.
We start by observing that if~\hbox{$(u,\psi)=(Q(x)e^{i\omega t},\Psi(x,z))$} is a solution of \eqref{schro}--\eqref{wave_schro}, then $(Q,\Psi)$ is a solution~of
\begin{subequations}
\begin{alignat}{11}
\label{soliton}&\ds -\frac{1}{2}\Delta_{x}Q+\omega Q+\left(\sigma_{1}\star_{x}\int\sigma_{2}\Psi\ud z\right)Q=0,&\hspace{1cm}x\in\R^{d},\\
\label{wave_soliton}&-c^{2}\Delta_{z}\Psi=-c^{2}\sigma_{2}(z)\left(\sigma_{1}\star_{x}Q^{2}\right)(x),&\hspace{1cm}x\in\R^{d},\;z\in\R^{n},
\end{alignat}
\end{subequations}
which is in fact independent of the parameter $c$. In turn, the profiles $(Q,\Psi)$ do not depend on $c$.
Moreover these particular solutions $(Q(x)e^{i\omega t},\Psi(x,z))$ are also solutions of the asymptotic system \eqref{Hartree}--\eqref{Poissonz_Hartree}.
It is therefore relevant to compare the behavior of the solutions of \eqref{schro}--\eqref{wave_schro}
and the solutions of \eqref{Hartree}--\eqref{Poissonz_Hartree} around the state $(Q(x)e^{i\omega t},\Psi(x,z))$:
this comparison
provides information on the action of the environment on the quantum particle.

According to the previous discussion, the expected behavior for the Schrödinger-wave system can be summarized as follows.
\begin{conj}
Let $(Q,\Psi)$ be a solution of \eqref{soliton}--\eqref{wave_soliton}.
If $u_{0}(x)=Q(x)e^{i\sfrac{p_{0}\cdot x}{M}}$ for some sufficiently small $p_{0}$ and if $(\psi_{0},\psi_{1})=(\Psi,0)$, then there exists two functions $x=x(t)$ and $\gamma=\gamma(t)$ such that
\begin{itemize}
\item the unique solution $(u,\psi)$ of \eqref{schro}--\eqref{wave_schro} associated to these initial conditions remains close (uniformly in time in some norms that have to be made precise) to $(Q(\cdot-x(t))e^{i\gamma(t)},\Psi(\cdot-x(t),\cdot))$;
\item $|\overbigdot{x}(t)|\leq C e^{-\lambda\sfrac{t}{c}}$ and $|x(t)-\bar{x}|\leq C e^{-\lambda\sfrac{t}{c}}$.
\end{itemize}
\end{conj}

In fact, we expect that the stability applies for particular solutions $(Q,\Psi)$, having minimal energy, the \emph{ground states} solutions of \eqref{soliton}--\eqref{wave_soliton}
The perturbation $e^{i\sfrac{p_{0}\cdot x}{M}}$ is particularly relevant because
it corresponds to provide an impulsion $p_0$ to the solitary wave
$Q$ and it gives rise
to a simple solution for the asymptotic model \eqref{hartree_bis}:
the solitary wave moves on a straight
line with a uniform momentum
since
\[
\tilde u(t,x)=Q(x-tp_0/M)e^{i(x-tp_0/M)\cdot p_0/M}e^{it(\omega +p_0^2/(2M^2))}
\]
is a solution of \eqref{hartree_bis}.
Hence it can be used as a reference state to understand the dynamics of
\eqref{schro}--\eqref{wave_schro}, which is expected to induce
quite intricate deformations of the solitary wave. This issue is further discussed
on numerical grounds in \cite{LVNum1}.

The orbital stability of solitary waves of nonlinear Schrödinger equations is a classical result for many years, see for instance \cite{Cazenave_Lions_orbital_stab,Weinstein1,Weinstein2}.
A general theory for abstract Hamiltonian systems $\partial_t u=\mathscr J \mathscr E'(u)$
in infinite dimension, where $\mathscr J$ is a skew-symmetric operator and $\mathscr E$ an ``energy'' functional
has been developed in \cite{GSS,GSS2}, and recently revisited in \cite{SRN1}.
However, there are several difficulties to justify the orbital stability in the present context, beyond the identification of $\mathscr J$ and $\mathscr E$, and these arguments cannot be applied directly.
The difficulties are
related to the fact that the nonlinearity is non local with, moreover,
the fact that the particle-wave unknown has a vectorial structure, so that the analysis does not reduce to the framework discussed in \cite{GSS,GSS2}.
In particular,
the ``Schrödinger part'' and the ``wave part'' of the system act on different variables ($x$ and $z$)
and this makes the functional framework, the spectral analysis and the group action much more intricate.
These issues are further discussed in the simpler context of plane waves solutions in \cite{GRN}.

Nevertheless, we can expect that structure properties of the simpler problem \eqref{hartree_bis} still apply to the system \eqref{schro}--\eqref{wave_schro}.
At first sight, assumption \eqref{h2} can be expected to make the problem easier than the
case where $\Sigma$ is replaced by the
kernel of the Poisson equation in dimension $d=3$, that is $\Sigma^0(x)=\sfrac{1}{|x|}$.
This specific case \eqref{hartree_bis} --- the Schrödinger-Newton equation --- has been investigated in detail by E.\,Lenzmann~\cite{Lenzmann}.
However, as reported above, while $\Sigma=\sigma_1\star\sigma_1$ has better regularity and support properties, it does not satisfy any
scale invariance.
It turns out that the analysis
of the Schrödinger-Newton equation exploits, in a quite crucial way, either explicit formulas or the scale invariance which are very specific to the kernel $\sfrac{1}{|x|}$.
For this reason, for establishing the first part of the conjecture (see Theorem~\ref{orbital_stab_2}), we shall use a quite indirect approach,
that relies on the perturbative arguments developed in \cite{Lenzmann}
for establishing spectral properties for the non relativistic Hartree equation.
The second part of the conjecture justifies that the environment acts on the quantum particle as a friction force and is the object of further investigations
on numerical grounds \cite{LVNum2,LV_PhD}.

\subsubsection*{Acknowledgements}
We are gratefully indebted to Stephan De Bièvre for many motivating discussions and warm encouragements. We also thank Enno Lenzmann for useful hints and kind advices. Finally, we thank David Chiron who indicated relevant improvements of the arguments.

\section{Main results}

As said above, the main objective is to discuss the existence and the stability of non trivial solutions (with finite mass and energy) of \eqref{schro}--\eqref{wave_schro}
with the form~\eqref{SolW}.
In~order to establish the existence, we start by observing that $(Q,\Psi)$ has to be a solution of \eqref{soliton}--\eqref{wave_soliton}. Then we can express $\Psi$ in terms of $Q$ as follows:\vspace*{-3pt}
\[
\Psi(x,z)=\Gamma(z)\,\sigma_{1}\star Q^{2}(x),
\]
where $\Gamma$ stands for the unique solution of \eqref{eqGamma}.
Coming back to \eqref{soliton}, we deduce that~$Q$ satisfies\vspace*{-3pt}
\begin{equation}\label{choquard}
-\frac{1}{2}\Delta_{x}Q+\omega Q-\kappa(\Sigma\star Q^{2})Q=0
\end{equation}
with the definition \eqref{defkapp}. This equation is known as the \emph{Choquard equation} and it has been intensively studied (see for example \cite{Lions_Choquard}, \cite{Lieb_Choquard} or \cite{Lenzmann} and the references therein). In particular, we already know from \cite{Lions_Choquard} that there exists infinitely many solitary waves.

\subsection{Ground states}

Nevertheless, we are only interested in \emph{stable} solitary waves: for this reason, we consider solitary waves that minimize the energy of the system under a mass constraint, a quantity preserved by the evolution equation. Such solitary waves are called \emph{ground states}.
The specific case of the Newtonian potential $\Sigma^0(x)=\sfrac{1}{|x|}$ in dimension $d=3$
has been studied in
\cite{Lieb_Choquard} which establishes the existence and uniqueness (up a change of phase and translation) of ground states for \eqref{hartree_bis}.
The existence part of \cite{Lieb_Choquard} still applies in the case where $\Sigma$ is a smooth, compactly supported,
radially symmetric, non increasing and non negative function.
However, the arguments for proving the uniqueness part of the statement rely strongly on the specific form of the Newtonian potential.
Besides, the definition of the energy
functional for the system \eqref{schro}--\eqref{wave_schro} differs from those of \eqref{hartree_bis}.
Therefore, one has to check that \eqref{schro}--\eqref{wave_schro} admits ground states. For that purpose we will need the following additional assumption on the form function $\sigma_{1}$.

\begin{enumerate}\renewcommand{\theenumi}{H\arabic{enumi}}\setcounter{enumi}{2}
\item\label{h3}
The form function $\sigma_{1}$ is non increasing.
\end{enumerate}

We interpret
the energy functional \eqref{energy_Sch} as depending on $u$, $\psi$ and $\chi=\partial_t\psi$. Namely, for $u:\mathbb R^d\to \mathbb C$,
$\psi, \chi: \mathbb R^d\times \mathbb R^n \to \mathbb R$, we set
\begin{multline*}
E(u,\psi,\chi)=
\ds\frac{1}{2}
\ds\int |\nabla_{x}u(x)|^2 \ud x+
\ds\frac{1}{2}\ds\int \Big(
\ds\frac{|\chi | ^ 2}{c^2}+ |\nabla_z \psi|^2\Big) (x,z) \ud z\ud x
\\
+
\ds \int
\sigma_{1} (x-y) \sigma_{2}(z)\psi(y,z) |u(x) |^{2}\ud z\ud y\ud x,
\end{multline*}
so that $\mathcal E_{ \mathrm{Sch}}(t)=E( u,\psi,\partial_t \psi)(t)$.
Similarly, we set
\begin{equation}\label{hamiltonian}
H(u)=\frac{1}{2}\ds\int |\nabla_{x}u(x)|^2 \ud x
-\frac{\kappa}{2}\ds\int \Sigma(x-y) |u(y) |^{2}|u(x) |^{2}\ud x\ud y,
\end{equation}
see \eqref{energy_Har}.
In order to establish the existence of ground states we will study the following three minimization problems.
\begin{subequations}
\begin{alignat}{1}
\label{IM}&\ds I_{M}:=\inf\bigl\{E(u,\psi,\chi)\sep(u,\psi,\chi)\in H_{x}^{1}\times L_{x}^{2}\overbigdot{H}_{z}^{1}\times L_{x}^{2}L_{z}^{2}\text{ and }\|u\|_{L_{x}^{2}}^{2}\leq M\bigr\},\\
\label{JM}&\ds J_{M}:=\inf\bigl\{E(u,\psi,\chi)\sep(u,\psi,\chi)\in H_{x}^{1}\times L_{x}^{2}\overbigdot{H}_{z}^{1}\times L_{x}^{2}L_{z}^{2}\text{ and }\|u\|_{L_{x}^{2}}^{2}=M\bigr\},\\
\label{KM}&\ds K_{M}:=\inf\bigl\{E(u,\Gamma\,\sigma_{1}\star|u|^{2},0)\sep u\in H_{x}^{1}\text{ and }\|u\|_{L_{x}^{2}}^{2}=M\bigr\}.
\end{alignat}
\end{subequations}
The interest of \eqref{KM} comes from the fact that $E(u,\Gamma\,\sigma_{1}\star|u|^{2},0)=H(u)$
since $\sigma_{1}$ is radially symmetric and therefore $\|\sigma_{1}\star|u|^{2}\|_{L_{x}^{2}}^{2}=\iint|u|^{2}(x)\Sigma(x-y)|u|^{2}(y)\ud x\ud y$. Then, if $K_{M}$ is reached at $u$,
$u$ is a ground state of \eqref{hartree_bis} and we will be able to compare ground states of \eqref{schro}--\eqref{wave_schro} with ground states of \eqref{hartree_bis}.
Section~\ref{GroundStates} is devoted to the proof of the following theorem.
\begin{theo}\label{existence_GS}
Let \eqref{h1}--\eqref{h3} be fulfilled.
\begin{enumeratei}
\item\label{existence_GSi} For every $M\geq 0$, $I_{M}$ is reached.
\item\label{existence_GSii}
For every $M\geq 0$, $I_{M}=J_{M}=K_{M}$.
\item\label{existence_GSiii}
There exists a mass threshold $M_{0}\geq 0$ such that for every $M\in[0,M_{0}]$, $J_{M}=0$ and for every $M>M_{0}$, $J_{M}<0$ is reached on $(u,\psi,\chi)=(u,\psi,0)$ with $u$ non negative, radially symmetric and non increasing. Moreover $(u,\psi)$ is a solution of \eqref{soliton}--\eqref{wave_soliton} for a certain $\omega>0$. In particular $\psi=\Gamma\,\sigma_{1}\star|u|^{2}$ is non positive, $u$ is an element of the Schwartz class $\mathcal{S}(\R^{d})$ and $K_{M}=J_{M}$ is reached at $u$.
\item\label{existence_GSiv}
If $d\geq 2$, then $M_{0}>0$.
\item\label{existence_GSv}
If $d=1$, then $M_{0}=0$.
\end{enumeratei}
\end{theo}

Note that we do not know whether the minimizer in item \eqref{existence_GSiii} is uniquely defined, up to a possible change of phase and translation.
Applying Lieb's method \cite{Lieb_Choquard}, we cannot even conclude whether or not the minimizers of $J_{M}$ are
radially symmetric, a preliminary step
to establish uniqueness, and strictly positive.
The alternative approach of L.\,Ma and L.\,Zhao \cite[\S 5]{MaZ} provides a positive answer to the strict positivity and radial symmetry of the minimizer, though.
Note also that the fourth item of this theorem is reminiscent to the fact that \eqref{schro}--\eqref{wave_schro} does not have a scale invariance. In contrast,
$M_{0}=0$ when $d=1$, a difference with the cases $d\geq 2$ which can be related with the discussion in Section~\ref{scaling}.

\subsection{Orbital stability}

The variational characterization will be used in Section~\ref{OrbitalStability1} to establish the following orbital stability result for these ground states. In this statement, for a given mass $M>0$, we denote by $S_{M}$ the space of all possible ground states
\[
S_{M}=\bigl\{(\widetilde{Q},\widetilde{\Psi})\in H_{x}^{1}\times L_{x}^{2}\overbigdot{H}_{z}^{1}\mid\|\widetilde{Q}\|_{L_{x}^{2}}^{2}=M\text{ and }E(\widetilde{Q},\widetilde{\Psi},0)=J_{M}\bigr\}.
\]

\begin{theo}\label{orbital_stab_1}
Let \eqref{h1}--\eqref{h3} be fulfilled. Let $M>M_{0}$ and $(Q,\Psi)$ be in $S_{M}$. For every $\varepsilon>0$ there exists $\delta_{\varepsilon}>0$ such that if $u_{0}\in H_{x}^{1}$, $\psi_{0}\in L_{x}^{2}\overbigdot{H}_{z}^{1}$ and $\chi_{0}\in L_{x}^{2}L_{z}^{2}$ with
\[
\|u_{0}-Q\|_{H_{x}^{1}}^{2}+\|\psi_{0}-\Psi\|_{L_{x}^{2}\overbigdot{H}_{z}^{1}}^{2}+\|\chi_{0}\|_{L_{x}^{2}L_{z}^{2}}^{2}<\delta_{\varepsilon},
\]
then the unique solution $(u,\psi,\chi=\partial_t\psi)$ of \eqref{schro}--\eqref{wave_schro} with initial data $(u_{0},\psi_{0},\chi_{0})$ satisfies
\[
\underset{t\geq 0}{\sup}\;\underset{(\widetilde{Q},\widetilde{\Psi})\in S_{M}}{\inf}\left(\|u(t)-\widetilde{Q}\|_{H_{x}^{1}}^{2}+\|\psi(t)-\widetilde{\Psi}\|_{L_{x}^{2}\overbigdot{H}_{z}^{1}}^{2}+\|\chi(t)\|_{L_{x}^{2}L_{z}^{2}}^{2}\right)<\varepsilon.
\]
\end{theo}
The proof is classical and based on the concentration-compactness lemma, see for instance \cite{Cazenave_Lions_orbital_stab, PLLcc1, PLLcc2} and the references therein.
Since we do not know whether the ground states are unique (up to the equation invariants),
the statement
only tells us that a perturbation of a ground state stay close (uniformly in time) to \emph{the manifold of all the possible ground states}.
This is weaker than the
expected conclusion which would assert that ``a perturbation of a given ground state stay close (uniformly in time) to \emph{the manifold generated by this ground state and the equation invariants (phase and translation)}''.

\subsection{Strengthened orbital stability}

A strengthened result can be obtained by using an alternative approach, based on the study of the linearization of the energy around a ground state (see
\cite{MaMe,Weinstein1,Weinstein2}; we also refer the reader to the lecture notes \cite[\S 2.6]{Martel} and the references therein). To be more specific, we fix $M>M_{0}$ and we consider a ground state $(Q,\Psi)$ of $J_{M}$ such that $Q$ is positive, radially symmetric and decreasing and such that $\|Q\|_{L_{x}^{2}}^{2}=M$. We introduce
\[
W(u,\psi,\chi)=E(u,\psi,\chi)+\omega\|u\|_{L_{x}^{2}}^{2}.
\]
Next, we linearize this quantity around $(Q,\Psi,0)$: for every $u\in H_{x}^{1}$, $\psi\in L_{x}^{2}\overbigdot{H}_{z}^{1}$ and $\chi\in L_{x}^{2}L_{z}^{2}$, we have\vspace*{-3pt}
\begin{align*}
\ds W(Q+u,&\Psi+\psi,\chi)=W(Q,\Psi,0)\\
&+ \ds \frac{1}{2}\int_{\R^{d}}\nabla_{x}Q\cdot(\nabla_{x}u+\nabla_{x}\bar{u})\ud x+\omega\int_{\R^{d}}Q(u+\bar{u})\ud x
\\
&+\ds\int_{\R^{d}}\left(\sigma_{1}\star\int_{\R^{n}}\sigma_{2}\Psi\ud z\right)Q(u+\bar{u})\ud x
\\
&+\ds\int_{\R^{d}}\left(\sigma_{1}\star\int_{\R^{n}}\sigma_{2}\psi\ud z\right)Q^{2}\ud x+\frac{1}{2}\iint_{\R^{d}\times\R^{n}}\nabla_{z}\Psi\cdot\nabla_{z}\psi\ud x\ud z
\\
&\ds +\frac{1}{2}\int_{\R^{d}}|\nabla_{x}u|^{2}\ud x+\omega\int_{\R^{d}}|u|^{2}\ud x+\int_{\R^{d}}\left(\sigma_{1}\star\int_{\R^{n}}\sigma_{2}\Psi\ud z\right)|u|^{2}\ud x\\
&\ds +\int_{\R^{d}}\left(\sigma_{1}\star\int_{\R^{n}}\sigma_{2}\psi\ud z\right)Q(u+\bar{u})\ud x+\frac{1}{2c^{2}}\iint_{\R^{d}\times\R^{n}}|\chi|^{2}\ud x\ud z
\\
&+\ds\frac{1}{2}\iint_{\R^{d}\times\R^{n}}|\nabla_{z}\psi|^{2}\ud x\ud z
+ \ds \int_{\R^{d}}\left(\sigma_{1}\star\int_{\R^{n}}\sigma_{2}\psi\ud z\right)|u|^{2}\ud x.
\end{align*}
We write this as $W(Q+u,\Psi+\psi,\chi)=W(Q,\Psi,0)+I_1+\cdots+I_{12}$.
Thanks to \eqref{soliton}, $I_1+I_2+I_3=0$ and thanks to \eqref{wave_soliton}, $I_4+I_5=0$. Let us denote\vspace*{-3pt}
\[
u=f+i g,\qquad f,g\in \mathbb R.
\]
We can rewrite\vspace*{-3pt}
\[
I_ 6+\cdots+I_{11}= \left\langle\mathcal{L}_{+}\begin{pmatrix} f\\\psi \end{pmatrix},\begin{pmatrix} f\\\psi \end{pmatrix}\right\rangle_{L_{x}^{2}\times L_{x}^{2}L_{z}^{2}}+\left\langle L_{-}g,g\right\rangle_{L_{x}^{2}}+\frac{1}{2c^{2}}\,\|\chi\|_{L_{x}^{2}L_{z}^{2}}^{2}
\]
where\vspace*{-3pt}
\begin{equation}\label{L+}
\ds \mathcal{L}_{+}=\begin{pmatrix} \ds-\frac{1}{2}\Delta_{x}+\omega+\left(\sigma_{1}\star\int_{\R^{n}}\sigma_{2}\Psi\ud z\right) & M_{1} \\ M_{2} & \ds-\frac{1}{2}\Delta_{z} \end{pmatrix}
\end{equation}
with\vspace*{-3pt}
\[
M_{1}\psi=\left(\sigma_{1}\star\int_{\R^{n}}\sigma_{2}\psi\ud z\right)Q, \qquad M_{2}f=\sigma_{2}\,(\sigma_{1}\star Qf),
\]
and\vspace*{-3pt}
\begin{equation}\label{L-}
L_{-}=-\frac{1}{2}\Delta_{x}+\omega+\biggl(\sigma_{1}\star\int_{\R^{n}}\sigma_{2}\Psi\ud z\biggr).
\end{equation}
Let us also introduce the operator $L_{+}$ defined by\vspace*{-3pt}
\begin{equation}\label{l+}
L_{+}f=-\frac{1}{2}\Delta_{x}f+\omega f-\kappa(\Sigma\star Q^{2})f-2\kappa(\Sigma\star Qf)Q,
\end{equation}
which will have an important role in the sequel: it is the analog to $\mathcal{L}_{+}$ for $\widetilde{W}(u)=H(u)+\omega\|u\|_{L_{x}^{2}}^{2}$.
We eventually obtain the following decomposition\vspace*{-5pt}
\begin{multline}\label{decomp_W}
W(Q+u,\Psi+\psi,\chi)=W(Q,\Psi,0)+\left\langle\mathcal{L}_{+}\begin{pmatrix} f\\\psi \end{pmatrix},\begin{pmatrix} f\\\psi \end{pmatrix}\right\rangle_{L_{x}^{2}\times L_{x}^{2}L_{z}^{2}}\\
+\left\langle L_{-}g,g\right\rangle_{L_{x}^{2}}
+\frac{1}{2c^{2}}\|\chi\|_{L_{x}^{2}L_{z}^{2}}^{2}+\int_{\R^{d}}\left(\sigma_{1}\star\int_{\R^{n}}\sigma_{2}\psi\ud z\right)|u|^{2}\ud x.
\end{multline}

\begin{rmk}
Relation \eqref{decomp_W} holds true when replacing, for some $\alpha\in\R$, $M_{1}$ and~$M_{2}$ in the definition of $\mathcal{L}_{+}$ by $\alpha M_{1}$ and $(2-\alpha)M_{2}$. However, $\mathcal{L}_{+}$ is self-adjoint only in the particular case $\alpha=1$.
\end{rmk}

The key argument to prove an orbital stability result is to characterize the kernel of~$L_{-}$ and $\mathcal{L}_{+}$ and to prove that these operators are coercive under some orthogonality conditions. The operator $L_{-}$ is a local operator, and we already
have at hand the following statement,
see for example \cite{Weinstein1}.

\begin{lemma}\label{Lemma_L-}
Let \eqref{h1}--\eqref{h3} be fulfilled. We have $\Ker(L_{-})=\mathrm{Span}\{Q\}$ and there exists a constant $\mu>0$ such that for every $g\in H_{x}^{1}$,
\begin{equation}\label{coercivity1}
\left\langle L_{-}g,g\right\rangle_{L_{x}^{2}}\geq\mu\|g\|_{H_{x}^{1}}^{2}-\frac{1}{\mu}\left|\langle g,Q\rangle_{H_{x}^{1}}\right|^{2}.
\end{equation}
\end{lemma}

The difficult part is to obtain an analogous statement for $\mathcal{L}_{+}$.
The method consists in working on the operator $L_{+}$: the knowledge of the kernel of $L_{+}$ will allow us to identify the kernel of $\mathcal{L}_{+}$ and a coercivity property for $L_{+}$ will provide a coercivity property for $\mathcal{L}_{+}$ too.
By direct inspection, it can be checked that
\[
\mathrm{Span}\{\partial_{x_{j}}Q\sep j=1,\ldots,d\} \subset \Ker (L_{+});
\]
we shall work further to establish the reverse inclusion and characterize $\Ker (L_{+})$. Since $L_{+}$ is a non-local operator, classical arguments based on Sturm-Liouville theory are not applicable. We shall need to develop alternative approaches and perturbative arguments, inspired from \cite{Lenzmann}.

We are going to exploit results known for some limiting cases, depending on the dimension $d$. Namely, in the case $d=1$ we will consider the case of the delta function
\begin{equation}\label{delta}
\Sigma^{0}=\delta_{0},
\end{equation}
while in dimension $d=3$ we will consider the case of the Newtonian potential
\begin{equation}\label{newt}\Sigma^0(x)=\ds\frac{1}{|x|}.
\end{equation}
Indeed, for these specific situations the following statement holds.
\begin{lemma}\label{Lemma_L+}

Let $d=1$ with the potential \eqref{delta} or $d=3$ with the potential \eqref{newt}.
Let $Q^0$ stand for the corresponding ground state, associated to the potential $\Sigma^0$ and mass $M>0$.
We have
$\Ker(L_{+})=\mathrm{Span}\{\partial_{x_{j}}Q^0,j=1,\ldots,d\}$. Moreover, there exists a constant $\nu^0>0$ such that for every $f\in H_{x}^{1}$,
\begin{equation}\label{coercivity2}
\left\langle L_{+}f,f\right\rangle_{L_{x}^{2}}\geq\nu^0\|f\|_{H_{x}^{1}}^{2}-\frac{1}{\nu^0}\biggl(\left|\langle f,Q^0\rangle_{L_{x}^{2}}\right|^{2}+\sum\limits_{j=1}^{d}\left|\langle f,\partial_{x_{j}}Q^0\rangle_{L_{x}^{2}}\right|^{2}\biggr).
\end{equation}
\end{lemma}

In the case $d=1$, the result is well known since the paper of M. Weinstein \cite{Weinstein1}. The analysis of the case $d=3$ is quite recent: the characterization of the kernel of $L_{+}$ has been obtained by E. Lenzmann in \cite[Th.\,4]{Lenzmann} and then, based on this characterization, P.\,D'Avenia and M.\,Squassina \cite[Lem.\,2.7]{Avenia_Squassina} established the coercivity property~\eqref{coercivity2}.
We need to extend such a property to potentials with the form~\hbox{$\Sigma=\sigma_{1}\star\sigma_{1}$}:
we~\hbox{denote} by $\mathscr{A}_{d}$ the set of \emph{admissible} form functions
$\sigma_{1}$ such that Lemma~\ref{Lemma_L+} applies in dimension $d$ when $\Sigma=\sigma_{1}\star\sigma_{1}$.
This is made clear by the following definition.

\begin{defin}\label{d:adm}
We say that $\sigma_{1}$ is an admissible form function if it satisfies \eqref{h2}--\eqref{h3} and if there exists a mass interval $I$ of non empty interior such that for every $M\in\nobreak I$ and every positive and radially symmetric minimizer $Q_{M}$ of $K_{M}$,
the conclusions of
Lemma~\ref{Lemma_L+} hold (with a constant $\nu$ which, now, depends on $\Sigma, \omega, M, Q$).
\end{defin}

That $\mathscr{A}_{d}$ is non empty is highly non trivial: in \cite{Lenzmann}, the characterization in Lemma~\ref{Lemma_L+} relies strongly
on the specific form of the Newtonian potential and the scale invariance property of equation \eqref{choquard} in this specific case. Let us temporarily postpone the treatment of this issue; in the meantime, we
state the following lemma which links the properties of the operator $\mathcal{L}_{+}$ to the properties of the operator $L_{+}$. Note that from now on we denote
\[
\mathscr{H}=\bigl \{(u,\psi) \in H_{x}^{1}\times L_{x}^{2}\overbigdot{H}_{z}^{1}\bigr\},
\]
which is a Hilbert space when endowed with the norm defined by
\[
\|(u,\psi) \|_{\mathscr{H}}^{2}=\|u\|_{H_{x}^{1}}^{2}+\|\psi\|_{L_{x}^{2}\overbigdot{H}_{z}^{1}}^{2}.
\]

\begin{lemma}\label{Lemma_LL+}
Assume \eqref{h1}--\eqref{h3}. Let $\sigma_{1}\in\mathscr{A}_{d}$ be an admissible form function and assume that the mass $M$ of the considered ground state $Q$ is in the interval $I$ of Definition~\ref{d:adm}. Then $\Ker(\mathcal{L}_{+})=\mathrm{Span}\{(\partial_{x_{j}}Q,\partial_{x_{j}}\Psi)^{t}, j=1,\ldots,d\}$ and there exists a constant $\tilde{\nu}>0$ such that for every $(f,\psi)\in\mathscr {H}$,
\begin{multline}\label{coercivity3}
\biggl\langle\mathcal{L}_{+}\begin{pmatrix} f\\\psi \end{pmatrix},\begin{pmatrix} f\\\psi \end{pmatrix}\biggr\rangle_{L_{x}^{2}\times L_{x}^{2}L_{z}^{2}}\\[
-5pt]
\geq\tilde{\nu}\|(f,\psi)\|_{\mathscr {H}}^{2}-\frac{1}{\tilde{\nu}}\biggl(\left|\langle f,Q\rangle_{L_{x}^{2}}\right|^{2}+\sum\limits_{j=1}^{d}\left|\langle f,\partial_{x_{j}}Q\rangle_{L_{x}^{2}}\right|^{2}\biggr).
\end{multline}
\end{lemma}

Lemma~\ref{Lemma_L-}, \ref{Lemma_L+} and \ref{Lemma_LL+}
can be interpreted as coercivity properties, up to
some orthogonality conditions, see \cite[\S 4]{Tao}, \cite[Prop.\,2.8 \& 2.9]{Weinstein1}, \cite[Eq.\,(2.14)]{Avenia_Squassina}:
$\mathcal L^+$ is coercive on
the subspace of $\mathscr H$ characterized by
$\langle f,Q\rangle_{L_{x}^{2}}=0=\langle f,\partial_{x_{j}}Q\rangle_{L_{x}^{2}}$.
Considering solutions of
the evolution problems, the latter will
be guaranteed by introducing some appropriate modulation factors $t\mto x(t)$ and $t\mto \gamma(t)$, see \cite [sp.\,\S3]{Weinstein2}. Then,
Lemma~\ref{Lemma_LL+}
is the key ingredient to prove the following orbital stability theorem that strengthens
Theorem~\ref{orbital_stab_1}. The proof is detailed in Section~\ref{OrbitalStability2}.

\begin{theo}\label{orbital_stab_2}
Assume \eqref{h1}--\eqref{h3}. Let $\sigma_{1}\in\mathscr{A}_{d}$ be an admissible form function and assume that $\|Q\|_{L_{x}^{2}}^{2}\in I$.
For every $(u_{0},\psi_{0},\chi_{0})\in H_{x}^{1}\times L_{x}^{2}\overbigdot{H}_{z}^{1}\times L_{x}^{2}L_{z}^{2}$ let us denote by $(u,\psi,\chi=\partial_t\psi)$ the unique solution of \eqref{schro} and \eqref{wave_schro} associated to the initial data $(u_{0},\psi_{0},\chi_{0})$. Let us assume $\|u_{0}\|_{L_{x}^{2}}=\|Q\|_{L_{x}^{2}}$. There exists $\varepsilon_{0}>0$ such that for every $\varepsilon\in(0,\varepsilon_{0})$ we can find $\eta(\varepsilon)>0$ and $\delta(\varepsilon)>0$ such that, if
\[
\|(u_{0}-Q,\psi_{0}-\Psi)\|_{\mathscr {H}}^{2}+\frac{1}{c^{2}}\|\chi_{0}\|_{L_{x}^{2}L_{z}^{2}}^{2}\leq\eta(\varepsilon)^{2} \hspace{0.25cm}\text{and}\hspace{0.25cm} W(u_{0},\psi_{0},\chi_{0})-W(Q,\Psi,0)\leq\delta(\varepsilon),
\]
then there exists two functions $x(t)$ and $\gamma(t)$, continuous in time, such that for every $t\geq 0$, $v(t,\cdot)=e^{-i\gamma(t)}u(t,\cdot+x(t))$ satisfies the following orthogonality conditions\vspace*{-3pt}
\begin{subequations}
\begin{align}
\label{ortho1}
&\left\langle\Re  v,\partial_{x_{j}}Q\right\rangle_{L_{x}^{2}}=0, \qquad j=1,\ldots,d,\\
\label{ortho2}
&\left\langle\Im  v,Q\right\rangle_{H_{x}^{1}}=0\\
\tag*{and}
&\underset{t\geq 0}{\sup}\,\left\|
\big(u(t)-e^{i\gamma(t)}Q(\cdot-x(t)),\psi(t)-\Psi(\cdot-x(t))\big)\right\|_{\mathscr {H}}^{2}+\frac{1}{c^{2}}\|\chi(t)\|_{L_{x}^{2}L_{z}^{2}}^{2}\leq\varepsilon^{2}.
\end{align}
\end{subequations}
\end{theo}

\begin{rmk}
Note that in the regime $c\gg 1/\varepsilon^{2}$, the theorem still applies if the perturbation $\chi_{0}$ is not close to zero.
It is also worth remarking that $\eta(\varepsilon)$ and $\delta(\varepsilon)$ are uniform with respect to $c$.
\end{rmk}

\begin{rmk} Note that it is still possible, in the spirit of results obtained in \cite{Faou} for the standard non linear Schrödinger equation, to justify that orbital stability holds for initial data with arbitrarily large momentum, $(u_0e^{-x\cdot p_0/M},\psi_0,\chi)$ being close to $(Q,\Psi,0)$, on finite time interval:
typically the solution remain at a distance $\epsilon$ of the manifold of the
ground state over time interval of order $\mathscr O(1/\sqrt\epsilon)$, see
\cite[Th.\,4.2.11 \& \S4.6]{LV_PhD}.
The argument relies on the dispersive properties
of the wave equation through Strichartz' estimates, combined to the conservation of the total momentum.
\end{rmk}

It is worth commenting the assumption on the mass of $u_{0}$ which did not appear in Theorem~\ref{orbital_stab_1}. Usually, assuming $\|u_{0}\|_{L_{x}^{2}}=\|Q\|_{L_{x}^{2}}$ is not a restriction. Indeed, as soon as the definition of the map $M\mto Q_{M}$ is meaningful (\emph{i.e.} when ground states are unique or at least locally unique) and defines a continuous map, any small perturbation $u_{0}$ of a ground state $Q_{M}$ is also a small perturbation of the ground state $Q_{\|u_{0}\|_{L_{x}^{2}}}$. Here, relaxing this assumption
requires to justify, first,
that the ground states are (at least locally) unique and, second, their continuity with respect to the mass~$M$. We decided not to
focus on the uniqueness issues in this work; nevertheless we~can provide some hints.
Our approach to find admissible form functions $\sigma_{1}$ is inspired by the strategy developed by E.\,Lenzmann \cite{Lenzmann} in order to prove the uniqueness of ground states (for almost every sufficiently small mass $M$) for the non relativistic Hartree equation.
Therefore, it is likely that a similar result applies here for almost every $M\in I$. Working in this direction may probably allow us to justify that the assumption on the mass of the perturbation $u_{0}$ is indeed not a restriction.

Theorem~\ref{orbital_stab_2} becomes fully meaningful if we are able to characterize the set of admissible form function $\mathscr{A}_{d}$, or at least to justify that $\mathscr{A}_{d}$ contains physically relevant form functions $\sigma_{1}$. This is the purpose of the following sections which contain the most original insights of the paper.
Our results cover the three-dimensional case $d=3$, which is the most relevant physically, and the one-dimensional case $d=1$
for which numerical investigations is more affordable \cite{LVNum2}.
It is worthwhile to see the role of the space dimension in the analysis, and we also provide some hints on the case
$d=2$.

\subsection{The case $d=3$}

Section~\ref{large_class_sigma1} is devoted to the construction of admissible form functions $\sigma_{1}$ in dimension $d=3$. The difficulty in identifying the class of admissible form functions $\sigma_{1}$ is a weakness of the method compared to the approach by concentration-compactness.
Nevertheless this additional restriction allows us to \hbox{obtain} the more precise orbital stability result of Theorem~\ref{orbital_stab_2} and we shall see in Section~\ref{large_class_sigma1} that we can find many form functions $\sigma_{1}$ that fit the physical framework introduced in \cite{BdB}.
We proceed in two steps. The idea is to boil down a perturbative approach for potentials $\Sigma$ close, in an appropriate sense, to $|\cdot|^{-1}$,
and then to push this result by suitable rescalings which allow us to identify physically relevant potentials $\Sigma=\sigma_1\star\sigma_1$ not necessarily close to $|\cdot|^{-1}$.
An important issue in this approach is to clarify the role of the mass constraint:
Theorem~\ref{orbital_stab_1}
applies to any ground state of mass $M>M_{0}$. Hence, we expect stability results that apply to a continuum of possible masses $M$, as stated in Definition~\ref{d:adm}.

\begin{proposition}\label{non_perturbative_L+}
The set $\mathscr{A}_{3}$ of admissible form functions is non empty.
\end{proposition}
We will indeed see in Section~\ref{large_class_sigma1} that the set $\mathscr{A}_{3}$ contains at least every form function~$\sigma_{1}$ satisfying \eqref{h2}--\eqref{h3} and such that $\Sigma^{\lambda,\mu}(x)=\mu\Sigma(\lambda x)$ is \emph{close enough} to $\Sigma^0=|\cdot|^{-1}$ for suitable rescaling parameters $\lambda,\mu>0$. As explained above, our strategy
to identify admissible form functions and to establish the orbital stability for the Schrödinger-wave system
is based on a perturbative analysis from $\Sigma^0$. For this purpose let us introduce the following more precise notations.

\begin{defin}\label{defKUpM}
For a given potential $\Sigma$ we denote $H^{\Sigma}$
and $K_{M}^{\Sigma}$ the corresponding energy defined by \eqref{hamiltonian}, and the minimization problem \eqref{KM}, respectively.
Then we denote by $Q_{M}^{\Sigma}$ a positive and radially symmetric minimizer of $K_{M}^{\Sigma}$ and by $\omega(\Sigma,Q_{M}^{\Sigma})$ the constant $\omega>0$ such that $Q_{M}^{\Sigma}$ is a solution of \eqref{choquard} with $\Sigma$ and $\omega=\omega(\Sigma,Q_{M}^{\Sigma})$. Note that the notation $Q_{M}^{\Sigma}$ could designate several minimizers since a priori we do not get the uniqueness of the minimizers of $K_{M}^{\Sigma}$.
Similarly, the Lagrange multiplier might depend on the selected minimizer $Q_{M}^{\Sigma}$, motivating the notation $\omega(\Sigma,Q_{M}^{\Sigma})$. Moreover, we make precise how the operator $L_{+}$ defined by \eqref{L+} depends on $\Sigma$, $Q$ and $\omega$. Since we will only consider cases where $\omega=\omega(\Sigma,Q)$ we will use the notation $L_{+}=L_{+}(\Sigma,Q)$.
\end{defin}

We consider a sequence $(\Sigma^{\varepsilon})_{\varepsilon>0}$ of smooth potentials satisfying the following assumption:

\begin{enumerate}\renewcommand{\theenumi}{H\arabic{enumi}}\setcounter{enumi}{3}
\item\label{h4}
For every $\varepsilon$ there exists $\sigma_{1}^{\varepsilon}$ satisfying \eqref{h2}--\eqref{h3} such that $\Sigma^{\varepsilon}=\sigma_{1}^{\varepsilon}\star\sigma_{1}^{\varepsilon}$ and the sequence $(\Sigma^{\varepsilon})_{\varepsilon>0}$ converges to $\Sigma^{0}=|\cdot|^{-1}$ in the following sense: for every $R>0$,
\begin{equation}\label{cv_potential}
\|(\Sigma^{\varepsilon}-\Sigma^{0})\mathbf{1}_{|x|\leq R}\|_{L_{x}^{3/2}}+\|(\Sigma^{\varepsilon}-\Sigma^{0})\mathbf{1}_{|x|> R}\|_{L_{x}^{\infty}}\underset{\varepsilon\to 0}{\to}\,0.
\end{equation}
\end{enumerate}
For such family we know that for each $\varepsilon>0$, there exists a mass threshold $M_{0}^{\varepsilon}>0$ such that $K_{M}^{\Sigma^{\varepsilon}}$ is achieved for every $M>M_{0}^{\varepsilon}$. In order to work with a fixed mass $M>0$ we will also assume that $\sup_{0<\varepsilon\leq \varepsilon_0}(M_{0}^{\varepsilon})<+\infty$, for some $\varepsilon_0>0$ and we will consider a mass $M$ such that $M>\sup_{0<\varepsilon\leq \varepsilon_0}(M_{0}^{\varepsilon})$. This assumption is quite reasonable since $\Sigma^{\varepsilon}\to\Sigma^{0}$ and there is no mass threshold in the case $\Sigma=\Sigma^{0}$. We refer the reader to Lemma~\ref{Lemma_mass_threshold} which ensures that this assumption is indeed always valid in the previous context and any mass $M$ can be reached that way.

Then we consider a sequence $(Q^{\varepsilon})_{\varepsilon>0}$ of smooth, positive, radially symmetric and decreasing functions and a sequence $(\omega^{\varepsilon})_{\varepsilon>0}$ of positive numbers such that $Q^{\varepsilon}=Q_{M}^{\Sigma^{\varepsilon}}$ and $\omega^{\varepsilon}=\omega(\Sigma^{\varepsilon},Q_{M}^{\Sigma^{\varepsilon}})$. In particular each $Q^{\varepsilon}$ is a solution of \eqref{choquard} with $\Sigma=\Sigma^{\varepsilon}$ and $\omega=\omega^{\varepsilon}$. We also consider $Q^{0}$, the unique positive and radially symmetric minimizer of $K_{M}^{\Sigma^{0}}$. Note that $Q^{0}$ is also decreasing and we can find $\omega^{0}>0$ such that $Q^{0}$ is a solution of \eqref{choquard} with $\Sigma=\Sigma^{0}$ and $\omega=\omega^{0}$.
Hence, the cornerstone of the analysis is given by the following result, established
in Section~\ref{Perturbative}.

\begin{proposition}\label{perturbative_L+}
Let \eqref{h1}--\eqref{h3} be fulfilled. With the previous notations and assuming moreover \eqref{h4}, the following properties hold.
\begin{enumeratei}
\item\label{perturbative_L+i}
\textup{Convergence.} For every $\delta>0$ there exists $\varepsilon_{0}>0$ such that for every $0<\varepsilon<\varepsilon_{0}$,
\[
\|Q^{\varepsilon}-Q^{0}\|_{H_{x}^{1}}+|\omega^{\varepsilon}-\omega^{0}|<\delta.
\]
\item\label{perturbative_L+ii}
\textup{Coercivity.} There exists $\bar{\varepsilon}_0>0$ such that for every $\varepsilon\in(0,\bar{\varepsilon}_{0})$, $Q^{\varepsilon}=Q_{M}^{\Sigma^{\varepsilon}}$ and $\omega^{\varepsilon}=\omega(\Sigma^{\varepsilon},Q_{M}^{\Sigma^{\varepsilon}})$ there exists $\nu(\Sigma^{\varepsilon},Q^{\varepsilon},\omega^{\varepsilon})>0$ satisfying, for every $f\in H_{x}^{1}$,
\[
\left\langle L_{+}(\Sigma^{\varepsilon},Q^{\varepsilon},\omega^{\varepsilon})f,f\right\rangle_{L_{x}^{2}}\!\geq\!
\nu(\Sigma^{\varepsilon},Q^{\varepsilon},\omega^{\varepsilon})\|f\|_{H_{x}^{1}}^{2}-\frac{1}{\nu^{0}}\biggl(\!\left|\langle f,Q^{\varepsilon}\rangle_{L_{x}^{2}}\right|^{2}+\sum_{j=1}^{3}\left|\langle f,\partial_{x_{j}}Q^{\varepsilon}\rangle_{L_{x}^{2}}\right|^{2}\!\biggr),
\]
where $\nu^{0}$ is the best constant possible in Lemma~\ref{Lemma_L+}. Moreover, $\nu(\Sigma^{\varepsilon},Q^{\varepsilon},\omega^{\varepsilon})\nearrow\nu^{0}$ when $\varepsilon\to 0$. This coercivity inequality insures that the kernel of $L_{+}(\Sigma^{\varepsilon},Q^{\varepsilon},\omega^{\varepsilon})$ is spanned by the $\partial_{x_{j}}Q^{\varepsilon}$ and Lemma~\ref{Lemma_L+} applies to the kernel $\Sigma^\varepsilon$ as well.
\end{enumeratei}
\end{proposition}

\begin{rmk}
In point \eqref{perturbative_L+i}, $\varepsilon_{0}$ depends on the chosen sequence $(Q^{\varepsilon})_{\varepsilon>0}$ whereas in point (ii), $\bar{\varepsilon}_{0}$ is the same for every sequence $(Q^{\varepsilon})_{\varepsilon>0}$. However, how the coercivity constant $\nu(\Sigma^{\varepsilon},Q^{\varepsilon},\omega^{\varepsilon})$ converges to $\nu^{0}$ depends on the considered sequence.
\end{rmk}

In this proposition, how small $\bar{\varepsilon}_{0}$ has to be depends on $M$; hence the result cannot be extended to consider, for a fixed potential $\Sigma^{\varepsilon}$ close to $\Sigma^{0}$, a continuum of possible masses $M$.
The statement applies for a given mass $M$ but it is not sufficient to justify that $\mathscr A_{3}$ is non empty. This issue is addressed in Section~\ref{large_class_sigma1}.

\begin{rmk}
Another relevant example with
a non local definition of the potential without scale invariance is the case of the Hartree equation with the Yukawa potential $\Sigma(x)=\sfrac{e^{-\mu|x|}}{|x|}$,
which corresponds to a coupling between the Schrödinger equation and the screened Poisson equation $\mu^{2}\Phi-\Delta_{x}\Phi=|u|^{2}$
for the potential.
The stability analysis for this problem is performed by a variational approach in \cite{Zhang}
and an improved statement has been obtained in \cite{jmaa} by using a perturbative approach next to $\mu=0$, in the spirit of the arguments we are developing here.
\end{rmk}

\subsection{The case $d=1$}
The characterization of $\mathscr{A}_{1}$ is much easier. This is related to the remarks made in Section~\ref{scaling}. Indeed, we can adapt the same strategy than developed for $d=3$, but now considering perturbations around $\delta_{0}$, and using the fact that the cubic nonlinear Schrödinger equation is $L^{2}$-subcritical for $d=1$. We obtain the following result.

\begin{proposition}\label{non_perturbative_L+_d=1}
If $\sigma_{1}$ satisfies \eqref{h2}--\eqref{h3}, then $\sigma_{1}\in \mathscr{A}_{1}$. Moreover there exists a mass $M^{*}>0$ such that $(0,M^{*})\subset I$. As a consequence, we obtain $M_{0}=0$.
\end{proposition}

Let $\sigma_{1}$ satisfy \eqref{h2}--\eqref{h3} and consider the sequence $(\Sigma^{\varepsilon})_{\varepsilon>0}$ of smooth potentials defined by
\[
\Sigma^{\varepsilon}(x)=\varepsilon^{-1}\Sigma(\varepsilon^{-1}x), \quad \Sigma=\sigma_{1}\star\sigma_{1}.
\]
This sequence converges to $\int \Sigma(y)\ud y\ \delta_{0}$, in the usual sense where
\[
\lim_{\varepsilon \to 0}\int \Sigma^\varepsilon(x)\phi(x)\ud x=\phi(0)\times \int \Sigma(y)\ud y
\]
holds for any $\phi\in C^0_c(\mathbb R)$. We know that for each $\varepsilon>0$, there exists a mass threshold $M_{0}^{\varepsilon}\geq 0$ such that $K_{M}^{\Sigma^{\varepsilon}}$ is achieved for every $M>M_{0}^{\varepsilon}$. As in the case $d=3$, we can prove (thanks to an easy adaptation of Lemma~\ref{Lemma_mass_threshold}) that $\sup_{0<\varepsilon<1}(M_{0}^{\varepsilon})<+\infty$.
Thanks to the scaling relations of Section~\ref{scaling}, applied with $\mu=\lambda=\varepsilon^{-1}$, we can express $M_{0}^{\varepsilon}$ in terms of $\varepsilon$ and $M_{0}^{1}$: $M_{0}^{\varepsilon}=\varepsilon^{-1}M_{0}^{1}$. Combining this relation to the boundedness of $\sup_{0<\varepsilon<1}(M_{0}^{\varepsilon})$ implies $M_{0}^{1}=0$ and then $M_{0}^{\varepsilon}=0$, as announced in Theorem~\ref{existence_GS}\eqref{existence_GSv}.

Hence, for a given mass $M>0$ we can consider a sequence of ground states $Q^{\varepsilon}=\nobreak Q_{M}^{\Sigma^{\varepsilon}}$ and Lagrange multipliers $\omega^{\varepsilon}=\omega(\Sigma^{\varepsilon},Q_{M}^{\Sigma^{\varepsilon}})$ and we can justify that the conclusions of Proposition~\ref{perturbative_L+} (where $Q^{0}$ is now the unique positive and even minimizer of $K_{M}^{\delta_{0}}$ and $\omega^{0}$ is the corresponding Lagrange multiplier) also hold in this case. We refer the reader to Example~3 in Section~\ref{large_class_sigma1} where we briefly justify how the conclusions of Proposition~\ref{perturbative_L+} allow us to obtain Proposition~\ref{non_perturbative_L+_d=1}.

\begin{rmk}
There is no major difficulty in order to adapt the proof of Proposition~\ref{perturbative_L+} to the case $d=1$.
The compact embedding $H_\mathrm{rad}^{1}(\R^{d})\to L^{p}(\R^{d})$, $p\in(2,p_{c})$ holds in dimension $d\geq 2$,
but we can exploit the fact that each $Q^{\varepsilon}$ is decreasing in order to recover some compactness result
(as in the proof of Theorem~\ref{existence_GS}\eqref{existence_GSi}).
\end{rmk}

\subsection{The case $d=2$}

One may naturally wonder what happens in dimension \hbox{$d=2$}.
This discussion is speculative, and it can be skipped without compromising
the under\-standing of the remainder of the paper.
The discussion in Section~\ref{scaling} supports the intuition that the case $d=2$
is likely more intricate than $d=1$ and
studying this situation can shed some light
on the restrictions on $\sigma_{1}$
adopted when $d=3$ (compare the characterization of the set of admissible form functions for $d=1$ in Proposition~\ref{non_perturbative_L+_d=1}
to the weaker statement in Proposition~\ref{non_perturbative_L+}).

The role of the dimension $d$ appears in the analysis of the minimization problem for $K_{M}$: for $d=1$, the mass threshold $M_0$ is zero, while it is strictly positive in higher dimensions.
We can thus expect to obtain useful information by
studying more precisely the value of $M_0$ and
considering
ground states having a mass close to $M_{0}$.
Going back to the proof of $M_{0}>0$ (see the proof of Lemma~\ref{basic_lemma}-f)), we are led to study
the best constant $C>0$ in the inequality
\[
\left|\int\big(\Sigma\star|u|^{2}\big)\,|u|^{2}\ud x\right|\leq C\|\nabla_{x}u\|_{L_{x}^{2}}^{2}\|u\|_{L_{x}^{2}}^{2}.
\]
This yields to the minimization problem\vspace*{-5pt}
\[
a^{\Sigma}:=\underset{\substack{u\in H_{x}^{1} \\ u\neq 0}}{\inf}\;A^{\Sigma}(u), \hspace{0.5cm} A^{\Sigma}(u)=\ds\frac{\ds\|\nabla_{x}u\|_{L_{x}^{2}}^{2}\|u\|_{L_{x}^{2}}^{2}}{\ds\int\big(\Sigma\star|u|^{2}\big)\,|u|^{2}\ud x}.
\]
Indeed, we can prove that $M_{0}=a^{\Sigma}/\kappa$ and we are thus led to \emph{compute} $a^{\Sigma}$.
Coming back to the scaling relations discussed in Section~\ref{scaling}, we set $\Sigma^{\varepsilon}(x)=\varepsilon^{-2}\Sigma(\varepsilon^{-1}x)$, $u^{\varepsilon}(x)=\varepsilon^{-1}u(\varepsilon^{-1}x)$ and we check that $A^{\Sigma^{\varepsilon}}(u^{\varepsilon})=A^{\Sigma}(u)$.
Accordingly, we have $a^{\Sigma^{\varepsilon}}=a^{\Sigma}$ for every $\varepsilon>0$. Passing \emph{formally} to the limit $\varepsilon\to 0$ in this relation, which amounts to saying $\Sigma^{\varepsilon}\to\delta_{0}$ (note that up to change the value of $\kappa$, we can always assume that $\|\Sigma\|_{L_{x}^{1}}=1$), would lead to identify $a^{\Sigma}$ and $a^{\delta_{0}}$ where $a^{\delta_{0}}$ stands for the best constant in the Gagliardo-Nirenberg inequality\vspace*{-3pt}
\[
\|u\|_{L_{x}^{4}}^{4}\leq \ds\frac{1}{a^{\delta_{0}}}\,\|\nabla_{x}u\|_{L_{x}^{2}}^{2}\|u\|_{L_{x}^{2}}^{2}, \hspace{0.5cm} a^{\delta_{0}}=\underset{\substack{u\in H_{x}^{1} \\ u\neq 0}}{\inf}\, A^{\delta_{0}}(u).
\]
It is well known that $A^{\delta_{0}}$ admits minimizers, see
the pioneering work \cite[Th.\,B]{weinstein1982} which points out the connection to the nonlinear Schrödinger equation,
and the recent reviews \cite{Bol, DPD}~; these minimizers are of arbitrary mass (thanks to the relation $A^{\Sigma}(\theta u)=A^{\Sigma}(u)$ for every $\theta\neq 0$) and solution of the equation\vspace*{-3pt}
\[
-\frac{1}{2}\Delta_{x}Q+\frac{1}{2}\frac{\|\nabla_{x}Q\|_{L_{x}^{2}}^{2}}{\|Q\|_{L_{x}^{2}}^{2}}\,Q-\frac{a^{\delta_{0}}}{\|Q\|_{L_{x}^{2}}^{2}}\,Q^{3}=0.
\]
By considering a minimizer of mass $a^{\delta_{0}}$ and thanks to the rescaling $Q^{\lambda}(x)=\lambda Q(\lambda x)$ (which leaves both the equation and the mass of the minimizer invariant) we can simply consider the equation\vspace*{-5pt}
\[
-\frac{1}{2}\Delta_{x}Q+Q-Q^{3}=0.
\]
It is well known that this equation admits a unique positive and radially symmetric solution, see \cite{Kwo,McL}. By denoting $Q^{\delta_{0}}$ this unique solution we eventually obtain $M_{0}=\|Q^{\delta_{0}}\|_{L_{x}^{2}}^{2}/\kappa$.
This discussion makes formally a bridge appear with the asymptotic system \eqref{nls} when $\Sigma\to\delta_{0}$; this intuition might be a guide for further analysis, which relies on the following known results for the cubic nonlinear Schrödinger equation in dimension $d=2$, where we keep the notations of
Definition~\ref{defKUpM}.
Details can be found in the seminal paper \cite{merle1993}, and in the in-depth review \cite{Raph}
which contains complete references.

\begin{theo}\label{th:218}
Let $d=2$ and $M_{0}=\|Q^{\delta_{0}}\|_{L_{x}^{2}}^{2}/\kappa$. The following assertion hold:
\begin{enumeratei}
\item\label{th:218i}
For every $0\leq M\leq M_{0}$, $K_{M}^{\delta_{0}}=0$ while $K_{M}^{\delta_{0}}=-\infty$ when $M>M_{0}$.
\item\label{th:218ii}
If $u\in H_{x}^{1}$ is such that $0<\|u\|_{L_{x}^{2}}^{2}\leq M_{0}$ and $H^{\delta_{0}}(u)=0$, then there exists $\lambda_{0}>0$, $x_{0}\in\R^{2}$ and $\gamma_{0}\in\R$ such that\vspace*{-3pt}
\[
u(x)=\frac{\lambda_{0}}{\sqrt{\kappa}}\,Q^{\delta_{0}}(\lambda_{0}x-x_{0})e^{i\gamma_{0}}.
\]
As a consequence $\|u\|_{L_{x}^{2}}^{2}=M_{0}$.
\item\label{th:218iii}
Let $L_{+}^{\delta_{0}}:=L_{+}(\delta_{0},Q^{\delta_{0}}/\sqrt{\kappa})$. There exists a universal constant $\nu>0$ such that for every $f\in H_{x}^{1}$,\vspace*{-5pt}
\begin{multline}\label{new_coercivity}
\bigl\langle L_{+}^{\delta_{0}}f,f\bigr\rangle_{L_{x}^{2}}
\geq\nu\|f\|_{H_{x}^{1}}^{2}\\[
-5pt]
-\frac{1}{\nu}\Bigl(\left|\langle f,Q^{\delta_{0}}\rangle_{L_{x}^{2}}\right|^{2}+\sum\limits_{j=1}^{2}\left|\langle f,\partial_{x_{j}}Q^{\delta_{0}}\rangle_{L_{x}^{2}}\right|^{2}+\left|\langle f,x\cdot\nabla_{x}Q^{\delta_{0}}+Q^{\delta_{0}}\rangle_{L_{x}^{2}}\right|^{2}\Bigr).
\end{multline}
\item\label{th:218iv}
If $\|u_{0}\|_{L_{x}^{2}}^{2}<M_{0}$, then the unique solution $u$ of \eqref{nls} with initial data $u_{0}$
satisfies the following
scattering estimate: there exists $u_{\infty}\in H_{x}^{1}$ such that
\[
\|u(t)-S(t)u_{\infty}\|_{H_{x}^{1}}\,\underset{t\to+\infty}{\to}\,0,
\]
where $S(t)u_{\infty}$ stands for the unique solution of the linear Schrödinger equation with initial data $u_{\infty}$.
\item\label{th:218v}
If $\|u_{0}\|_{L_{x}^{2}}^{2}=M_{0}$, then there are only three possible scenario for the unique solution $u$ of \eqref{nls} associated to the initial data $u_{0}$:
\begin{itemize}
\item $u$ is a solitary wave (up to the equation's invariants),
\item $u$ blows up in finite time,
\item $u$ is globally defined in time and satisfies the scattering property.
\end{itemize}
\end{enumeratei}
\end{theo}

Here, we have obtained the analogue of point \eqref{th:218i} when a smooth potential $\Sigma$ replaces the delta function $\delta_{0}$: in this case the only difference is that $K_{M}^{\Sigma}$ is finite and strictly negative when $M>M_{0}$. Point \eqref{th:218ii} gives the characterization of the manifold of all possible ground states of mass $M_{0}$. Compared to the case $d=1$, here the manifold is parametrized by an additional parameter ($\lambda_{0}\in\R_{+}^{*}$) which is the translation of the $L^{2}$-criticality of this case. Hence the coercivity relation in \eqref{th:218iii} naturally involves an additional orthogonality condition, compared to \eqref{coercivity2}. From this discussion, we can address the following questions for future investigations.

\begin{question}
Does $A^{\Sigma}$ admit minimizers? If it is so, what is the dimension of the manifold of all minimizers? It involves at least three free parameters (one for the phase and two for the translation), but does it need an additional parameter~$\lambda_{0}$? In~other words, does it exist a transformation $\mathcal{T}_{\lambda_{0}}$, continuous with respect to the one-dimensional parameter $\lambda_{0}$, which does not correspond to a translation or a change of phase, and such that if $Q$ is a ground state then $\mathcal{T}_{\lambda_{0}}Q$ is a ground state too~? Moreover, if such a transformation exists, does it preserve the $L^{2}$-norm of its argument (for every~$\lambda_{0}$, $\|\mathcal{T}_{\lambda_{0}}Q\|_{L_{x}^{2}}=\|Q\|_{L_{x}^{2}}$)?
\end{question}

Depending on the answers, it will be possible to obtain a coercivity relation of the form \eqref{coercivity2} or with an additional orthogonality relation as in \eqref{new_coercivity}. Then, such results will allow us to justify the (un-)stability of ground states of mass $M_{0}$. Note that the existence of a transformation $\mathcal{T}_{\lambda_{0}}$ is quite natural. Indeed, the continuity of the ground states with respect to their mass is at least expected in any dimension $d$. The main issue is to determine whether
or not the transformation also preserves the mass, as
it does for the formal limit case $\delta_{0}$.

\begin{question}
Is it possible to extend the conclusions for ground states of mass $M_{0}$ (if they do exist) to ground states of mass $M>M_{0}$ close to $M_{0}$?
\end{question}

\begin{question}
Does the analogue of point \eqref{th:218iv} still hold true when $u$ is a solution of~\eqref{hartree_bis}?
\end{question}

Thanks to the smoothness of the potential $\Sigma$, we already know that every solution of \eqref{hartree_bis} is globally defined in time. This excludes the scenario where solutions blow up in finite time.
This is a major difference between the dynamics of \eqref{nls} and \eqref{hartree_bis}. Nevertheless,
the similar structure of the infimum of their energy when $M<M_{0}$ suggests that solutions of \eqref{hartree_bis} with a mass strictly less than $M_{0}$
obey the scattering property. If it is so, this major difference with the case $d=1$ (for which ground states exist for any mass) would indicate that dynamics specific to $L^{2}$-critical or $L^{2}$-super-critical equations also hold when $d\geq 2$ and \eqref{hartree_bis} is considered with a smooth potential~$\Sigma$. As a consequence, obtaining positive results of stability when $d\geq 2$ seems much more challenging than in the case $d=1$.

In this paper this difficulty is treated at the price of restricting to potentials $\Sigma$ \emph{close} to $|\cdot|^{-1}$
(instead of being close to $\delta_0$).
This viewpoint takes advantage
of the fact that \eqref{hartree_bis} is $L^{2}$-subcritical when $d=3$ and $\Sigma=|\cdot|^{-1}$, which allows us
to proceed with a perturbative analysis.
Equation \eqref{hartree_bis} is equally $L^{2}$-subcritical when $\Sigma=|\cdot|^{-\alpha}$ with $1\leq\alpha<2$ (the case $\alpha=2$ being $L^{2}$-critical): up to the knowledge of stability results for these cases, the strategy developed in the paper could be adapted to any potential $\Sigma$ close to $|\cdot|^{-\alpha}$, when $1\leq \alpha<2$.
For the same reason, these arguments equally can be used to deal with higher dimensions $d>3$,
since the results of \cite{Lenzmann} likely extend to this case.

\section{Existence of ground states: proof of Theorem~\ref{existence_GS}}\label{GroundStates}

Let us gather the basic properties of $I_{M}$, $J_{M}$ and $K_{M}$ in the following lemma, which is further illustrated by Figure~\ref{figIM}.

\begin{lemma}\label{basic_lemma}
Let \eqref{h1}--\eqref{h2} be fulfilled. The following assertions hold:
\begin{enumeratea}
\item\label{basic_lemmaa}
$M\mto I_{M}$ is non increasing.
\item\label{basic_lemmab}
$I_{0}=J_{0}=0$ are reached at $(u,\psi,\chi)=(0,0,0)$ and $K_{0}=0$ is reached at $u=0$.
\item\label{basic_lemmac}
For every $M\geq 0$, $-\infty<I_{M}\leq J_{M}\leq K_{M}\leq 0$.
\item\label{basic_lemmad}
There exists a mass threshold $M_{0}\geq 0$ such that $I_{M}=0$ for $M\in[0,M_{0}]$ and $I_{M}<0$ for $M>M_{0}$.
\item\label{basic_lemmae}
If $I_{M}<0$ is reached at $(u,\psi,\chi)$, then $\|u\|_{L_{x}^{2}}^{2}=M$ and $J_{M}=I_{M}$ is reached at $(u,\psi,\chi)$. Moreover $\chi=0$, $\psi=\Gamma\,\sigma_{1}\star|u|^{2}$ and $u\in\mathcal{S}(\R^{d})$ is a solution of \eqref{choquard} for a certain $\omega>0$. In particular $K_{M}=J_{M}$ is reached at $u$.
\item\label{basic_lemmaf}
If $d\geq 2$, then $M_{0}>0$.
\end{enumeratea}
\end{lemma}

Before proving this lemma let us make several remarks:
\begin{itemize}
\item Points \eqref{basic_lemmac} and \eqref{basic_lemmad} tell us that $K_{M}=J_{M}=I_{M}=0$ when $0\leq M\leq M_0$;
moreover, when $M>M_0$, points \eqref{basic_lemmad}, \eqref{basic_lemmae} together
with Theorem~\ref{existence_GS}\eqref{existence_GSi} imply $K_{M}=J_{M}=I_{M}<0$.
Hence, we have $K_{M}=J_{M}=I_{M}$ for every $M\geq 0$.
\item Points d) and e) coupled with Theorem~\ref{existence_GS}\eqref{existence_GSi} imply that $J_{M}$ is reached for $M>M_{0}$ and improve also point \eqref{basic_lemmaa}: $I_{M}=0$ for $M\in[0,M_{0}]$ and
$M\mto I_{M}$ is decreasing on $(M_{0},+\infty)$.
Indeed, if $M\leq M'$, we have $\{u\in H^1_x\sep \|u\|_{L^2}\leq M\}\subset \{u\in H^1_x\sep \|u\|_{L^2}\leq M'\}$
so that $I_M\geq I_{M'}$: $M\mto I_M$ is non increasing.
From d), we infer that if $M\leq M'\leq M_0$, then $I_M=I_{M'}=0$ holds and if $M_0< M< M'$, then both $I_M$ and $I_{M'}$ are negative.
In the latter case, by Theorem~\ref{existence_GS}\eqref{existence_GSi} and \eqref{basic_lemmae}, the constraint is saturated by minimizers.
Suppose $I_M=I_{M'}$. By monotonicity, we get $I_M=I_{M'}=I_\mu<0$ for any $M<\mu<M'$. Hence, the minimum $I_{M'}$ would be reached by a minimizer with mass $\mu<M'$, contradicting e).
Therefore $M\mto I_M$ is decreasing on $(M_{0},+\infty)$.
\item The proof of point \eqref{basic_lemmaf} will give us the following additional information on $M_{0}$:
\begin{equation}\label{estimation_M1}
0<\frac{1}{\kappa C^{2}\|\Sigma\|_{L_{x}^{\sfrac{d}{2}}}}\leq M_{0},
\end{equation}
with $C$ a constant related to Gagliardo-Nirenberg estimates.
\end{itemize}

\begin{figure}[htb]
\definecolor{ffqqqq}{rgb}{1,0,0}
\definecolor{qqqqff}{rgb}{0,0,1}
\definecolor{zzqqzz}{rgb}{0.6,0,0.6}
\definecolor{qqqqcc}{rgb}{0,0,0.8}
\begin{center}
\resizebox{0.4\linewidth}{!}{
\begin{tikzpicture}[line cap=round,line join=round,>=triangle 45,x=2.0cm,y=2.0cm,scale=.65]
\draw[->,color=black] (-0.2,0) -- (4.38,0);
\foreach \x in {-0.5,0.5,1,1.5,2,2.5,3,3.5,4}
\draw[shift={(\x,0)},color=black] (0pt,-2pt);
\draw[->,color=black] (0,-2.33) -- (0,1.7);
\foreach \y in {-2,-1.5,-1,-0.5,0.5,1,1.5}
\draw[shift={(0,\y)},color=black] (2pt,0pt);
\clip(-0.58,-2.33) rectangle (4.38,1.7);
\draw[ffqqqq, smooth,samples=100,domain=2.0:2.5] plot(\x,{0-(\x)*(\x)+2*(\x)});
\draw[dash pattern=on 6pt off 6pt,color=ffqqqq, smooth,samples=100,domain=2.5:2.8] plot(\x,{0-(\x)*(\x)+2*(\x)});
\draw [color=ffqqqq] (0,0)-- (2,0);
\draw (4.,-0.02) node[anchor=north west] {$ M $};
\draw [color=ffqqqq](0.5,-1.5) node[anchor=north west] {$ I_{M}=J_{M}=K_{M}$};
\draw [color=qqqqcc](1.85,0.4) node[anchor=north west] {$ M_{0}$};
\draw (-0.3,-0.02) node[anchor=north west] {$ 0 $};
\begin{scriptsize}
\fill [color=qqqqff] (2,0) circle (1.5pt);
\end{scriptsize}
\end{tikzpicture}}

\caption[A possible graph representing $I_M, J_M, K_M$ as a function of the mass $M$]{A possible graph representing $I_M, J_M, K_M$ as a function of the mass $M$.
Note that nothing ensures that these functions are differentiable as the picture might indicate.
\label{figIM}}
\end{center}
\end{figure}

\begin{proof}
\textit{Items \eqref{basic_lemmaa} and \eqref{basic_lemmab}} are direct consequences of the definition of $I_{M}$, $J_{M}$ and $K_{M}$. The non trivial parts of c) are to prove that $E(u,\psi,\chi)$ is bounded from below under the mass constraint $\|u\|_{L_{x}^{2}}^{2}=M$ and that $K_{M}\leq 0$. Since for every $(u,\psi,\chi)$,\vspace*{-5pt}
\bgroup
\multlinegap0pt
\begin{multline}\label{gradient_borne}
E(u,\psi,\chi)\geq\frac{1}{2}\|\nabla_{x}u\|_{L_{x}^{2}}^{2}-\left|\int_{\R^{d}}\left(\sigma_{1}\star\int_{\R^{n}}\sigma_{2}\psi\ud z\right)|u|^{2}\ud x\right| \\
\shoveright{+\frac{1}{2}\|\nabla_{z}\psi\|_{L_{x}^{2}L_{z}^{2}}^{2}+\frac{1}{2c^{2}}\|\chi\|_{L_{x}^{2}L_{z}^{2}}^{2}}\\
\geq\frac{1}{2}\|\nabla_{x}u\|_{L_{x}^{2}}^{2}-M\,\|\sigma_{1}\|_{L_{x}^{2}}\|\sigma_{2}\|_{L_{z}^{2n/(n+2)}}\|\psi\|_{L_{x}^{2}L_{z}^{2n/(n-2)}}
+\frac{1}{2}\|\nabla_{z}\psi\|_{L_{x}^{2}L_{z}^{2}}^{2}+\frac{1}{2c^{2}}\|\chi\|_{L_{x}^{2}L_{z}^{2}}^{2},
\end{multline}
\egroup
the Sobolev inequality $\|f\|_{L_{z}^{2n/(n-2)}}\lesssim\|\nabla_{z}f\|_{L_{z}^{2}}$, see e.g. \cite[Th., p.\,125]{Nir} allows us to conclude that $I_{M}>-\infty$. In order to prove $K_{M}\leq 0$ we use the immediate estimate $H(u)\leq \|\nabla_{x}u\|_{L_{x}^{2}}^{2}/2$. Then, for every $u\in H_{x}^{1}$, by setting $u_{\lambda}(x)=\lambda^{d/2}u(\lambda x)$ we get $\|u_{\lambda}\|_{L_{x}^{2}}=\|u\|_{L_{x}^{2}}$ and
\[
H(u_{\lambda})\leq\frac{1}{2}\|\nabla_{x}u_{\lambda}\|_{L_{x}^{2}}^{2}=\frac{\lambda^{2}}{2}\|\nabla_{x}u\|_{L_{x}^{2}}^{2}\underset{\lambda\to 0}{\to}0.
\]

\textit{Item \eqref{basic_lemmad}}. For every $(u,\psi)$ such that the intersections $\mathrm{supp}(u)\cap\mathrm{supp}(\sigma_{1})$ and $\mathrm{supp}(\psi)\cap\mathrm{supp}(\sigma_{1})\times\mathrm{supp}(\sigma_{2})$ are non empty and for every $a\in\R$, we have
\begin{multline*}
E(au,a|\psi|,0)
\\
= a^{2}\left(\frac{1}{2}\|\nabla_{x}u\|_{L_{x}^{2}}^{2}-a\int\left(\sigma_{1}\star\int\sigma_{2}|\psi|\ud z\right)|u|^{2}\ud x+\frac{1}{2}\|\nabla_{z}|\psi|\|_{L_{x}^{2}L_{z}^{2}}^{2}\right)\underset{a\to+\infty}{\to}-\infty
\end{multline*}
and $\|au\|_{L_{x}^{2}}^{2}=a^{2}\|u\|_{L_{x}^{2}}^{2}$. We conclude by using that $I_{M}\leq 0$ and $M\mto I_{M}$ is non increasing.

\textit{Item \eqref{basic_lemmae}}. We argue by contradiction: we suppose that $E(u,\psi,\chi)=I_{M}$ with $\|u\|_{L_{x}^{2}}^{2}=\nobreak m$ and $0<m<M$ (note that $I_{M}<0$ implies $m\neq 0$). We first remark that $I_{M}<0$ implies
\[
\int\left(\sigma_{1}\star\int\sigma_{2}\psi\ud z\right)|u|^{2}\ud x<0.
\]
Then, by considering $v=(M/m)^{1/2}u$, $\varphi=(M/m)^{1/2}\psi$ and $\zeta=(M/m)^{1/2}\chi$ we
have $\|v\|_{L^2}^2=M$ so that $I_{M}\leq E(v,\varphi,\zeta)$
and we get\vspace*{-3pt}
\begin{align*}
E(v,\varphi,\zeta)
&=\frac{M}{m}\biggl(\frac{1}{2}\|\nabla_{x}u\|_{L_{x}^{2}}^{2}+\underbrace{\sqrt{\frac{M}{m}}}_{>1}\underbrace{\int\left(\sigma_{1}\star\int\sigma_{2}\psi\ud z\right)|u|^{2}\ud x}_{<0}
\\[
-5pt]
&\hspace*{5cm}+\frac{1}{2c^{2}}\|\chi\|_{L_{x}^{2}L_{z}^{2}}^{2}+\frac{1}{2}\|\nabla_{z}\psi\|_{L_{x}^{2}L_{z}^{2}}^{2}\biggr)\\[
-10pt]
& <\frac{M}{m}E(u,\psi,\chi)=\frac{M}{m}I_{M}<I_{M},
\end{align*}
a contradiction. Since $(u,\psi,\chi)$ is a minimizer of $J_{M}$, the Euler-Lagrange relations imply the existence of a Lagrange multiplier $\lambda_{u,\psi,\chi}$ such that $\nabla_{u,\psi,\chi}E(u,\psi,\chi)=\lambda_{u,\psi,\chi}\nabla_{u,\psi,\chi}(u\mto\|u\|_{L_{x}^{2}}^{2})=2\lambda_{u,\psi,\chi}(u,0,0)^{t}$. The first two components of this vectorial relation imply that $(u,\psi)$ is a solution of \eqref{soliton}--\eqref{wave_soliton} with $\omega=-\lambda_{u,\psi,\chi}$ and the third component implies that $\chi=0$. Then $\psi=\Gamma\,\sigma_{1}\star|u|^{2}$ (which implies that $K_{M}=J_{M}$ is reached at $u$) and $u$ is a solution of \eqref{choquard} with $\omega=-\lambda_{u,\psi,\chi}$. Moreover, by multiplying \eqref{choquard} by $u$ and integrating over $\R^{d}$ we get\vspace*{-3pt}
\[
\frac{1}{2}\|\nabla_{x}u\|_{L_{x}^{2}}^{2}+\omega\|u\|_{L_{x}^{2}}^{2}-\kappa\iint|u|^{2}(x)\Sigma(x-y)|u|^{2}(y)\ud x\ud y=0.
\]
It follows that\vspace*{-5pt}
\begin{multline*}
0>J_{M}=K_{M}=\frac{1}{2}\|\nabla_{x}u\|_{L_{x}^{2}}^{2}-\frac{\kappa}{2}\iint|u|^{2}(x)\Sigma(x-y)|u|^{2}(y)\ud x\ud y\\
=-\omega\|u\|_{L_{x}^{2}}^{2}+\frac{\kappa}{2}\iint|u|^{2}(x)\Sigma(x-y)|u|^{2}(y)\ud x\ud y
\end{multline*}
and thus $\omega>0$. Eventually, thanks to the fact that $\omega$ is a positive number, one can prove by standard arguments that $u$ is in the Schwartz class (we refer the reader to \cite[Th.\,8]{Lieb_Choquard} and its proof in \cite[Rem.\,1]{Lions_Choquard}).

\textit{Item \eqref{basic_lemmaf}}. Let us denote by $C$ the optimal constant of the homogeneous Sobolev embedding $\|f\|_{L_{x}^{2d/(d-2)}}\leq C\|\nabla_{x}f\|_{L_{x}^{2}}$ (note that this estimate requires $d\geq 3$). Since $E(u,\Gamma\,\sigma_{1}\star|u|^{2},0)=H(u)$ and by using the estimate\vspace*{-3pt}
\begin{align*}
\iint|u|^{2}(x)\Sigma(x-y)|u|^{2}(y)\ud x\ud y&\leq\|\Sigma\star|u|^{2}\|_{L_{x}^{\infty}}\|u\|_{L_{x}^{2}}^{2}\\
&\leq\|\Sigma\|_{L_{x}^{\sfrac{d}{2}}}\|\,|u|^{2}\,\|_{L_{x}^{\spfrac{d}{d-2}}}\|u\|_{L_{x}^{2}}^{2}\\
&=\|\Sigma\|_{L_{x}^{\sfrac{d}{2}}}\|u\|_{L_{x}^{\spfrac{2d}{d-2}}}^{2}\|u\|_{L_{x}^{2}}^{2}\\
&\leq C^{2}\|\Sigma\|_{L_{x}^{\sfrac{d}{2}}}\|\nabla_{x}u\|_{L_{x}^{2}}^{2}\|u\|_{L_{x}^{2}}^{2},
\end{align*}
we eventually obtain\vspace*{-3pt}
\[
E(u,\Gamma\,\sigma_{1}\star|u|^{2},0)\geq\frac{1}{2}\Bigl(1-\kappa C^{2}\|\Sigma\|_{L_{x}^{\sfrac{d}{2}}}\|u\|_{L_{x}^{2}}^{2}\Bigr)\|\nabla_{x}u\|_{L_{x}^{2}}^{2},
\]
and $K_{M}$ is non negative as soon as $1-\kappa C^{2}\|\Sigma\|_{L_{x}^{d/2}}M>0$. The case of the dimension $d=2$ can be treated as follows:\vspace*{-5pt}
\begin{multline*}
\ds\iint|u|^{2}(x)\Sigma(x-y)|u|^{2}(y)\ud x\ud y\leq\|\Sigma\star|u|^{2}\|_{L_{x}^{2}}\|\,|u|^{2}\|_{L_{x}^{2}}\\
\ds\leq\|\Sigma\|_{L_{x}^{1}}\|\,|u|^{2}\,\|_{L_{x}^{2}}\|\,|u|^{2}\|_{L_{x}^{2}}=\|\Sigma\|_{L_{x}^{1}}\|u\|_{L_{x}^{4}}^{4}\leq \widetilde{C}^{2}\|\Sigma\|_{L_{x}^{1}}\|\nabla_{x}u\|_{L_{x}^{2}}^{2}\|u\|_{L_{x}^{2}}^{2},
\end{multline*}
where the last estimate is obtained thanks to the Gagliardo-Nirenberg inequality.
We~deduce that $K_{M}$ is non negative provided $1-\kappa \tilde C^{2}\|\Sigma\|_{L_{x}^{d/2}}M>0$.
In any case $d\geq 2$, we conclude that $K_M\geq 0$ for small $M$'s; combined to \textit{c)}, it tells us that $K_M=0$ is such situations.
\end{proof}

Thanks to the previous arguments, Theorem~\ref{existence_GS}\eqref{existence_GSiii} follows from Theorem~\ref{existence_GS}\eqref{existence_GSi}: in the proof we will construct a minimizer such that $u$ is non negative, radially symmetric and non increasing. We are thus left with the task of proving Theorem~\ref{existence_GS}\eqref{existence_GSi}.

\begin{ProofOf}{Theorem~\ref{existence_GS}\eqref{existence_GSi}}
We fix $M>0$ and we consider a minimizing sequence $(u_{\nu},\psi_{\nu},\chi_{\nu})_{\nu\in\mathbb N}$ of $I_{M}$. We start by constructing from this sequence another minimizing sequence with specific properties. Since $E(u_{\nu},\psi_{\nu},0)\leq E(u_{\nu},\psi_{\nu},\chi_{\nu})$, we can take $\chi_{\nu}=0$ for every $\nu$. Moreover, owing to
the diamagnetic inequality \cite[Th.\,7.21]{LL}, we~have
$E(|u_{\nu}|,-|\psi_{\nu}|,0)\leq E(u_{\nu},\psi_{\nu},0)$ and we can suppose $u_{\nu}\geq 0$ and $\psi_{\nu}\leq 0$. Finally, the density of linear combinations of tensor product in $L_{x}^{2}\overbigdot{H}_{z}^{1}$ allows us to assume that every $\psi_{\nu}$ can be written as\vspace*{-3pt}
\[
\psi_{\nu}(x,z)=-\sum\limits_{i=0}^{N_{\nu}}f_{i}^{\nu}(x)g_{i}^{\nu}(z),
\]
where $f_{i}^{\nu}\in L_{x}^{2}$ and $g_{i}^{\nu}\in\overbigdot{H}_{z}^{1}$ are positive functions.
Possibly at the price of decomposing the
$g_{i}^{\nu}$'s on a Hilbert basis of $\overbigdot{H}_{z}^{1}$, we can suppose that for each $\nu$, $(g_{i}^{\nu})_{i\in \mathbb N}$ forms an orthogonal family and we obtain\vspace*{-8pt}
\begin{multline*}
E(u_{\nu},\psi_{\nu},0)=\frac{1}{2}\|\nabla_{x}u_{\nu}\|_{L_{x}^{2}}^{2}\\
-\sum\limits_{i=0}^{N_{\nu}}\left(\int_{\R^{n}}\sigma_{2}(z)g_{i}^{\nu}(z)\ud z\right)\left(\iint_{\R^{d}\times\R^{d}}|u_{\nu}(x)|^{2}\sigma_{1}(x-y)f_{i}^{\nu}(y)\ud x\ud y\right)\\[
-8pt]
+\sum\limits_{i=0}^{N_{\nu}}\|f_{i}^{\nu}\|_{L_{x}^{2}}^{2}\|g_{i}^{\nu}\|_{\overbigdot{H}_{z}^{1}}^{2}.
\end{multline*}
From here we can apply the symmetric decreasing rearrangement theory in order to obtain the inequalities, see \cite[Chap.\,3]{LL}, $\|u_{\nu}^{*}\|_{L_{x}^{2}}^{2}=\|u_{\nu}\|_{L_{x}^{2}}^{2}$,
$\|\nabla_{x}u_{\nu}^{*}\|_{L_{x}^{2}}^{2}\leq\|\nabla_{x}u_{\nu}\|_{L_{x}^{2}}^{2}$, $\|f_{i}^{\nu,*}\|_{L_{x}^{2}}^{2}=\|f_{i}^{\nu}\|_{L_{x}^{2}}^{2}$ and\vspace*{-3pt}
\[
\iint_{\R^{d}\times\R^{d}}|u_{\nu}(x)|^{2}\sigma_{1}(x-y)f_{i}^{\nu}(y)\ud x\ud y\leq \iint_{\R^{d}\times\R^{d}}|u_{\nu}^{*}(x)|^{2}\sigma_{1}^{*}(x-y)f_{i}^{\nu,*}(y)\ud x\ud y,
\]
where $\cdot^{*}$ stands for the symmetric decreasing rearrangement of a given function. Since~$\sigma_{1}$ is assumed non negative, radially symmetric and non increasing, $\sigma_{1}^{*}=\sigma_{1}$ and since\vspace*{-5pt}
\[
\sum\limits_{i=0}^{N_{\nu}}\|f_{i}^{\nu,*}\|_{L_{x}^{2}}^{2}\|g_{i}^{\nu}\|_{\overbigdot{H}_{z}^{1}}^{2}=\biggl\|\sum\limits_{i=0}^{N_{\nu}}f_{i}^{\nu,*}g_{i}^{\nu}\biggr\|_{L_{x}^{2}\overbigdot{H}_{z}^{1}}^{2},
\]
we eventually obtain $E(u_{\nu}^{*},\widetilde{\psi}_{\nu},0)\leq E(u_{\nu},\psi_{\nu},0)$,
where $\widetilde{\psi}_{\nu}=\sum_{i=0}^{N_{\nu}}f_{i}^{\nu,*}g_{i}^{\nu}$
and $u_{\nu}^{*}$ has mass $\leq M$. From now on, we will use the abuse of notation $u_{\nu}=u_{\nu}^{*}$ and $\psi_{\nu}=\widetilde{\psi}_{\nu}$.

Having disposed of these preliminaries, we enter into the heart of the proof. Thanks to \eqref{gradient_borne} we know that $(u_{\nu})_{\nu\in \mathbb N}$ is bounded in $H_{x}^{1}$ and $(\psi_{\nu})_{\nu\in \mathbb N}$ is bounded in $L_{x}^{2}\overbigdot{H}_{z}^{1}$. Hence we can suppose, possibly at the price of extracting subsequences, that $(u_{\nu})_{\nu\in \mathbb N}$ converges weakly to $u$ in $H_{x}^{1}$, $(\psi_{\nu})_{\nu\in \mathbb N}$ converges weakly to $\psi$ in $L_{x}^{2}\overbigdot{H}_{z}^{1}$. We have
\[
\|u\|_{L_{x}^{2}}^{2}\leq M,\quad \|\nabla_{x}u\|_{L_{x}^{2}}^{2}\leq\liminf_{\nu\to \infty}\|\nabla_{x}u_{\nu}\|_{L_{x}^{2}}^{2},\quad \|\psi\|_{L_{x}^{2}\overbigdot{H}_{z}^{1}}^{2}\leq\liminf_{\nu\to \infty}\|\psi_{\nu}\|_{L_{x}^{2}\overbigdot{H}_{z}^{1}}^{2}.
\]
In order to conclude the proof it only remains to prove that\vspace*{-3pt}
\begin{equation}\label{last_convergence}
\int_{\R^{d}}\left(\sigma_{1}\star\int_{\R^{n}}\sigma_{2}\psi_{\nu}\ud z\right)|u_{\nu}(x)|^{2}\ud x\,\underset{\nu\to +\infty}{\to}\,\int_{\R^{d}}\left(\sigma_{1}\star\int_{\R^{n}}\sigma_{2}\psi\ud z\right)|u(x)|^{2}\ud x.
\end{equation}
Indeed, \eqref{last_convergence} now implies $E(u,\psi,0)\leq\liminf_{\nu\to\infty} E(u_{\nu},\psi_{\nu},0)=I_{M}$ and we eventually conclude that $I_{M}$ is reached at $(u,\psi,0)$.

We turn to \eqref{last_convergence}. On the one hand, in the case $d\geq 2$,
we can use the symmetry property of the functions $u_{\nu}\in H_{rad}^{1}$,
which are thus uniformly bounded in $H^1(\mathbb R^d)$ and radially symmetric,
in order to justify the strong convergence of $u_{\nu}$ to $u$ in $L_{x}^{p}$ for $2<p<p_{c}$ (where $p_{c}=2d/(d-2)$ if $d\geq 3$ and $p_{c}=+\infty$ if $d=2$):
see \cite[Prop.\,1.1]{Pllsc}, \cite[Radial Lem.\,1]{Strauss} for such compactness statements based on symmetry properties. On the other hand, in the case $d=1$,
the sequence $(u_{\nu})_{\nu\in \mathbb N}$ is bounded in $H^1(\mathbb R)$
and made of even functions, non increasing on $(0,\infty)$.
With the compact embedding
$H^{1}(\R_{x})\subset L^{p}((-R,R))$, for every $1\leq p<+\infty$ and $0<R<\infty$
and by using a diagonal argument, extracting subsequences if necessary, we know that $(u_{\nu})_{\nu\in\mathbb N}$ converges pointwise to $u$, and strongly in $L^p((-R,R))$ for any finite $R$.
Since for every $\nu$, $u_{\nu}$ is a non negative even function with a non increasing profile, for almost every
$x\in\R$ we get
\[
2|x|\,|u_{\nu}(x)|^{2}\leq\int_{-|x|}^{|x|}|u_{\nu}(y)|^{2}\ud y\leq M \hspace{0.2cm}\text{ and then }\hspace{0.2cm} |u_{\nu}(x)|\leq\sqrt{\frac{M}{2|x|}}\lesssim|x|^{-1/2}.
\]
Thanks to this uniform estimate with respect to $\nu$, we can justify that the sequence $(|u_{\nu}|^{p})_{\nu\in \mathbb N}$ is tight for every $2<p<+\infty$.
It
allows us to justify that the sequence $(u_{\nu})_{\nu\in \mathbb N}$ converges strongly to $u$ in any $L_{x}^{p}$ with $2<p<+\infty$.

We can now conclude the proof as follows:
\bgroup
\multlinegap0pt
\begin{multline*}
\int_{\R^{d}}\left(\sigma_{1}\star\int_{\R^{n}}\sigma_{2}\psi_{\nu}\ud z\right)|u_{\nu}|^{2}\ud x=\int_{\R^{d}}(\sigma_{1}\star|u_{\nu}|^{2})\left(\int_{\R^{n}}\sigma_{2}\psi_{\nu}\ud z\right)\ud x\\
=\int_{\R^{d}}\hspace*{-1mm}\left[(\sigma_{1}\star|u_{\nu}|^{2})-(\sigma_{1}\star|u|^{2})\right]\biggl(\int_{\R^{n}}\hspace*{-2mm}\sigma_{2}\psi_{\nu}\ud z\biggr)\ud x+\int_{\R^{d}}(\sigma_{1}\star|u|^{2})\biggl(\int_{\R^{n}}\hspace*{-2mm}\sigma_{2}\psi_{\nu}\ud z\biggr)\ud x,
\end{multline*}
\egroup
where
\begin{multline*}
\left|\int_{\R^{d}}\left[(\sigma_{1}\star|u_{\nu}|^{2})-(\sigma_{1}\star|u|^{2})\right]\;\left(\int_{\R^{n}}\sigma_{2}\psi_{\nu}\ud z\right)\ud x\right|\\
\lesssim\left\|(\sigma_{1}\star|u_{\nu}|^{2})-(\sigma_{1}\star|u|^{2})\right\|_{L_{x}^{2}}\|\psi_{\nu}\|_{L_{x}^{2}\overbigdot{H}_{z}^{1}}.
\end{multline*}
Note that the weak convergence of $\psi_{\nu}$ to $\psi$ in $L_{x}^{2}\overbigdot{H}_{z}^{1}$ implies the convergence of the second term of the right hand side to $\int(\sigma_{1}\star\int\sigma_{2}\psi\ud z)|u|^{2}\ud x$. Indeed
\bgroup
\multlinegap0pt
\begin{multline*}
\int_{\R^{d}}(\sigma_{1}\star|u|^{2})\left(\int_{\R^{n}}\sigma_{2}\psi_{\nu}\ud z\right)\ud x\\
=\iint_{\R^{d}\times\R^{n}}(\sigma_{1}\star|u|^{2})\sigma_{2}\,\psi_{\nu}\ud x\ud z
=\iint_{\R^{d}\times\R^{n}} (\sigma_{1}
\star|u|^{2})(x)\frac{\widehat{\sigma}_{2}(\zeta)}{ |\zeta|} \,|\zeta|\overline{\widehat{\psi}_{\nu}(x,\zeta)}\ud x\ud \zeta\\
\underset{\nu\to +\infty}{\to}\,\iint_{\R^{d}\times\R^{n}}
(\sigma_{1}\star|u|^{2})(x)\frac{\widehat{\sigma}_{2}(\zeta)}{ |\zeta|}\, |\zeta|\overline{\widehat{\psi}(x,\zeta)}\ud x\ud\zeta
=\int_{\R^{d}}\left(\sigma_{1}\star\int_{\R^{n}}\sigma_{2}\psi\ud z\right)|u|^{2}\ud x,
\end{multline*}
\egroup
where we used $n\geq 3$ in order to justify that $\zeta\mto\sfrac{\widehat{\sigma}_{2}(\zeta)}{|\zeta|}$ is an element of $L_{\zeta}^{2}$.
(Indeed, owing to the fast decay of $\zeta\mto \widehat{\sigma}_{2}(\zeta)$ as $|\zeta|\to \infty$ induced by \eqref{h2},
the only point to be discussed is the integrability near the origin where we use
\[
\frac{|\widehat{\sigma}_{2}(\zeta)|^2}{|\zeta|^2}\sim_{|\zeta|\to 0} \frac{|\widehat{\sigma}_{2}(0)|^2}{|\zeta|^2}
=\biggl(\int \sigma_2(z)\ud z\biggr)^2\frac{1}{|\zeta|^2},
\]
which is integrable --- \resp non integrable --- over $|\zeta|\leq R$ when $n\geq 3$ --- \resp when $n=1$ and $n=2$.)
Thus, it only remains to prove that $\sigma_{1}\star|u_{\nu}|^{2}$ converges strongly to $\sigma_{1}\star|u|^{2}$ in $L_{x}^{2}$. To this end, we remark that
\[
\sigma_{1}\star|u_{\nu}|^{2}-\sigma_{1}\star|u|^{2}=\sigma_{1}\star\left(\,|u_{\nu}-u+u|^{2}-|u|^{2}\right)=\sigma_{1}\star\left(\,|u_{\nu}-u|^{2}+2\Re (u_{\nu}-u)\bar{u}\right).
\]
By using Young's inequalities \cite[Th.\,4.2]{LL} we obtain for every $1\leq p,q\leq+\infty$ with $1/p+1/q=1+1/2$
\begin{align*}
\left\|(\sigma_{1}\star|u_{\nu}|^{2})-(\sigma_{1}\star|u|^{2})\right\|_{L_{x}^{2}}&\leq\|\sigma_{1}\|_{L_{x}^{p}}\left\|\,|u_{\nu}-u|^{2}+2\Re (u_{\nu}-u)\bar{u}\right\|_{L_{x}^{q}}\\
&\leq\|\sigma_{1}\|_{L_{x}^{p}}\bigl(\|u_{\nu}-u\|_{L_{x}^{2q}}^{2}+2\|u_{\nu}-u\|_{L_{x}^{2q}}\|u\|_{L_{x}^{2q}}\bigr).
\end{align*}
Then, since $q$ can be chosen arbitrarily in $[1,2]$, we can always pick $q$ such that $2q\in(2,p_{c})$ and the strong convergence of $u_{\nu}$ to $u$ in $L_{x}^{q}$ for every $q\in(2,p_{c})$ allows us to conclude.
Note that, by construction, the obtained minimizer $u$ is radially symmetric.
\end{ProofOf}

A few comments about the solutions of the minimization problem $J_{M}$ are worthwhile.
As soon as $J_{M}$ is reached at $(u,\psi,\chi)$,
we have
$\chi=0$, $\psi=\Gamma\,\sigma_{1}\star|u|^{2}$ and $K_{M}=J_{M}$ is reached at $u$. Hence, that $J_{M}$ admits a unique minimizer is equivalent to the existence-uniqueness of a minimizer for $K_{M}$.
In \cite{Lieb_Choquard}, E.\,Lieb fully answers the question of the uniqueness of the minimizer of $K_{M}$ for the Newtonian kernel $\Sigma^0(x)=\sfrac{1}{|x|}$ in dimension $d=3$.
The argument proceeds into
two steps.
The former justifies
that minimizers of $K_M$ are positive, radially symmetric and decreasing, up to a translation and a change of phase,
by using that $r\mto 1/r$ is decreasing, see \cite[Lem.\,3 \& Cor.\,4]{Lieb_Choquard}.
The latter proves uniqueness in this class of functions \cite[Th.\,10]{Lieb_Choquard},
and the proof relies crucially on the specific properties of the kernel $\Sigma^0(x)=\sfrac{1}{|x|}$.
Therefore, this analysis does not apply to the present context (note that here $\sigma_{1}$ is non increasing).
Nevertheless, the recent result of L. Ma-L. Zhao \cite[\S 5]{MaZ}
answers positively to the first step, by showing that any non negative solution of \eqref{choquard} is strictly positive, radially symmetric and decreasing.
The idea in \cite{MaZ} consists in writing \eqref{choquard} as a system
\[
\Big(\omega -\ds\frac{1}{2}\Delta\Big) Q= Q X,\qquad X= \kappa\Sigma\star Q^2.
\]
By using the Bessel potential \cite[Chap.\,V, \S3]{Stein}
\[
\mathscr J(x)=\ds\frac{1}{4\pi}\ds\int_0^\infty e^{-\pi x^2/ t}e^{-t/(4\pi)} t^{-(d-2)/2}\ds\frac{\ud t}{t},
\]
$Q$ appears as the solution of an integral equation
\[
Q=\mathscr J\star (QX),\qquad X= \kappa\Sigma\star Q^2.
\]
Using
that $ \mathscr J(x)>0$ for any $x\in\mathbb R^d$,
since we already know that $Q$ is non negative, we deduce that actually $Q$ is positive.
Moreover $\mathscr J$ is decreasing, $\Sigma$ is non increasing, which allows us to adapt the moving plane strategy of
\cite[Th.\,2 \& \S3]{MaZ}: we conclude that $Q$ is radially symmetric,
and monotone decreasing in the radial direction.

\section{Orbital stability: concentration-compactness approach}
\label{OrbitalStability1}

Theorem~\ref{orbital_stab_1} is a consequence of the following lemma.

\begin{lemma}\label{key_lemma_stab1}
Let $M>M_0$.
If $(u_{\nu},\psi_{\nu},\chi_{\nu})_{\nu\in\mathbb N}\subset H_{x}^{1}\times L_{x}^{2}\overbigdot{H}_{z}^{1}\times L_{x}^{2}L_{z}^{2}$ is a sequence such that
\[
\|u_{\nu}\|_{L_{x}^{2}}^{2}\underset{\nu\to+\infty}{\to}M \hspace{0.2cm}\text{ and }\hspace{0.2cm} E(u_{\nu},\psi_{\nu},\chi_{\nu})\underset{\nu\to+\infty}{\to}J_{M},
\]
then there exists a sequence $(x_{\nu})_{\nu\in\mathbb N}$
of elements of
$\R^{d}$ and $(\widetilde{Q},\widetilde{\Psi})\in S_{M}$ such that, up to a sub-sequence,
\[
\|u_{\nu}(\cdot-x_{\nu})-\widetilde{Q}\|_{H_{x}^{1}}^{2}+\|\psi_{\nu}(\cdot-x_{\nu},\cdot)-\widetilde{\Psi}\|_{L_{x}^{2}\overbigdot{H}_{z}^{1}}^{2}+\|\chi_{\nu}\|_{L_{x}^{2}L_{z}^{2}}^{2}\,\underset{\nu\to +\infty}{\to}\,0.
\]
\end{lemma}

Let us first explain how this lemma implies Theorem~\ref{orbital_stab_1}. We argue by contradiction. Let us assume the existence of $\varepsilon>0$ and a sequence of initial data $(u_{0}^{\nu},\psi_{0}^{\nu},\chi_{0}^{\nu})_{\nu\in\mathbb N}$ satisfying
\[
\|u_{0}^{\nu}-Q\|_{H_{x}^{1}}^{2}+\|\psi_{0}^{\nu}-\Psi\|_{L_{x}^{2}\overbigdot{H}_{z}^{1}}^{2}+\|\chi_{0}^{\nu}\|_{L_{x}^{2}L_{z}^{2}}^{2}\,\underset{\nu\to +\infty}{\to}\,0,
\]
and such that for any $\nu\in\mathbb{N}$, the unique solution $(u^{\nu},\psi^{\nu},\chi^{\nu})$ of \eqref{schro}--\eqref{wave_schro}
with initial data $(u_{0}^{\nu},\psi_{0}^{\nu},\chi_{0}^{\nu})$ satisfies for some $t_{\nu}>0$,
\[
\underset{(\widetilde{Q},\widetilde{\Psi})\in S_{M}}{\inf}\,\left(\|u^{\nu}(t_{\nu})-\widetilde{Q}\|_{H_{x}^{1}}^{2}+\|\psi^{\nu}(t_{\nu})-\widetilde{\Psi}\|_{L_{x}^{2}\overbigdot{H}_{z}^{1}}^{2}+\|\chi^{\nu}(t_{\nu})\|_{L_{x}^{2}L_{z}^{2}}^{2}\right)>\varepsilon.
\]
The strong convergence of $u_{0}^{\nu}$ to $Q$ in $H_{x}^{1}$ implies $\|u_{0}^{\nu}\|_{L_{x}^{2}}^{2}\to M$ while the continuity of the energy functional $E$ with respect to $u\in H_{x}^{1}$, $\psi\in L_{x}^{2}\overbigdot{H}_{z}^{1}$ and $\chi\in L_{x}^{2}L_{z}^{2}$ implies
\[
E(u_{0}^{\nu},\psi_{0}^{\nu},\chi_{0}^{\nu})\,\underset{\nu\to+\infty}{\to}\,E(Q,\Psi,0)=J_{M}.
\]
By using the property of mass and energy conservations, we check that the sequence $(u^{\nu}(t_{\nu}) , \psi^{\nu}(t_{\nu}) , \chi^{\nu}(t_{\nu}) )_{\nu\in\mathbb N}$ fulfills the assumptions of Lemma~\ref{key_lemma_stab1} and we eventually obtain the required contradiction.

The proof of Lemma~\ref{key_lemma_stab1} is based on the concentration compactness lemma. In~order to apply this lemma let us state and prove the following result on $J_{M}$.

\skpt
\begin{lemma}\label{lem:42}
\begin{enumeratei}
\item\label{lem:42i}
For every $M>M_{0}$ and for every $\theta>1$, $J_{\theta M}<\theta J_{M}$.
\item\label{lem:42ii}
For every $M>M_{0}$ and for every $\alpha\in(0,1)$,
\begin{equation}\label{sub_additivity}
J_{M}<J_{\alpha M}+J_{(1-\alpha)M}.
\end{equation}
\end{enumeratei}
\end{lemma}

\begin{proof}
\emph{Item \eqref{lem:42i}.} The proof follows the strategy of proof of Lemma~\ref{basic_lemma}\eqref{basic_lemmae}. Since $M>M_{0}$ there exists $(u,\psi,\chi)$ such that $\|u\|_{L_{x}^{2}}^{2}=M$ and $J_{M}=E(u,\psi,\chi)$. Hence, defining for $\theta>1$, $v=\sqrt{\theta}u$, $\varphi=\sqrt{\theta}\psi$ and $\zeta=\sqrt{\theta}\chi$ and following the proof of Lemma~\ref{basic_lemma}\eqref{basic_lemmae} we are led to
\[
J_{\theta M}\leq E(v,\varphi,\zeta)<\theta E(u,\psi,\chi)=\theta J_{M}.
\]

\emph{Item \eqref{lem:42ii}.} Let us distinguish two cases. The first case is $\alpha M\!\leq\!M_{0}$ or \hbox{$(1-\alpha)M\!\leq\!M_{0}$}. The case where these two conditions are satisfied is obvious:\vspace*{-3pt}
\[
J_{M}<0=J_{\alpha M}+J_{(1-\alpha)M}.
\]
Hence let us assume, without loss of generality, that $\alpha M\!\leq\! M_{0}$ and \hbox{$(1-\alpha)M\!>\!M_{0}$}. Since $M\mto J_{M}$ is strictly decreasing on $(M_{0},+\infty)$ (see the remarks after the statement of Lemma~\ref{basic_lemma}) we obtain\vspace*{-3pt}
\[
J_{M}<J_{(1-\alpha)M}=J_{\alpha M}+J_{(1-\alpha)M}.
\]
The second case is $\alpha M>M_{0}$ and $(1-\alpha)M>M_{0}$. In this case we apply the previous item (with $\theta =1/\alpha$ and $\theta=1/(1-\alpha)$) as follows:\vspace*{-3pt}
\begin{align*}
J_{M}=\alpha J_{M}+(1-\alpha)J_{M}&=\alpha J_{\frac{1}{\alpha}\alpha M}+(1-\alpha)J_{\frac{1}{1-\alpha}(1-\alpha)M}\\
&<\alpha\frac{1}{\alpha}J_{\alpha M}+(1-\alpha)\frac{1}{1-\alpha}J_{(1-\alpha)M}=J_{\alpha M}+J_{(1-\alpha)M}.\qedhere
\end{align*}
\end{proof}

\begin{ProofOf}{Lemma~\ref{key_lemma_stab1}}
To start with, owing to \eqref{gradient_borne}, $(u_{\nu})_{\nu\in\mathbb N}$ is bounded in $H_{x}^{1}$ and $(\psi_{\nu})_{\nu\in\mathbb N}$ is bounded in $L_{x}^{2}\overbigdot{H}_{z}^{1}$.
Without loss of generality, we can suppose that for every $\nu\in\N$, $\|u_{\nu}\|_{L_{x}^{2}}^{2}=M$. Indeed, let us set\vspace*{-3pt}
\[
\widetilde{u_{\nu}}=\frac{\sqrt{M}}{\|u_{\nu}\|_{L_{x}^{2}}}\,u_{\nu};
\]
since $\|u_{\nu}\|_{L_{x}^{2}}^{2}\to M$ implies $\|\widetilde{u_{\nu}}-u_{\nu}\|_{H_{x}^{1}}\to 0$, if the conclusion of Lemma~\ref{key_lemma_stab1} holds for the sequence $(\widetilde{u_{\nu}})_{\nu\in \mathbb N}$, then it equally holds with the sequence $(u_{\nu})_{\nu\in \mathbb N}$.
Since $J_{M}\leq E(u_{\nu},\psi_{\nu},0)\leq E(u_{\nu},\psi_{\nu},\chi_{\nu})$ and $E(u_{\nu},\psi_{\nu},\chi_{\nu})\to J_{M}$ when $\nu\to+\infty$ we obtain\vspace*{-3pt}
\[
\frac{1}{2c}\|\chi_{\nu}\|_{L_{x}^{2}L_{z}^{2}}^{2}=E(u_{\nu},\psi_{\nu},\chi_{\nu})-E(u_{\nu},\psi_{\nu},0)\,\underset{\nu\to +\infty}{\to}\,0.
\]
The concentration compactness lemma \cite{PLLcc1, PLLcc2} --- here we use the version that can be found in \cite[Prop.\,1.7.6]{Caze} --- ensures that there are only three different possible scenarios for the behavior of the sequence $(u_{\nu})_{\nu\in\mathbb N}$.
We are going to exclude the first two --- evanescence and dichotomy --- by contradiction, and the last option corresponds to the asserted conclusion.
By the way, this approach proves again the existence of minimizers for $J_M$.

\subsubsection*{Scenario 1: Evanescence} Up to a sub-sequence, for every $2<q<2^{*}$, $(u_{\nu})_{\nu\in\mathbb N}$ converges strongly to $0$ in $L_{x}^{q}$,
where $2^{*}=+\infty$ if $d=1$ or $2$ and $2^{*}=2d/(d-2)$ if $d\geq 3$. Let us assume $d\geq 3$; we have\vspace*{-3pt}
\begin{multline*}
\left|\int_{\R^{d}}\left(\sigma_{1}\star\int_{\R^{n}}\sigma_{2}\psi_{\nu}\ud z\right)|u_{\nu}|^{2}\ud x\right|\leq\left\|\sigma_{1}\star\int_{\R^{n}}\sigma_{2}\psi_{\nu}\ud z\right\|_{L_{x}^{d-1}}\|\,|u_{\nu}|^{2}\|_{L_{x}^{(d-1)/(d-2)}}\\
\leq\|\sigma_{1}\|_{L_{x}^{2(d-1)/(d+1)}}\|\sigma_{2}\|_{L_{z}^{2n/(n+2)}}\|\psi_{\nu}\|_{L_{x}^{2}L_{z}^{2n/(n-2)}}\lesssim\|\psi_{\nu}\|_{L_{x}^{2}\overbigdot{H}_{z}^{1}}\|u_{\nu}\|_{L_{x}^{2(d-1)/(d-2)}}^{2}.
\end{multline*}
Since $(\psi_{\nu})_{\nu\in\mathbb N}$ is bounded in $L_{x}^{2}\overbigdot{H}_{z}^{1}$ and $2<2(d-1)/(d-2)<2^{*}$, we even\-tual\-ly obtain\vspace*{-3pt}
\[
\int_{\R^{d}}\left(\sigma_{1}\star\int_{\R^{n}}\sigma_{2}\psi_{\nu}\ud z\right)|u_{\nu}|^{2}\ud x\,\underset{\nu\to +\infty}{\to}\,0.
\]
Then
\[
J_{M}=\underset{\nu\to+\infty}{\lim}E(u_{\nu},\psi_{\nu},0)=\underset{\nu\to+\infty}{\lim}\Bigl(\frac{1}{2}\|\nabla_{x}u_\nu\|_{L_{x}^{2}}^{2}+\frac{1}{2}\|\nabla_{z}\psi_{\nu}\|_{L_{x}^{2}L_{z}^{2}}^{2}\Bigr)\geq 0,
\]
which contradicts Lemma~\ref{basic_lemma}: since $M>M_0$, we have $J_{M}=I_M<0$.

\subsubsection*{Scenario 2: Dichotomy} Up to possible extraction, there exists two sequences $(v_{\nu})_{\nu\in\mathbb N}$ and $(w_{\nu})_{\nu\in\mathbb N}$, bounded in $H_{x}^{1}$
and such that the following assertions hold
\begin{enumeratei}
\item\label{enum:i}
$\ds \exists\alpha\in(0,1)\text{ such that }\|v_{\nu}\|_{L_{x}^{2}}^{2}\underset{\nu\to+\infty}{\to}\alpha M\text{ and }\|w_{\nu}\|_{L_{x}^{2}}^{2}\underset{\nu\to+\infty}{\to}(1-\alpha)M$,
\item
\label{enum:ii}
$\forall q$ such that $2\leq q<2^{*}$,\quad$\ds\|u_{\nu}\|_{L_{x}^{q}}^{q}-\|v_{\nu}\|_{L_{x}^{q}}^{q}-\|w_{\nu}\|_{L_{x}^{q}}^{q}\underset{\nu\to+\infty}{\to}0$,
\item
\label{enum:iii}
$\ds\underset{\nu\to+\infty}{\liminf}\bigl(\|\nabla_{x}u_{\nu}\|_{L_{x}^{2}}^{2}-\|\nabla_{x}v_{\nu}\|_{L_{x}^{2}}^{2}-\|\nabla_{x}w_{\nu}\|_{L_{x}^{2}}^{2}\bigr)\geq0$.
\end{enumeratei}
\noindent
With \eqref{enum:ii}, we infer
\begin{multline}\label{application_ii}
\left|\int_{\R^{d}}\left(\sigma_{1}\star\int_{\R^{n}}\sigma_{2}\psi_{\nu}\ud z\right)\left(|u_{\nu}|^{2}-|v_{\nu}|^{2}-|w_{\nu}|^{2}\right)\ud x\right|\\
\leq\|\sigma_{1}\|_{L_{x}^{2}}\|\sigma_{2}\|_{L_{z}^{2n/(n+2)}}\|\psi_{\nu}\|_{L_{x}^{2}\overbigdot{H}_{z}^{1}}\left(\int_{\R^{d}}\left||u_{\nu}|^{2}-|v_{\nu}|^{2}-|w_{\nu}|^{2}\right|\ud x\right)\,\underset{\nu\to +\infty}{\to}\,0.
\end{multline}
Note that we can apply \eqref{enum:ii} because in the proof of the concentration compactness lemma
\cite[Prop.\,1.7.6, sp. (1.7.11) \& (1.7.12)]{Caze}
\cite[Proof of Lem.\,I.1 \& Lem.\,III.1]{PLLcc1},
$v_{\nu}$~and~$w_{\nu}$ are built in such way that $|u_{\nu}|^{2}-|v_{\nu}|^{2}-|w_{\nu}|^{2}\geq 0$. Since\vspace*{-10pt}
\bgroup
\multlinegap0pt
\begin{multline*}
E(u_{\nu},\psi_{\nu},0)=\frac{1}{2}\left(\|\nabla_{x}u_{\nu}\|_{L_{x}^{2}}^{2}-\|\nabla_{x}v_{\nu}\|_{L_{x}^{2}}^{2}-\|\nabla_{x}w_{\nu}\|_{L_{x}^{2}}^{2}\right)\\
+\int_{\R^{d}}\left(\sigma_{1}\star\int_{\R^{n}}\sigma_{2}\psi_{\nu}\ud z\right)\left(|u_{\nu}|^{2}-|v_{\nu}|^{2}-|w_{\nu}|^{2}\right)\ud x
+E(v_{\nu},\psi_{\nu},0)+E(w_{\nu},\psi_{\nu},0),
\end{multline*}
\egroup
combining \eqref{application_ii}, \eqref{enum:iii} and \eqref{enum:i}, yields
\begin{multline*}
J_{M}=\underset{\nu\to +\infty}{\lim}\,E(u_{\nu},\psi_{\nu},0)\geq\underset{\nu\to+\infty}{\liminf}\,\left(E(v_{\nu},\psi_{\nu},0)+E(w_{\nu},\psi_{\nu},0)\right)\\
\geq\underset{\nu\to+\infty}{\liminf}\,E(v_{\nu},\psi_{\nu},0)+\underset{\nu\to+\infty}{\liminf}\,E(w_{\nu},\psi_{\nu},0)\geq J_{\alpha M}+J_{(1-\alpha)M},
\end{multline*}
which is a contradiction with \eqref{sub_additivity}, satisfied for $M\in (M_0,2M_0)$.

\subsubsection*{Scenario 3: Compactness} Up to a sub-sequence, there exists a sequence $(x_{\nu})_{\nu\in\mathbb N}$ in~$\R^{d}$ such that $v_{\nu}(x)=u_{\nu}(x-x_{\nu})$ converges weakly to $u$ in $H_{x}^{1}$ and strongly to $u$ in $L_{x}^{q}$ for any $q$ such that $2\leq q<2^{*}$. The sequence $\varphi_{\nu}(x,z)=\psi_{\nu}(x-x_{\nu},z)$ is bounded in $L_{x}^{2}\overbigdot{H}_{z}^{1}$ (note that $\|\varphi_{\nu}\|_{L_{x}^{2}\overbigdot{H}_{z}^{1}}=\|\psi_{\nu}\|_{L_{x}^{2}\overbigdot{H}_{z}^{1}}$) and then, up to a subsequence, $(\varphi_{\nu})_{\nu\in\mathbb N}$ converges weakly to $\psi$ in $L_{x}^{2}\overbigdot{H}_{z}^{1}$. Since $(v_{\nu})_{\nu\in\mathbb N}$ converges strongly to $u$ in $L_{x}^{2}$ we have $\|u\|_{L_{x}^{2}}^{2}=M$ and then $E(u,\psi,0)\geq J_{M}$. Moreover, reasoning as in \eqref{last_convergence} we get
\begin{equation}\label{cv1}
\int_{\R^{d}}\left(\sigma_{1}\star\int_{\R^{n}}\sigma_{2}\varphi_{\nu}\ud z\right)|v_{\nu}|^{2}\ud x\,\underset{\nu\to+\infty}{\to}\,\int_{\R^{d}}\left(\sigma_{1}\star\int_{\R^{n}}\sigma_{2}\psi\ud z\right)|u|^{2}\ud x,
\end{equation}
which allows us to justify that $(u,\psi)$ lies in $S_{M}$:
\bgroup
\multlinegap0pt
\begin{multline*}
J_{M}=\underset{\nu\to+\infty}{\lim}\,E(v_{\nu},\varphi_{\nu},0)\geq\underset{\nu\to+\infty}{\liminf}\Bigl(\frac{1}{2}\|\nabla_{x}v_{\nu}\|_{L_{x}^{2}}^{2}\Bigr)\\
+\underset{\nu\to +\infty}{\liminf}\left(\int_{\R^{d}}\left(\sigma_{1}\star\int_{\R^{n}}\sigma_{2}\varphi_{\nu}\ud z\right)|v_{\nu}|^{2}\ud x\right)+\underset{\nu\to+\infty}{\liminf}\Bigl(\frac{1}{2}\|\nabla_{z}\varphi_{\nu}\|_{L_{x}^{2}L_{z}^{2}}^{2}\Bigr)\geq E(u,\psi,0).
\end{multline*}
\egroup
In order to conclude the proof it only remains to justify the strong convergence of $(v_{\nu},\varphi_{\nu})_{\nu\in\mathbb N}$ to $(u,\psi)$ in $H_{x}^{1}\times L_{x}^{2}\overbigdot{H}_{z}^{1}$. We already know that this convergence holds weakly. We combine
\[
E(u,\psi,0)=J_{M}=\underset{\nu\to+\infty}{\lim}\,E(v_{\nu},\varphi_{\nu},0)
\]
and \eqref{cv1} to deduce that
\[
\frac{1}{2}\|\nabla_{x}v_{\nu}\|_{L_{x}^{2}}^{2}+\frac{1}{2}\|\nabla_{z}\varphi_{\nu}\|_{L_{x}^{2}L_{z}^{2}}^{2}\,\underset{\nu\to+\infty}{\to}\,\frac{1}{2}\|\nabla_{x}u\|_{L_{x}^{2}}^{2}+\frac{1}{2}\|\nabla_{z}\psi\|_{L_{x}^{2}L_{z}^{2}}^{2},
\]
holds, which ends the proof.
\end{ProofOf}

\section{Strengthened orbital stability: approach by linearization}\label{OrbitalStability2}

In this Section, we explain how Lemma~\ref{Lemma_L-} and Lemma~\ref{Lemma_LL+} imply Theorem~\ref{orbital_stab_2}.

\subsubsection*{Step 1} The first step of the proof consists in checking that, up to the invariants of the equation, any $v\in H_{x}^{1}$ close enough to $Q$ satisfies the orthogonality conditions \eqref{ortho1}--\eqref{ortho2}. For that purpose, let us introduce the function $F:H_{x}^{1}\times\R^{d+1}\to\R^{d+1}$ defined by
\begin{align*}
\ds F_{j}\left(v,(y,\theta)\right)&=\bigl\langle\Re e^{-i\theta}v(\cdot+y),\partial_{x_{j}}Q\bigr\rangle_{L_{x}^{2}},\qquad j=1,\ldots,d,\\
\ds F_{d+1}\left(v,(y,\theta)\right)&=\bigl\langle\Im e^{-i\theta}v(\cdot+y),Q\bigr\rangle_{H_{x}^{1}}.
\end{align*}
Direct computations show that $F\left(Q,(0,0)\right)=0$ and $\mathrm{D}_{y,\theta}\,F\left(Q,(0,0)\right)$ is an invertible diagonal matrix (indeed $\partial_{y_{j}}F_{j}\left(Q,(0,0)\right)=\|\partial_{x_{j}}Q\|_{L_{x}^{2}}^{2}$ and $\partial_{\theta}F_{d+1}\left(Q,(0,0)\right)=-\|Q\|_{H_{x}^{1}}^{2}$). The implicit function theorem provides the existence of $\varepsilon_{0}>0$ and a $C^{1}$\nobreakdash-diffeo\-morphism $G:B_{H_{x}^{1}}(Q,2\varepsilon_{0})\to U_{\varepsilon_{0}}\subset\R^{d+1}$, $G(v)=(x,\gamma)$ such that for every $v\in B_{H_{x}^{1}}(Q,2\varepsilon_{0})$ and every $(y,\theta)\in U_{\varepsilon_{0}}$, $F\left(v,(y,\theta)\right)=0$ if and only if $(y,\theta)=G(v)$. Moreover, since
\[
|(x,\gamma)|=|G(v)-G(Q)|\leq\Bigl(\sup_{v\in B_{H_{x}^{1}}(Q,2\varepsilon_{0})}\,\|\mathrm{D}_{v}G\|\Bigr)\|v-Q\|_{H_{x}^{1}},
\]
for every $\varepsilon\in(0,\varepsilon_{0})$ there exists $\eta(\varepsilon)>0$ such that
\[
\left\|(v-Q,\varphi-\Psi)\right\|_{\mathscr{H}}^{2}+\frac{1}{c^{2}}\|\chi\|_{L_{x}^{2}L_{z}^{2}}^{2}\leq\eta(\varepsilon)^{2}
\]
implies for $(x,\gamma)=G(v)$,
\[
\left\|(e^{-i\gamma}v(\cdot+x)-Q,\varphi(\cdot+x)-\Psi)\right\|_{\mathscr{H}}^{2}+\frac{1}{c^{2}}\|\chi\|_{L_{x}^{2}L_{z}^{2}}^{2}\leq\varepsilon^{2}.
\]

\subsubsection*{Step 2} In this second step we show that, if for a given time $t_{0}\in[0,+\infty)$, there exists $(x_{0},\gamma_{0})\in\R^{d+1}$ such that $v=e^{-i\gamma_{0}}u(t_{0},\cdot+x_{0})$ satisfies the orthogonality conditions \eqref{ortho1}--\eqref{ortho2} and the estimate\vspace*{-3pt}
\[
\left\|(e^{-i\gamma_{0}}u(t_{0},\cdot+x_{0})-Q,\psi(t_{0},\cdot+x_{0})-\Psi)\right\|_{\mathscr{H}}^{2}+\frac{1}{c^{2}}\|\chi(t_{0})\|_{L_{x}^{2}L_{z}^{2}}^{2}\leq\varepsilon^{2}<\varepsilon_{0}^{2},
\]
then there exists a time $T^{\star}>t_{0}$ and two functions $x(t)$ and $\gamma(t)$ continuous in time such that $(x(t_{0}),\gamma(t_{0}))=(x_{0},\gamma_{0})$ and, for every $t\in[t_{0},T^{\star})$,
\begin{enumeratei}
\item\label{enum2:i}
$(x(t)-x_{0},\gamma(t)-\gamma_{0})\in U_{\varepsilon_{0}}$,
\item\label{enum2:ii}
$v=e^{-i\gamma(t)}u(t,\cdot+x(t))$ satisfies the orthogonality conditions \eqref{ortho1}--\eqref{ortho2},
\item\label{enum2:iii}
$\ds\bigl\|(e^{-i\gamma(t)}u(t,\cdot+x(t))-Q,\psi(t,\cdot+x(t))-\Psi)\bigr\|_{\mathscr{H}}^{2}+\frac{1}{c^{2}}\|\chi(t)\|_{L_{x}^{2}L_{z}^{2}}^{2}\leq\varepsilon^{2}.$
\end{enumeratei}
First, thanks to the time continuity of $t\mto(e^{-i\gamma_{0}}u(t,\cdot+x_{0}),\psi(t,\cdot+x_{0}))\in\mathscr{H}$, there exists a time $T^{\star}>t_{0}$ such that for every $t\in[t_{0},T^{\star})$\vspace*{-3pt}
\[
\bigl\|\bigl(e^{-i\gamma_{0}}u(t,\cdot+x_{0})-Q,\psi(t,\cdot+x_{0})-\Psi\bigr)\bigr\|_{\mathscr{H}}^{2}\leq 4\varepsilon^{2}<4\varepsilon_{0}^{2}.
\]
Next, for every $t\in[t_{0},T^{\star})$ we can apply the first step to $v=e^{-i\gamma_{0}}u(t,\cdot+x_{0})$ and we obtain the existence of $x(t)$ and $\gamma(t)$ such that $(x(t_{0}),\gamma(t_{0}))=(x_{0},\gamma_{0})$ and such that~\eqref{enum2:i} and~\eqref{enum2:ii} hold. Moreover the continuity of $t\mto e^{-i\gamma_{0}}u(t,\cdot+x_{0})$ implies the continuity of $t\mto x(t)$ and $t\mto\gamma(t)$. We notice also that we can extend by continuity $x(t)$ and $\gamma(t)$ at time $T^{\star}$ and this extension is such that $v=e^{-i\gamma(T^{\star})}u(T^{\star},\cdot+x(T^{\star}))$ still satisfies the orthogonality conditions \eqref{ortho1}--\eqref{ortho2}.

We can now apply Lemma~\ref{Lemma_L-} and \ref{Lemma_LL+} as follows. Thanks to the conservation of mass and energy and
to the invariance by translation and phase of these quantities we get\vspace*{-3pt}
\begin{align*}
W(u_{0},\psi_{0},\chi_{0})&=W(u(t),\psi(t),\chi(t))
=W\bigl(e^{-i\gamma(t)}u(t,\cdot+x(t)),\psi(t,\cdot+x(t)),\chi(t)\bigr)\\
&=W(Q+u^{\varepsilon}(t),\Psi+\psi^{\varepsilon}(t),\chi(t)),
\end{align*}
where\vspace*{-5pt}
\[
u^{\varepsilon}(t)=e^{-i\gamma(t)}u(t,\cdot+x(t))-Q\hspace{0.5cm}\text{ and }\hspace{0.5cm}\psi^{\varepsilon}(t)=\psi(t,\cdot+x(t))-\Psi.
\]
We make use of the decomposition \eqref{decomp_W} combined with Lemma~\ref{Lemma_L-} and \ref{Lemma_LL+}; we obtain\vspace*{-5pt}
\begin{multline*}
\bar{\nu}\|(\Re u^{\varepsilon},\psi^{\varepsilon})\|_{\mathscr{H}}^{2}+\mu\|\Im u^{\varepsilon}\|_{H_{x}^{1}}^{2}+\frac{1}{2c^{2}}\|\chi(t)\|_{L_{x}^{2}L_{z}^{2}}^{2}\\[
-3pt]
\leq W(u_{0},\psi_{0},\chi_{0})-W(Q,\Psi,0)+\frac{1}{\bar{\nu}}\Bigl(\left|\langle \Re u^{\varepsilon},Q\rangle_{L_{x}^{2}}\right|^{2}+\sum_{j=1}^{d}\left|\langle \Re u^{\varepsilon},\partial_{x_{j}}Q\rangle_{L_{x}^{2}}\right|^{2}\Bigr)\\[
-2pt]
+\frac{1}{\mu}\left|\langle \Im u^{\varepsilon},Q\rangle_{H_{x}^{1}}\right|^{2}-\int_{\R^{d}}\left(\sigma_{1}\star\int_{\R^{n}}\sigma_{2}\psi^{\varepsilon}(t)\ud z\right)|u^{\varepsilon}(t)|^{2}\ud x.
\end{multline*}
Since $e^{-i\gamma(t)}u(t,\cdot+x(t))$ and $Q$ satisfy the orthogonality conditions \eqref{ortho1}--\eqref{ortho2} we know that $u^{\varepsilon}$ also satisfies these conditions. Moreover $\|Q\|_{L_{x}^{2}}=\|u(t)\|_{L_{x}^{2}}=\|u^{\varepsilon}+Q\|_{L_{x}^{2}}$ leads to\vspace*{-3pt}
\[
\|Q\|_{L_{x}^{2}}^{2}=\|u^{\varepsilon}\|_{L_{x}^{2}}^{2}+\|Q\|_{L_{x}^{2}}^{2}+2\langle\Re u^{\varepsilon},Q\rangle_{L_{x}^{2}} \hspace{0.25cm}\text{and then}\hspace{0.25cm} \langle\Re u^{\varepsilon},Q\rangle_{L_{x}^{2}}=-\frac{1}{2}\|u^{\varepsilon}\|_{L_{x}^{2}}^{2},
\]
which implies\vspace*{-3pt}
\[
\left|\langle \Re u^{\varepsilon},Q\rangle_{L_{x}^{2}}\right|^{2}\leq\frac{1}{4}\|u^{\varepsilon}\|_{L_{x}^{2}}^{4}\leq 4\,\varepsilon^{4}.
\]
We also get
\bgroup
\multlinegap0pt
\begin{multline*}
\left|\int_{\R^{d}}\biggl(\sigma_{1}\star\int_{\R^{n}}\!\!\sigma_{2}\psi^{\varepsilon}(t)\ud z\biggr)|u^{\varepsilon}(t)|^{2}\ud x\right|\leq\|\sigma_{1}\|_{L_{x}^{2}}\|\sigma_{2}\|_{L_{z}^{2n/(n+2)}}\|\psi^{\varepsilon}(t)\|_{L_{x}^{2}\overbigdot{H}_{z}^{1}}\|u^{\varepsilon}(t)\|_{L_{x}^{2}}^{2}\\
\leq\|\sigma_{1}\|_{L_{x}^{2}}\|\sigma_{2}\|_{L_{z}^{2n/(n+2)}}\|(u^{\varepsilon}(t),\psi^{\varepsilon}(t))\|_{\mathscr{H}}^{3}\leq 8\,\|\sigma_{1}\|_{L_{x}^{2}}\|\sigma_{2}\|_{L_{z}^{2n/(n+2)}}\;\varepsilon^{3}.
\end{multline*}
\egroup

Gathering these estimates leads eventually to (we recall that the inequality $W(u_{0},\psi_{0},\chi_{0})-W(Q,\Psi,0)\leq\delta(\varepsilon)$ holds)
\begin{multline*}
\|(\Re u^{\varepsilon},\psi^{\varepsilon})\|_{\mathscr{H}}^{2}+\|\Im u^{\varepsilon}\|_{H_{x}^{1}}^{2}+\frac{1}{c^{2}}\|\chi(t)\|_{L_{x}^{2}L_{z}^{2}}^{2}\\
\leq\frac{1}{\min\big(\bar{\nu},\mu,\sfrac{1}{2}\big)}\,\Bigl(\delta(\varepsilon)+\frac{4}{\bar{\nu}}\,\varepsilon^{4}+8\,\|\sigma_{1}\|_{L_{x}^{2}}\|\sigma_{2}\|_{L_{z}^{2n/(n+2)}}\;\varepsilon^{3}\Bigr).
\end{multline*}
By taking
\[
\delta(\varepsilon)=\frac{\varepsilon^{2}}{2\,\min\big(\bar{\nu},\mu,\sfrac{1}{2}\big)},
\]
and possibly at the price of picking a smaller $\varepsilon_{0}$, we eventually obtain \eqref{enum2:iii} for every $t\in[t_{0},T^{\star}]$.

\subsubsection*{Conclusion} We apply the first step with $v=u_{0}$, which ensures the existence of $x(0)$ and $\gamma(0)$ such that we can apply the second step at time $t=0$. Thus, since $T^{\star}>0$ and since we took care to justify that the conclusions of second step is also valid at time $t=T^{\star}$, a classical argument on connected space allows us to conclude that $T^{\star}=+\infty$.\qed

\section{Coercivity of \texorpdfstring{$\mathcal{L}_{+}$}{L}: proof of Lemma~\ref{Lemma_LL+}}\label{Proof_Lemma_LL+}

This section is dedicated to the proof of Lemma~\ref{Lemma_LL+}, which is a key ingredient of the proof of Theorem~\ref{orbital_stab_2}.
We assume that $\sigma_1$ is admissible:
the kernel of $\mathcal{L}_{+}$ can be identified by using
the fact that the conclusions of
Lemma~\ref{Lemma_L+} apply. Indeed, since $(f,\psi)^{t}\in\Ker(\mathcal{L}_{+})$ implies
\[
-\frac{1}{2}\Delta_{z}\psi+\sigma_{2}\,(\sigma_{1}\star Qf)=0,
\]
we can express $\psi$ in term of $f$ as follows: $\psi=2\Gamma\,(\sigma_{1}\star Qf)$. Moreover the relation
\begin{equation}\label{relation_LL+_L+}
\mathcal{L}_{+}\begin{pmatrix} f \\ 2\Gamma\,(\sigma_{1}\star Qf) \end{pmatrix}=\begin{pmatrix} L_{+}f\\0 \end{pmatrix}
\end{equation}
allows us to identify the kernel of $\mathcal{L}_{+}$
by using the knowledge of the kernel of $L_{+}$: we~eventually get
\[
\Ker(\mathcal{L}_{+})=\mathrm{Span}\{(\partial_{x_{j}}Q,\partial_{x_{j}}\Psi)^{t}\sep j=1,\ldots,d\}.
\]
In order to prove the coercivity relations \eqref{coercivity3}, we need the following two lemmas.

\begin{lemma}\label{positivity_LL+}
For every $(f,\psi)\in\mathscr{H}$ such that $\langle f,Q\rangle_{L_{x}^{2}}=0$, we have
\[
\left\langle\mathcal{L}_{+}\begin{pmatrix} f\\\psi \end{pmatrix},\begin{pmatrix} f\\\psi \end{pmatrix}\right\rangle_{L_{x}^{2}\times L_{x}^{2}L_{z}^{2}}\geq 0.
\]
Moreover, since $\Ker(\mathcal{L}_{+})=\{(\partial_{x_{j}}Q,\partial_{x_{j}}\Psi)^{t}\sep j=1,\ldots,d\}$ and $\langle \partial_{x_{j}}Q,Q\rangle_{L_{x}^{2}}=0$, we~know that this inequality cannot be strict.
\end{lemma}

\begin{lemma}\label{compactness_LL+}
Let $(f_{\nu},\psi_{\nu})_{\nu\in\mathbb N}$ be a bounded sequence of $\mathscr{H}$ which converges weakly to $(\bar{f},\bar{\psi})$ in $\mathscr{H}$. Then, up to a sub-sequence if needed, we have the following two convergences:
\begin{equation}\label{lcv1}
\int_{\R^{d}}\left(\sigma_{1}\star\int_{\R^{n}}\sigma_{2}\Psi\ud z\right)|f_{\nu}|^{2}\ud x\,\underset{\nu\to+\infty}{\to}\,\int_{\R^{d}}\left(\sigma_{1}\star\int_{\R^{n}}\sigma_{2}\Psi\ud z\right)|\bar{f}|^{2}\ud x
\end{equation}
and
\begin{equation}\label{lcv2}
\int_{\R^{d}}\left(\sigma_{1}\star\int_{\R^{n}}\sigma_{2}\psi_{\nu}\ud z\right)Qf_{\nu}\ud x\,\underset{\nu\to +\infty}{\to}\,\int_{\R^{d}}\left(\sigma_{1}\star\int_{\R^{n}}\sigma_{2}\bar{\psi}\ud z\right)Q\bar{f}\ud x.
\end{equation}
\end{lemma}

\begin{ProofOf}{Lemma~\ref{positivity_LL+}}
Let $f$ be a real valued function of $H_{x}^{1}$ such that $\langle f,Q\rangle_{L_{x}^{2}}=0$, let $\psi$ be a function of $L_{x}^{2}\overbigdot{H}_{z}^{1}$ and let $u$ be the function defined on $\R$ by
\[
u(s)=\frac{\|Q\|_{L_{x}^{2}}}{\|Q+sf\|_{L_{x}^{2}}}\,(Q+sf).
\]
One can check that $u(s)$ is a real valued function of $H_{x}^{1}$ and $\|u(s)\|_{L_{x}^{2}}=\|Q\|_{L_{x}^{2}}$ for every $s\in\R$, $u$ is smooth, $u(0)=Q$ and
\[
u'(0)=f-\frac{\langle f,Q\rangle_{L_{x}^{2}}}{\|Q\|_{L_{x}^{2}}^{2}}\,Q=f.
\]
Since $(Q,\Psi,0)$ is a minimizer of $J_{M}$, we know that for every $s\in\R$, $W(Q,\Psi,0)\leq W(u(s),\Psi+s\psi,0)$. Moreover \eqref{decomp_W} leads to
\bgroup
\multlinegap0pt
\begin{multline*}
0\leq W(u(s),\Psi+s\psi,0)-W(Q,\Psi,0)=\left\langle\mathcal{L}_{+}\begin{pmatrix} u(s)-Q\\s\psi \end{pmatrix},\begin{pmatrix} u(s)-Q\\s\psi \end{pmatrix}\right\rangle_{L_{x}^{2}\times L_{x}^{2}L_{z}^{2}}\\
+\int_{\R^{d}}\left(\sigma_{1}\star\int_{\R^{n}}\sigma_{2}s\psi\ud z\right)|u(s)-Q|^{2}\ud x.
\end{multline*}
\egroup
Since $u(s)-Q=u(s)-u(0)=sf+o(s)$ (when $s$ goes to $0$), we eventually obtain
\[
0\leq s^{2}\left\langle\mathcal{L}_{+}\begin{pmatrix} f\\\psi \end{pmatrix},\begin{pmatrix} f\\\psi \end{pmatrix}\right\rangle_{L_{x}^{2}\times L_{x}^{2}L_{z}^{2}}+o(s^{2}),
\]
which concludes the proof.
\end{ProofOf}

\begin{ProofOf}{Lemma~\ref{compactness_LL+}}
The proof uses in several places a basic result of integration theory,
consequence of
Egoroff's theorem \cite[Prop.\,3.9]{Tay}: if a sequence $(g_{\nu})_{\nu\in\mathbb N}\subset L^{p}(\R^{d})$ converges weakly to some $\bar{g}$ in $L^{p}(\R^{d})$ where $1\leq p<+\infty$ and if this sequence converges also a.e.\ to some $g$, then $\bar{g}=g$.

Here, the sequence $(f_{\nu})_{\nu\in \mathbb N}$ is bounded in $H^{1}(\R^{d})$
and the compact embedding $H^{1}(\Omega)\to L^{2}(\Omega)$
which holds for any bounded open set $\Omega\subset\R^{d}$ implies that, up to a sub-sequence, $(f_{\nu})_{\nu\in\mathbb N}$ converges strongly to $\bar{f}$ in $L^{2}(\Omega)$ and thus converges, up to a further sub-sequence, a.e.\ in $\Omega$ to $\bar{f}$.
A diagonal argument yields the a.e.\ convergence of $(f_{\nu})_{\nu\in\mathbb N}$ to $\bar{f}$ in $\R^{d}$.
Moreover, by using the homogeneous Sobolev embedding in dimension $d=3$, the boundedness of $(f_{\nu})_{\nu\in\mathbb N}$ in $H_{x}^{1}$ implies its boundedness in $L_{x}^{2}$ and $L_{x}^{6}$ and, by interpolation, in any $L_{x}^{p}$ with $2\leq p\leq 6$. Consequently, the sequence $(|f_{\nu}|^{2})_{\nu\in\mathbb N}$ is bounded in $L_{x}^{3}$ and, up to a sub-sequence, converges weakly in $L_{x}^{3}$ to some $g$. Since this sequence converges also a.e.\ to $|\bar{f}|^{2}$, we have indeed $g=|\bar{f}|^{2}$.

To prove \eqref{lcv1} we proceed as follows.
Since $\Psi=\Gamma\,\sigma_{1}\star Q^{2}$ with $Q$ lying in the Schwartz class, the weak convergence of $(|f_{\nu}|^{2})_{n\in \mathbb N}$ to $|f|^{2}$ in $L_{x}^{3}$ yields
\begin{multline*}
\int\left(\sigma_{1}\star\int\sigma_{2}\Psi\ud z\right)|f_{\nu}|^{2}\ud x=-\kappa\int\left(\Sigma\star Q^{2}\right)|f_{\nu}|^{2}\ud x\\
\underset{\nu\to+\infty}{\to}\,-\kappa\int\left(\Sigma\star Q^{2}\right)|\bar{f}|^{2}\ud x=\int\left(\sigma_{1}\star\int\sigma_{2}\Psi\ud z\right)|\bar{f}|^{2}\ud x.
\end{multline*}

We turn to \eqref{lcv2}. We split
\begin{multline*}
\int\left(\sigma_{1}\star\int\sigma_{2}\psi_{\nu}\ud z\right)Qf_{\nu}\ud x=\iint\sigma_{2}\left(\sigma_{1}\star Qf_{\nu}\right)\psi_{\nu}\ud x\ud z\\
=\iint\sigma_{2}\left(\sigma_{1}\star Q(f_{\nu}-\bar{f})\right)\psi_{\nu}\ud x\ud z+\iint\sigma_{2}\left(\sigma_{1}\star Q\bar{f}\right)\psi_{\nu}\ud x\ud z.
\end{multline*}
The weak convergence of $(\psi_{\nu})_{\nu\in\mathbb N}$ to $\bar{\psi}$ in $L_{x}^{2}\overbigdot{H}_{z}^{1}$ (note that $\sigma_{2}$ smooth and $n\geq 3$ imply $\sigma_{2}\in\overbigdot{H}_{z}^{-1}$) directly implies that
the second term of the right hand side converges to $\int(\sigma_{1}\star\int\sigma_{2}\bar{\psi}\ud z)Q\bar{f}\ud x$.
It only remains to prove that the first term of the right hand side converges to $0$.
To this end, we are going to show
that $(Qf_{\nu})_{\nu\in\mathbb N}$ converges strongly to $Q\bar{f}$ in $L_{x}^{3/2}$. Indeed, $(|f_{\nu}|^{3/2})_{\nu\in\mathbb N}$ is bounded in $L_{x}^{2}$ and, up to a sub-sequence it converges weakly to $g=|\bar{f}|^{3/2}$ in $L_{x}^{2}$. Since $Q^{3/2}\in L_{x}^{2}$, we get $\|Qf_{\nu}\|_{L_{x}^{3/2}}\to\|Q\bar{f}\|_{L_{x}^{3/2}}$ as $\nu\to \infty$. Moreover the sequence $(Qf_{\nu})_{\nu\in\mathbb N}$ is also bounded in $L_{x}^{3/2}$ and, up to a further sub-sequence if needed, it converges weakly to $Q\bar{f}$ in $L_{x}^{3/2}$. Thus we get the announced strong convergence. We combine this strong convergence with the boundedness of $(\psi_{\nu})_{n\in \mathbb N}$ in $L_{x}^{2}\overbigdot{H}_{z}^{1}$ and we conclude as follows:
\bgroup
\multlinegap0pt
\begin{multline*}
\left|\iint\sigma_{2}\left(\sigma_{1}\star Q(f_{\nu}-\bar{f})\right)\psi_{\nu}\ud x\ud z\right|\leq\|\sigma_{2}\|_{L_{z}^{2n/(n+2)}}\|\psi_{\nu}\|_{L_{x}^{2}\overbigdot{H}_{z}^{1}}\|\sigma_{1}\star Q(f_{\nu}-\bar{f})\|_{L_{x}^{2}}\\
\leq\|\sigma_{2}\|_{L_{z}^{2n/(n+2)}}\|\psi_{\nu}\|_{L_{x}^{2}\overbigdot{H}_{z}^{1}}\|\sigma_{1}\|_{L_{x}^{6/5}}\|Qf_{\nu}-Q\bar{f}\|_{L_{x}^{3/2}}\,\underset{\nu\to+\infty}{\to}\,0.\qedhere
\end{multline*}\egroup
\end{ProofOf}

We are now able to prove the coercivity relation \eqref{coercivity3}.

\begin{ProofOf}{\eqref{coercivity3}}
We argue by contradiction, assuming the existence of a sequence of positive numbers $(\tilde{\nu}_{k})_{k\in\mathbb N}$ which converges to $0$ and the existence of a sequence $(f_{k},\psi_{k})_{k\in\mathbb N}$ in $\mathscr{H}$ such that for every $k$,
\begin{multline}\label{non_coercivity}
\left\langle\mathcal{L}_{+}\begin{pmatrix} f_{k}\\\psi_{k} \end{pmatrix},\begin{pmatrix} f_{k}\\\psi_{k} \end{pmatrix}\right\rangle_{L_{x}^{2}\times L_{x}^{2}L_{z}^{2}}
\\
<\tilde{\nu}_{k}\|(f_{k},\psi_{k})\|_{\mathscr{H}}^{2}-\frac{1}{\tilde{\nu}_{k}}\Bigl(\left|\langle f_{k},Q\rangle_{L_{x}^{2}}\right|^{2}+\sum\limits_{j=1}^{d}\left|\langle f_{k},\partial_{x_{j}}Q\rangle_{L_{x}^{2}}\right|^{2}\Bigr).
\end{multline}
We can assume that $\| (f_{k},\psi_{k})\|_{\mathscr{H}}=1$ and thus, that there exists $\bar{f}\in H_{x}^{1}$ and $\bar{\psi}\in L_{x}^{2}\overbigdot{H}_{z}^{1}$ such that $(f_{k})_{k\in\mathbb N}$ converges weakly to $\bar{f}$ in $H_{x}^{1}$ and $(\psi_{k})_{k\in\mathbb N}$ converges weakly to $\bar{\psi}$ in $L_{x}^{2}\overbigdot{H}_{z}^{1}$. On the one hand, thanks to the weak convergence of $(f_{k})_{k\in\mathbb N}$, we get
\[
\langle f_{k},Q\rangle_{L_{x}^{2}}\,\underset{k\to+\infty}{\to}\,\langle\bar{f},Q\rangle_{L_{x}^{2}} \hspace{0.25cm}\text{ and }\hspace{0.25cm} \langle f_{k},\partial_{x_{j}}Q\rangle_{L_{x}^{2}}\,\underset{k\to+\infty}{\to}\,\langle\bar{f},\partial_{x_{j}}Q\rangle_{L_{x}^{2}},
\]
while on the other hand \eqref{non_coercivity} implies
\[
0\leq\left|\langle f_{k},Q\rangle_{L_{x}^{2}}\right|^{2}+\sum\limits_{j=1}^{d}\left|\langle f_{k},\partial_{x_{j}}Q\rangle_{L_{x}^{2}}\right|^{2}<\bar{\nu}_{k}^{2}-\bar{\nu}_{k}\left\langle\mathcal{L}_{+}\begin{pmatrix} f_{k}\\\psi_{k} \end{pmatrix},\begin{pmatrix} f_{k}\\\psi_{k} \end{pmatrix}\right\rangle_{L_{x}^{2}\times L_{x}^{2}L_{z}^{2}}\underset{k\to+\infty}{\to}0,
\]
bearing in mind that $\langle\mathcal{L}_{+} h,h\rangle\leq K\|h\|_{\mathscr{H}}^{2}$. We eventually obtain $\langle\bar{f},Q\rangle_{L_{x}^{2}}=0$ and $\langle\bar{f},\partial_{x_{j}}Q\rangle_{L_{x}^{2}}=0$. Knowing that $\bar{f}$ is orthogonal to $Q$, we can apply Lemma~\ref{positivity_LL+} in order to obtain
\[
\left\langle\mathcal{L}_{+}\begin{pmatrix} \bar{f}\\\bar{\psi} \end{pmatrix},\begin{pmatrix} \bar{f}\\\bar{\psi} \end{pmatrix}\right\rangle_{L_{x}^{2}\times L_{x}^{2}L_{z}^{2}}\geq 0.
\]
On the other hand, the relation
\begin{multline*}
\left\langle\mathcal{L}_{+}\begin{pmatrix} f_{k}\\\psi_{k} \end{pmatrix},\begin{pmatrix} f_{k}\\\psi_{k} \end{pmatrix}\right\rangle_{L_{x}^{2}\times L_{x}^{2}L_{z}^{2}}
\\
=\frac{1}{2}\|\nabla_{x}f_{k}\|_{L_{x}^{2}}^{2}+\omega\|f_{k}\|_{L_{x}^{2}}^{2}+\int\left(\sigma_{1}\star\int\sigma_{2}\Psi\ud z\right)|f_{k}|^{2}\ud x\\
\qquad +2\int\left(\sigma_{1}\star\int\sigma_{2}\psi_{k}\ud z\right)Qf_{k}\ud x+\frac{1}{2}\|\nabla_{z}\psi_{k}\|_{L_{x}^{2}L_{z}^{2}}^{2},
\end{multline*}
coupled with Lemma~\ref{compactness_LL+} and \eqref{non_coercivity} leads to
\begin{align*}
\left\langle\mathcal{L}_{+}\begin{pmatrix} \bar{f}\\\bar{\psi} \end{pmatrix},\begin{pmatrix} \bar{f}\\\bar{\psi} \end{pmatrix}\right\rangle_{L_{x}^{2}\times L_{x}^{2}L_{z}^{2}}&\leq\underset{k\to+\infty}{\liminf}\,\left\langle\mathcal{L}_{+}\begin{pmatrix} f_{k}\\\psi_{k} \end{pmatrix},\begin{pmatrix} f_{k}\\\psi_{k} \end{pmatrix}\right\rangle_{L_{x}^{2}\times L_{x}^{2}L_{z}^{2}}\\
&\leq\underset{k\to+\infty}{\limsup}\,\left\langle\mathcal{L}_{+}\begin{pmatrix} f_{k}\\\psi_{k} \end{pmatrix},\begin{pmatrix} f_{k}\\\psi_{k} \end{pmatrix}\right\rangle_{L_{x}^{2}\times L_{x}^{2}L_{z}^{2}}\\
&\leq\underset{k\to +\infty}{\limsup}\,\biggl\{\frac{1}{\bar{\nu}_{k}}\Bigl(\left|\langle f_{k},Q\rangle_{L_{x}^{2}}\right|^{2}
+\sum\limits_{j=1}^{d}\left|\langle f_{k},\partial_{x_{j}}Q\rangle_{L_{x}^{2}}\right|^{2}\Bigr)
\\
&\hspace*{3cm}+\left\langle\mathcal{L}_{+}\begin{pmatrix} f_{k}\\\psi_{k} \end{pmatrix},\begin{pmatrix} f_{k}\\\psi_{k} \end{pmatrix}\right\rangle_{L_{x}^{2}\times L_{x}^{2}L_{z}^{2}}\biggr\}\\
&\leq\underset{k\to+\infty}{\limsup}\,\bar{\nu}_{k}=0.
\end{align*}
We eventually deduce
\begin{equation}\label{cv_LL+}
\underset{k\to+\infty}{\lim}\,\left\langle\mathcal{L}_{+}\begin{pmatrix} f_{k}\\\psi_{k} \end{pmatrix},\begin{pmatrix} f_{k}\\\psi_{k} \end{pmatrix}\right\rangle_{L_{x}^{2}\times L_{x}^{2}L_{z}^{2}}\,=\, \left\langle\mathcal{L}_{+}\begin{pmatrix} \bar{f}\\\bar{\psi} \end{pmatrix},\begin{pmatrix} \bar{f}\\\bar{\psi} \end{pmatrix}\right\rangle_{L_{x}^{2}\times L_{x}^{2}L_{z}^{2}}=0
\end{equation}
and thus $(\bar{f},\bar{\psi})$ is a minimizer of
\begin{equation}\label{min_LL+}
\underset{\langle f,Q\rangle_{L_{x}^{2}}=0}{\inf}\;\left\langle\mathcal{L}_{+}\begin{pmatrix} f\\\psi \end{pmatrix},\begin{pmatrix} f\\\psi \end{pmatrix}\right\rangle_{L_{x}^{2}\times L_{x}^{2}L_{z}^{2}}.
\end{equation}
We can now conclude as follows. First of all, the relation \eqref{cv_LL+} coupled with Lemma~\ref{compactness_LL+} leads to the norm convergence
\[
\frac{1}{2}\|\nabla_{x}f_{k}\|_{L_{x}^{2}}^{2}+\omega\|f_{k}\|_{L_{x}^{2}}^{2}+\frac{1}{2}\|\psi_{k}\|_{L_{x}^{2}\overbigdot{H}_{z}^{1}}^{2}\,\underset{k\to +\infty}{\to}\,\frac{1}{2}\|\nabla_{x}\bar{f}\|_{L_{x}^{2}}^{2}+\omega\|\bar{f}\|_{L_{x}^{2}}^{2}+\frac{1}{2}\|\bar{\psi}\|_{L_{x}^{2}\overbigdot{H}_{z}^{1}}^{2}.
\]
It implies the strong convergence of $(f_{k},\psi_{k})_{k\in\mathbb N}$ to $(\bar{f},\bar{\psi})$ in $\mathscr{H}$. In particular we know that $\|(\bar{f},\bar{\psi})\|_{\mathscr{H}}=1$. Second of all, $(\bar{f},\bar{\psi})$ is a minimizer of \eqref{min_LL+} and the Euler-Lagrange relation ensures the existence of a real number $\lambda$ such that
\[
\mathcal{L}_{+}\begin{pmatrix} \bar{f}\\\bar{\psi} \end{pmatrix}=\lambda\begin{pmatrix} Q\\0 \end{pmatrix}.
\]
The second component of this vectorial relation leads to $\bar{\psi}=2\Gamma\,(\sigma_{1}\star Q\bar{f})$. From this relation we obtain the contradiction as follows: owing to \eqref{relation_LL+_L+}, the fact that $\sigma_1$ is admissible, and since $\bar{f}$ is orthogonal to $Q$ and $\partial_{x_{j}}Q$, we get
\begin{align*}
0&=\left\langle\mathcal{L}_{+}\begin{pmatrix} \bar{f}\\\bar{\psi} \end{pmatrix},\begin{pmatrix} \bar{f}\\\bar{\psi} \end{pmatrix}\right\rangle_{L_{x}^{2}\times L_{x}^{2}L_{z}^{2}}=\left\langle\begin{pmatrix} L_{+}\bar{f}\\0 \end{pmatrix},\begin{pmatrix} \bar{f}\\\bar{\psi} \end{pmatrix}\right\rangle_{L_{x}^{2}\times L_{x}^{2}L_{z}^{2}}\\
&=\left\langle L_{+}\bar{f},\bar{f}\right\rangle_{L_{x}^{2}}\geq\nu\|\bar{f}\|_{H_{x}^{1}}^{2}-\frac{1}{\nu}\Bigl(\left|\langle \bar{f},Q\rangle_{L_{x}^{2}}\right|^{2}+\sum\limits_{j=1}^{d}\left|\langle \bar{f},\partial_{x_{j}}Q\rangle_{L_{x}^{2}}\right|^{2}\Bigr)=\nu\|\bar{f}\|_{H_{x}^{1}}^{2}.
\end{align*}
Thus $(\bar{f},\bar{\psi})=(0,0)$, which contradicts $\|(\bar{f},\bar{\psi})\|_{\mathscr{H}}=1$.
\end{ProofOf}

\section{Perturbation analysis: proof of Proposition~\ref{perturbative_L+}}\label{Perturbative}

In this section, since there is no ambiguity, we will use the following shorthand notations, see Definition~\ref{defKUpM}, $H^{\varepsilon}=H^{\Sigma^{\varepsilon}}$, $K_{M}^{\varepsilon}=K_{M}^{\Sigma^{\varepsilon}}$, $L_{+}^{\varepsilon}=L_{+}(\Sigma^{\varepsilon},Q^{\varepsilon})$, $H^{0}=H^{\Sigma^{0}}$, $K_{M}^{0}=K_{M}^{\Sigma^{0}}$ and $L_{+}^{0}=L_{+}(\Sigma^{0},Q^{0})$. Before proving Proposition~\ref{perturbative_L+} let us check that $\sup_{0<\varepsilon\leq 1}(M_{0}^{\varepsilon})<+\infty$. We remind the reader that the sequence of ground states $(Q^{\varepsilon})_{\varepsilon>0}$ is well defined only if this supremum is finite.

\begin{lemma}\label{Lemma_mass_threshold}
Let \eqref{h4} be fulfilled. For every $M>0$ there exists $\varepsilon_{0}>0$ such that for every $\varepsilon\in(0,\varepsilon_{0})$, $M_{0}^{\varepsilon}<M$.
\end{lemma}
\begin{proof}
We start by showing that for every $u\in H_{x}^{1}$,
\[
H^{\varepsilon}(u)\underset{\varepsilon\to 0}{\to}\,H^{0}(u).
\]
Indeed, thanks to the Hölder inequality we have
\[
\left|H^{\varepsilon}(u)-H^{0}(u)\right|=\left|\int|u|^{2}\star(\Sigma^{\varepsilon}-\Sigma^{0})(x)\,|u|^{2}(x)\ud x\right|\leq\||u|^{2}\star(\Sigma^{\varepsilon}-\Sigma^{0})\|_{L_{x}^{\infty}}\|u\|_{L_{x}^{2}}^{2},
\]
and thanks to the homogeneous Sobolev embedding in dimension $d=3$ we get
\begin{align*}
\||u|^{2}\star{}&(\Sigma^{\varepsilon}-\Sigma^{0})\|_{L_{x}^{\infty}}\\
&\leq \|(\Sigma^{\varepsilon}-\Sigma^{0})\mathbf{1}_{|x|\leq R}\|_{L_{x}^{3/2}}\|\,|u|^{2}\|_{L_{x}^{3}}+\|(\Sigma^{\varepsilon}-\Sigma^{0})\mathbf{1}_{|x|>R}\|_{L_{x}^{\infty}}\|\,|u|^{2}\|_{L_{x}^{1}}\\
&\leq C\|(\Sigma^{\varepsilon}-\Sigma^{0})\mathbf{1}_{|x|\leq R}\|_{L_{x}^{3/2}}\|\nabla_{x}u\|_{L_{x}^{2}}^{2}+\|(\Sigma^{\varepsilon}-\Sigma^{0})\mathbf{1}_{|x|>R}\|_{L_{x}^{\infty}}\|u\|_{L_{x}^{2}}^{2}.
\end{align*}
Thus, assumption \eqref{h4} leads to the required convergence. We conclude as follows. By using the results of E.\,Lieb in \cite{Lieb_Choquard} we know that $K_{M}^{0}<0$ is achieved at a unique positive and radially symmetric function $Q^{0}$. Then $H^{\varepsilon}(Q^{0})\to H^{0}(Q^{0})=K_{M}^{0}<0$ implies $K_{M}^{\varepsilon}<0$ as soon as $\varepsilon$ is sufficiently small. Eventually Lemma~\ref{basic_lemma}-(d) and (e) allows us to conclude.
\end{proof}


We turn to the proof of Proposition~\ref{perturbative_L+}.

\begin{ProofOf}{\eqref{perturbative_L+i} Convergence. Step 1}
We prove that for every $u\in H_{x}^{1}$ and for every $\delta, R>0$, there exists $\varepsilon_{0}>0$ such that for every $0<\varepsilon<\varepsilon_{0}$,
\begin{equation}\label{estimation_H_epsilon}
H^{\varepsilon}(u)\geq\frac{1}{2}\|\nabla_{x}u\|_{L_{x}^{2}}^{2}-\frac{\kappa C}{2}\left(\delta+cR\right)\|u\|_{L_{x}^{2}}^{2}\|\nabla_{x}u\|_{L_{x}^{2}}^{2}-\frac{\kappa}{2}\Bigl(\delta+\frac{1}{R}\Bigr)\|u\|_{L_{x}^{2}}^{4},
\end{equation}
where $C$ denotes the best constant in the homogeneous Sobolev embedding in dimension $d=3$ and $c>0$ is a constant. Since
\begin{multline*}
H^{\varepsilon}(u)=\frac{1}{2}\|\nabla_{x}u\|_{L_{x}^{2}}^{2}-\frac{\kappa}{2}\iint|u|^{2}(x)\Sigma^{\varepsilon}(x-y)|u|^{2}(y)\ud x\ud y\\
\geq\frac{1}{2}\|\nabla_{x}u\|_{L_{x}^{2}}^{2}-\frac{\kappa}{2}\left|\iint|u|^{2}(x)\Sigma^{\varepsilon}(x-y)|u|^{2}(y)\ud x\ud y\right|
\end{multline*}
we only have to estimate the last term of the right hand side. Again, we use the Hölder inequality and the homogeneous Sobolev embedding and we obtain
\begin{align*}
\biggl|\iint|u|^{2}&(x)\Sigma^{\varepsilon}(x-y)|u|^{2}(y)\ud x\ud y\biggr|
\\
&\leq C\|\Sigma^{\varepsilon}\mathbf{1}_{|x|\leq R}\|_{L_{x}^{3/2}}\|u\|_{L_{x}^{2}}^{2}\|\nabla_{x}u\|_{L_{x}^{2}}^{2}+\|\Sigma^{\varepsilon}\mathbf{1}_{|x|>R}\|_{L_{x}^{\infty}}\|u\|_{L_{x}^{2}}^{4}\\
&\leq C\bigl(\|(\Sigma^{\varepsilon}-\Sigma^{0})\mathbf{1}_{|x|\leq R}\|_{L_{x}^{3/2}}+\|\Sigma^{0}\mathbf{1}_{|x|\leq R}\|_{L_{x}^{3/2}}\bigr)\|u\|_{L_{x}^{2}}^{2}\|\nabla_{x}u\|_{L_{x}^{2}}^{2}\\
&\hspace*{2cm}+\left(\|(\Sigma^{\varepsilon}-\Sigma^{0})\mathbf{1}_{|x|>R}\|_{L_{x}^{\infty}}+\|\Sigma^{0}\mathbf{1}_{|x|>R}\|_{L_{x}^{\infty}}\right)\|u\|_{L_{x}^{2}}^{4}.
\end{align*}
The quantities $\|\Sigma^{0}\mathbf{1}_{|x|\leq R}\|_{L_{x}^{3/2}}$ and $\|\Sigma^{0}\mathbf{1}_{|x|>R}\|_{L_{x}^{\infty}}$ can be evaluated
explicitly. Combined with the convergence \eqref{cv_potential}, it allows us to obtain \eqref{estimation_H_epsilon} for every $\delta>0$ provided $\varepsilon>0$ is sufficiently small.

\subsubsection*{Step 2} Estimate \eqref{estimation_H_epsilon} has two consequences: firstly, the sequence $(Q^{\varepsilon})_{\varepsilon>0}$ is bounded in $H_{x}^{1}$ and, secondly, the sequence $(K_{M}^{\varepsilon})_{\varepsilon>0}$ is bounded from below (at~least for $\varepsilon>0$ sufficiently small) by $-\kappa(\delta+1/R)M^{2}/2$. Indeed, we already know that $\|Q^{\varepsilon}\|_{L_{x}^{2}}^{2}=M$ and for $\delta+cR>0$ sufficiently small (that means $\varepsilon>0$ is also sufficiently small), we have $\kappa C(\delta+cR)M/2\leq 1/4$. Hence, \eqref{estimation_H_epsilon} with $u=Q^{\varepsilon}$ becomes
\[
H^{\varepsilon}(Q^{\varepsilon})\geq\frac{1}{4}\|\nabla_{x}Q^{\varepsilon}\|_{L_{x}^{2}}^{2}-\frac{\kappa}{2}\Bigl(\delta+\frac{1}{R}\Bigr)M^{2}.
\]
Since $H^{\varepsilon}(Q^{\varepsilon})=K_{M}^{\varepsilon}<0$ is negative for every $\varepsilon>0$ we eventually deduce that $\|\nabla_{x}Q^{\varepsilon}\|_{L_{x}^{2}}$ is bounded. Moreover, it is clear that the sequence $(K_{M}^{\varepsilon})_{\varepsilon>0}$ is bounded from below by $-\kappa(\delta+1/R)M^{2}/2$, as soon as $\varepsilon>0$ is sufficiently small.

Therefore, we know that $(Q^{\varepsilon})_{\varepsilon>0}$ is bounded in $H_{x}^{1}$, and we also know the existence of two constant $a,A>0$ such that for every $\varepsilon>0$ sufficiently small, \hbox{$-A\leq J_{M}^{\varepsilon}\leq -a$} (the existence of $a$ comes from the proof of Lemma~\ref{Lemma_mass_threshold} where we proved that \hbox{$K_{M}^{\varepsilon}\leq H^{\varepsilon}(Q^{0})\to H^{0}(Q^{0})=K_{M}^{0}<0$}). Moreover, since $Q^{\varepsilon}$ is a solution of \eqref{choquard} with $\Sigma=\Sigma^{\varepsilon}$ and $\omega=\omega^{\varepsilon}$, by multiplying this equation by $Q^{\varepsilon}$ and integrating over $\R^{3}$ we~get
\[
\omega^{\varepsilon}M=-\frac{1}{2}\|\nabla_{x}Q^{\varepsilon}\|_{L_{x}^{2}}^{2}+\kappa\iint|Q^{\varepsilon}|^{2}(x)\Sigma^{\varepsilon}(x-y)|Q^{\varepsilon}|^{2}(y)\ud x\ud y.
\]
In turn, the sequence $(\omega^{\varepsilon})_{\varepsilon>0}$ is bounded:
\begin{align*}
0<\frac{a}{M}\leq\omega^{\varepsilon}&=-\frac{K_{M}^{\varepsilon}}{M}+\frac{\kappa}{2M}\iint|Q^{\varepsilon}|^{2}(x)\Sigma^{\varepsilon}(x-y)|Q^{\varepsilon}|^{2}(y)\ud x\ud y\\
&\leq\frac{A}{M}+\frac{\kappa C}{2M}\left(\delta+cR\right)\|Q^{\varepsilon}\|_{L_{x}^{2}}^{2}\|\nabla_{x}Q^{\varepsilon}\|_{L_{x}^{2}}^{2}+\frac{\kappa}{2M}\Bigl(\delta+\frac{1}{R} \Bigr)\|Q^{\varepsilon}\|_{L_{x}^{2}}^{4}.
\end{align*}
There exists $\widetilde{Q}\in H_{x}^{1}$ and $\widetilde{\omega}>0$ such that, up to a subsequence, $(Q^{\varepsilon})_{\varepsilon>0}$ converges weakly to $\widetilde{Q}$ in $H_{x}^{1}$ and $(\omega^{\varepsilon})_{\varepsilon>0}$ converges to $\widetilde{\omega}$. Since the functions $Q^{\varepsilon}$ are positive and radially symmetric,
we also know that $\widetilde{Q}$ is positive and radially symmetric, and $(Q^{\varepsilon})_{\varepsilon>0}$ converges strongly to $\widetilde{Q}$ in $L_{x}^{p}$ for $2<p<6$,
see \cite{Pllsc,Strauss} for such compactness statements based on symmetry properties.

\subsubsection*{Step 3} We are going to prove that $\widetilde{Q}=Q^{0}$ and $\widetilde{\omega}=\omega^{0}$. To this end, it is sufficient to prove that $\widetilde{Q}$ is a solution of the Choquard equation \eqref{choquard} with $\Sigma=\Sigma^{0}$, $\omega=\nobreak\widetilde{\omega}$ and $\|\widetilde{Q}\|_{L_{x}^{2}}^{2}=M$. Indeed, we know that the Choquard equation with $\Sigma=\Sigma^{0}$ admits a unique positive, radially symmetric solution for $\omega=1$ (see for instance \cite{Lieb_Choquard} or \cite[App.\,A]{Lenzmann}).
This result can extended by a scaling argument for every $\omega>0$.
Hence, we~can justify the following assertion: if two positive and radially symmetric solutions~$Q_{1}$ and~$Q_{2}$ of \eqref{choquard} with $\Sigma=\Sigma^{0}$, $\omega=\omega_{1}$ and $\omega=\omega_{2}$ have the same mass, then $Q_{1}=Q_{2}$ and $\lambda_{1}=\lambda_{2}$.

For every $\varepsilon>0$ and for every $\varphi\in C_{c}^{\infty}(\R_{x}^{3})$, we have
\[
\frac{1}{2}\int\nabla_{x}Q^{\varepsilon}\cdot\nabla_{x}\varphi\ud x+\omega^{\varepsilon}\int Q^{\varepsilon}\varphi\ud x-\kappa\iint Q^{\varepsilon}\varphi(x)\Sigma^{\varepsilon}(x-y)|Q^{\varepsilon}|^{2}(y)\ud x\ud y=0.
\]
It is obvious that the first two terms converge respectively to $(\int\nabla_{x}\widetilde{Q}\cdot\nabla_{x}\varphi\ud x)/2$ and $\widetilde{\omega}\int\widetilde{Q}\varphi\ud x$ (note that for the second term we use the fact that $\|Q^{\varepsilon}\|_{L_{x}^{2}}$ is bounded with respect to $\varepsilon$). Let us now show that the third term converges to
\[
-\kappa\iint\widetilde{Q}\varphi(x)\Sigma^{0}(x-y)|\widetilde{Q}|^{2}(y)\ud x\ud y.
\]
For that purpose we decompose the difference as follows
\[
\begin{array}{l}
\ds \left|\iint Q^{\varepsilon}\varphi(x)\Sigma^{\varepsilon}(x-y)|Q^{\varepsilon}|^{2}(y)\ud x\ud y-\iint\widetilde{Q}\varphi(x)\Sigma^{0}(x-y)|\widetilde{Q}|^{2}(y)\ud x\ud y\right|\\
\ds\qquad\qquad \leq\underbrace{\left|\iint Q^{\varepsilon}\varphi(x)\left(\Sigma^{\varepsilon}(x-y)-\Sigma^{0}(x-y)\right)|Q^{\varepsilon}|^{2}(y)\ud x\ud y\right|}_{=\textup{I}}\\
\ds\qquad\qquad\qquad \qquad+\underbrace{\left|\iint \left(Q^{\varepsilon}(x)-\widetilde{Q}(x)\right)\varphi(x)\Sigma^{0}(x-y)|Q^{\varepsilon}|^{2}(y)\ud x\ud y\right|}_{=\textup{II}}\\
\ds\qquad\qquad\qquad \qquad\qquad\qquad \qquad +\underbrace{\biggl|\iint \widetilde{Q}\varphi(x)\Sigma^{0}(x-y)\left(|Q^{\varepsilon}|^{2}-|Q^{0}|^{2}\right)(y)\ud x\ud y\biggr|}_{=\textup{III}}.
\end{array}
\]
The convergence of $I$ follows from the boundedness of $(Q^{\varepsilon})_{\varepsilon>0}$ in $H_{x}^{1}$ together with the convergence \eqref{cv_potential}:
\begin{multline*}
\textup{I}\leq\|Q^{\varepsilon}\varphi\|_{L_{x}^{1}}\|(\Sigma^{\varepsilon}-\Sigma^{0})\star|Q^{\varepsilon}|^{2}\|_{L_{x}^{\infty}}\\
\leq\|Q^{\varepsilon}\|_{L_{x}^{2}}\|\varphi\|_{L_{x}^{2}}\Bigl(C\|(\Sigma^{\varepsilon}-\Sigma^{0})\mathbf{1}_{|x|\leq R}\|_{L_{x}^{3/2}}\|\nabla_{x}Q^{\varepsilon}\|_{L_{x}^{2}}^{2}
\\+\|(\Sigma^{\varepsilon}-\Sigma^{0})\mathbf{1}_{|x|>R}\|_{L_{x}^{\infty}}\|Q^{\varepsilon}\|_{L_{x}^{2}}^{2}\Bigr).
\end{multline*}
The boundedness of $(Q^{\varepsilon})_{\varepsilon>0}$ in $L_{x}^{2}$ and the strong convergence of $Q^{\varepsilon}$ to $\widetilde{Q}$ in $L_{x}^{p}$ for $2<p<6$ with $p=4$ and $p=8/3$ imply the convergence of $II$ (we use that $\Sigma^{0}\mathbf{1}_{|x|\leq R}$ lies in $L_{x}^{q}$ for $1\leq q<3$ and $\Sigma^{0}\mathbf{1}_{|x|>R}$ lies in $L_{x}^{q}$ for $q>3$):
\bgroup
\multlinegap0pt
\begin{multline*}
\textup{II}\leq\|\Sigma^{0}\star(Q^{\varepsilon}-\widetilde{Q})\varphi\|_{L_{x}^{\infty}}\|Q^{\varepsilon}\|_{L_{x}^{2}}^{2}\\
\shoveleft{\hphantom{\textup{II}}\leq\left(\|\Sigma^{0}\mathbf{1}_{|x|\leq R}\|_{L_{x}^{2}}\|(Q^{\varepsilon}-\widetilde{Q})\varphi\|_{L_{x}^{2}}+\|\Sigma^{0}\mathbf{1}_{|x|>R}\|_{L_{x}^{4}}\|(Q^{\varepsilon}-\widetilde{Q})\varphi\|_{L_{x}^{4/3}}\right)\|Q^{\varepsilon}\|_{L_{x}^{2}}^{2}}\\
\leq\left(\|\Sigma^{0}\mathbf{1}_{|x|\leq R}\|_{L_{x}^{2}}\|Q^{\varepsilon}-\widetilde{Q}\|_{L_{x}^{4}}\|\varphi\|_{L_{x}^{4}}+\|\Sigma^{0}\mathbf{1}_{|x|>R}\|_{L_{x}^{4}}\|Q^{\varepsilon}-\widetilde{Q}\|_{L_{x}^{8/3}}\|\varphi\|_{L_{x}^{8/3}}\right)\|Q^{\varepsilon}\|_{L_{x}^{2}}^{2}.
\end{multline*}
\egroup
For the last term we use almost the same strategy than for $II$. We write
\begin{multline*}
\textup{III}\leq\|\widetilde{Q}\varphi\|_{L_{x}^{1}}\|\Sigma^{0}\star(|Q^{\varepsilon}|^{2}-|\widetilde{Q}|^{2})\|_{L_{x}^{\infty}}\\
\leq\|\widetilde{Q}\|_{L_{x}^{2}}\|\varphi\|_{L_{x}^{2}}\Bigl(\left\|\Sigma^{0}\mathbf{1}_{|x|\leq R}\right\|_{L_{x}^{2}}\bigl\|\,|Q^{\varepsilon}|^{2}-|\widetilde{Q}|^{2}\bigr\|_{L_{x}^{2}}\\
+\left\|\Sigma^{0}\mathbf{1}_{|x|>R}\right\|_{L_{x}^{4}}\bigl\|\,|Q^{\varepsilon}|^{2}-|\widetilde{Q}|^{2}\bigr\|_{L_{x}^{4/3}}\Bigr).
\end{multline*}
Since $|Q^{\varepsilon}|^{2}-|\widetilde{Q}|^{2}=|Q^{\varepsilon}-\widetilde{Q}|^{2}+2(Q^{\varepsilon}-\widetilde{Q})\widetilde{Q}$ we eventually obtain
\begin{align*}
\bigl\|\,|Q^{\varepsilon}|^{2}-|\widetilde{Q}|^{2}\bigr\|_{L_{x}^{2}}&\leq\bigl\|\,|Q^{\varepsilon}-\widetilde{Q}|^{2}\bigr\|_{L_{x}^{2}}+2\bigl\|(Q^{\varepsilon}-\widetilde{Q})\widetilde{Q}\bigr\|_{L_{x}^{2}}\\
&\leq\bigl\|Q^{\varepsilon}-\widetilde{Q}\bigr\|_{L_{x}^{4}}^{2}+2\bigl\|Q^{\varepsilon}-\widetilde{Q}\bigr\|_{L_{x}^{4}}\bigl\|\widetilde{Q}\bigr\|_{L_{x}^{4}}
\end{align*}
and
\begin{align*}
\bigl\|\,|Q^{\varepsilon}|^{2}-|\widetilde{Q}|^{2}\bigr\|_{L_{x}^{4/3}}&\leq\bigl\|\,|Q^{\varepsilon}-\widetilde{Q}|^{2}\bigr\|_{L_{x}^{4/3}}+2\bigl\|(Q^{\varepsilon}-\widetilde{Q})\widetilde{Q}\bigr\|_{L_{x}^{4/3}}\\
&\leq\bigl\|Q^{\varepsilon}-\widetilde{Q}\bigr\|_{L_{x}^{8/3}}^{2}+2\bigl\|Q^{\varepsilon}-\widetilde{Q}\bigr\|_{L_{x}^{8/3}}\bigl\|\widetilde{Q}\bigr\|_{L_{x}^{8/3}}.
\end{align*}
These convergences allow us to obtain that $\widetilde{Q}$ is a solution of \eqref{choquard} with $\Sigma=\Sigma^{0}$ and $\omega=\widetilde{\omega}$. It only remains to prove that $\|\widetilde{Q}\|_{L_{x}^{2}}^{2}=M$: the weak-$L_{x}^{2}$ convergence of $Q^{\varepsilon}$ already implies $\|\widetilde{Q}\|_{L_{x}^{2}}^{2}\leq M$.

We multiply by $Q^{\varepsilon}$ the Choquard equation satisfied
by $Q^{\varepsilon}$ and we integrate over~$\R_{x}^{3}$; it yields
\[
-\omega^{\varepsilon}M=\frac{1}{2}\|\nabla_{x}Q^{\varepsilon}\|_{L_{x}^{2}}^{2}-\kappa\iint|Q^{\varepsilon}|^{2}(x)\Sigma^{\varepsilon}(x-y)|Q^{\varepsilon}|^{2}(y)\ud x\ud y.
\]
Taking $\liminf_{\varepsilon\to 0}$ leads to
\[
-\widetilde{\omega}M\geq\frac{1}{2}\|\nabla_{x}\widetilde{Q}\|_{L_{x}^{2}}^{2}-\kappa\,\underset{\varepsilon\to 0}{\limsup}\,\iint|Q^{\varepsilon}|^{2}(x)\Sigma^{\varepsilon}(x-y)|Q^{\varepsilon}|^{2}(y)\ud x\ud y.
\]
We justify as before that the last term converges to $\iint|\widetilde{Q}|^{2}(x)\Sigma^{0}(x-y)|\widetilde{Q}|^{2}(y)\ud x\ud y$. Since $\widetilde{Q}$ is a solution of \eqref{choquard} with $\Sigma=\Sigma^{0}$ and $\omega=\widetilde{\omega}$ we obtain
\[
-\widetilde{\omega}M\geq\frac{1}{2}\|\nabla_{x}\widetilde{Q}\|_{L_{x}^{2}}^{2}-\kappa\iint|\widetilde{Q}|^{2}(x)\Sigma^{0}(x-y)|\widetilde{Q}|^{2}(y)\ud x\ud y=-\widetilde{\omega}\|\widetilde{Q}\|_{L_{x}^{2}}^{2}.
\]
Since $\widetilde{\omega}>0$, we eventually obtain $M\leq\|\bar{Q}\|_{L_{x}^{2}}^{2}$ and thus $\widetilde{Q}=Q^{0}$ and $\widetilde{\omega}=\omega^{0}$.

\subsubsection*{Step 4} In order to conclude the proof it only remains to justify that the weak convergence of (a sub-sequence of) $(Q^{\varepsilon})_{\varepsilon>0}$ to $Q^{0}$ in $H_{x}^{1}$ actually holds strongly (then, thanks to the uniqueness of $Q^{0}$, one can extend this convergence to the entire sequence). We already know that $\|Q^{0}\|_{L_{x}^{2}}^{2}=M=\|Q^{\varepsilon}\|_{L_{x}^{2}}^{2}$, which implies the strong convergence of $(Q^{\varepsilon})_{\varepsilon>0}$ in $L_{x}^{2}$. We turn to the strong convergence of $(\nabla_{x}Q^{\varepsilon})_{\varepsilon>0}$ in~$L_{x}^{2}$. The end of the previous step tells us that
\[
\underset{\varepsilon\to 0}{\lim}\,\|\nabla_{x}Q^{\varepsilon}\|_{L_{x}^{2}}^{2}=
2\left(
-\omega^{0}M+\kappa\iint|Q^{0}|^{2}(x)\Sigma^{0}(x-y)|Q^{0}|^{2}(y)\ud x\ud y \right)=\|\nabla_{x}Q^{0}\|_{L_{x}^{2}}^{2},
\]
which finishes the proof.
\end{ProofOf}

\begin{ProofOf}{\eqref{perturbative_L+ii} Coercivity}
We fix $\varepsilon>0$ and we consider a positive and radially symmetric minimizer $Q^{\varepsilon}$ of $K_{M}^{\varepsilon}$. Proposition~\ref{Lemma_L+} gives
\[
\left\langle L_{+}^{0}f,f\right\rangle_{L_{x}^{2}}\geq\nu^{0}\|f\|_{H_{x}^{1}}^{2}-\frac{1}{\nu^{0}}\Bigl(\left|\langle f,Q^{0}\rangle_{L_{x}^{2}}\right|^{2}+\sum\limits_{j=1}^{d}\left|\langle f,\partial_{x_{j}}Q^{0}\rangle_{L_{x}^{2}}\right|^{2}\Bigr).
\]
Next, we compute $\langle L_{+}^{\varepsilon}f,f\rangle$ as follows:
\begin{align*}
\bigl\langle L_{+}^{\varepsilon}f&,f\bigr\rangle_{L_{x}^{2}}=\left\langle L_{+}^{0}f,f\right\rangle_{L_{x}^{2}}+\left\langle (L_{+}^{\varepsilon}-L_{+}^{0})f,f\right\rangle_{L_{x}^{2}}\\
&\geq\nu^{0}\|f\|_{H_{x}^{1}}^{2}-\frac{1}{\nu^{0}}\Bigl(\left|\langle f,Q^{0}\rangle_{L_{x}^{2}}\right|^{2}+\sum\limits_{j=1}^{d}\left|\langle f,\partial_{x_{j}}Q^{0}\rangle_{L_{x}^{2}}\right|^{2}\Bigr)-\bigl|\left\langle(L_{+}^{\varepsilon}-L_{+}^{0})f,f\right\rangle_{L_{x}^{2}}\bigr|
\\
&\geq\nu^{0}\|f\|_{H_{x}^{1}}^{2}-\frac{1}{\nu^{0}}\Bigl(\left|\langle f,Q^{\varepsilon}\rangle_{L_{x}^{2}}\right|^{2}+\sum\limits_{j=1}^{d}\left|\langle f,\partial_{x_{j}}Q^{\varepsilon}\rangle_{L_{x}^{2}}\right|^{2}\Bigr)
\\
&\hspace*{5cm}-\frac{1}{\nu^{0}}R^{\varepsilon}-\bigl|\left\langle(L_{+}^{\varepsilon}-L_{+}^{0})f,f\right\rangle_{L_{x}^{2}}\bigr|,
\end{align*}
where\vspace*{-5pt}
\begin{multline*}
R^{\varepsilon}=\left|\langle f,Q^{0}-Q^{\varepsilon}\rangle_{L_{x}^{2}}\right|^{2}+\sum\limits_{j=1}^{d}\left|\langle f,\partial_{x_{j}}Q^{0}-\partial_{x_{j}}Q^{\varepsilon}\rangle_{L_{x}^{2}}\right|^{2}\\[
-8pt]
+2\left|\langle f,Q^{0}-Q^{\varepsilon}\rangle_{L_{x}^{2}}\right|\left|\langle f,Q^{\varepsilon}\rangle_{L_{x}^{2}}\right|+2\sum\limits_{j=1}^{d}\left|\langle f,\partial_{x_{j}}Q^{0}-\partial_{x_{j}}Q^{\varepsilon}\rangle_{L_{x}^{2}}\right|\bigl|\langle f,\partial_{x_{j}}Q^{\varepsilon}\rangle_{L_{x}^{2}}\bigr|.
\end{multline*}
Then we infer the following estimate: $R^{\varepsilon}\!\leq\!\alpha(Q^{\varepsilon})\|f\|_{H_{x}^{1}}^{2}$ where $\alpha(Q)\!>\!0$ and \hbox{$\alpha(Q)\!\to\! 0$} when $\|Q-Q^{0}\|_{H_{x}^{1}}\to 0$. Moreover\vspace*{-5pt}
\begin{multline*}
\left\langle(L_{+}^{\varepsilon}-L_{+}^{0})f,f\right\rangle_{L_{x}^{2}}=\left(\omega^{\varepsilon}-\omega^{0}\right)\|f\|_{L_{x}^{2}}^{2}-\kappa\int\left(\Sigma^{\varepsilon}\star|Q^{\varepsilon}|^{2}-\Sigma^{0}\star|Q^{0}|^{2}\right)|f|^{2}\ud x\\
-2\kappa\iint\left(Q^{\varepsilon}f(x)\Sigma^{\varepsilon}(x-y)Q^{\varepsilon}f(y)-Q^{0}f(x)\Sigma^{0}(x-y)Q^{0}f(y)\right)\ud x\ud y,
\end{multline*}
and from this expression we can obtain (thanks to a reasoning similar to the proof of point (i)) the following estimate
\[
\bigl|\left\langle(L_{+}^{\varepsilon}-L_{+}^{0})f,f\right\rangle_{L_{x}^{2}}\bigr|\leq\beta(\Sigma^{\varepsilon},Q^{\varepsilon},\omega^{\varepsilon})\|f\|_{H_{x}^{1}}^{2},
\]
where $\beta(\Sigma,Q,\omega)>0$ and $\beta(\Sigma,Q,\omega)\to 0$ when
\[
\|(\Sigma-\Sigma^{0})\mathbf{1}_{|x|\leq R}\|_{L_{x}^{3/2}}+\|(\Sigma-\Sigma^{0})\mathbf{1}_{|x|> R}\|_{L_{x}^{\infty}}+\|Q-Q^{0}\|_{H_{x}^{1}}+|\omega-\omega^{0}|\to 0.
\]
This assertion applies for any $R>0$;
here $R$ is fixed once for all (not necessarily small as in the proof of convergence).
Gathering these two estimates leads to\vspace*{-5pt}
\begin{multline*}
\left\langle L_{+}^{\varepsilon}f,f\right\rangle_{L_{x}^{2}}\geq
\Bigl(\nu^{0}-\frac{\alpha(Q^{\varepsilon})}{\nu^{0}}-\beta(\Sigma^{\varepsilon},Q^{\varepsilon},\omega^{\varepsilon})\Bigr)\|f\|_{H_{x}^{1}}^{2}\\
-\frac{1}{\nu^{0}}\Bigl(\left|\langle f,Q^{\varepsilon}\rangle_{L_{x}^{2}}\right|^{2}+\sum\limits_{j=1}^{d}\left|\langle f,\partial_{x_{j}}Q^{\varepsilon}\rangle_{L_{x}^{2}}\right|^{2}\Bigr).
\end{multline*}
The announced coercivity property holds for the ground state $Q^{\varepsilon}$ provided $\alpha(Q^{\varepsilon})/\nu^{0}+\beta(\Sigma^{\varepsilon},Q^{\varepsilon},\omega^{\varepsilon})<\nu^{0}$. Since $\alpha(Q)$ and $\beta(\Sigma,Q,\omega)$ converge to zero when we have
\[
\|(\Sigma-\Sigma^{0})\mathbf{1}_{|x|\leq R}\|_{L_{x}^{3/2}}+\|(\Sigma-\Sigma^{0})\mathbf{1}_{|x|> R}\|_{L_{x}^{\infty}}+\|Q-Q^{0}\|_{H_{x}^{1}}+|\omega-\omega^{0}|\to 0,
\]
there exists $\delta>0$ such that
\[
\|(\Sigma-\Sigma^{0})\mathbf{1}_{|x|\leq R}\|_{L_{x}^{3/2}}+\|(\Sigma-\Sigma^{0})\mathbf{1}_{|x|> R}\|_{L_{x}^{\infty}}+\|Q-Q^{0}\|_{H_{x}^{1}}+|\omega-\omega^{0}|<\delta
\]
implies $\alpha(Q)/\nu^{0}+\beta(\Sigma,Q,\omega)<\nu^{0}$. Thanks to \eqref{h4} we can find $\bar{\varepsilon}_{0}>0$ such that for every $\varepsilon\in(0,\bar{\varepsilon}_{0})$,
\[
\|(\Sigma-\Sigma^{0})\mathbf{1}_{|x|\leq R}\|_{L_{x}^{3/2}}+\|(\Sigma-\Sigma^{0})\mathbf{1}_{|x|> R}\|_{L_{x}^{\infty}}<\frac{\delta}{2}.
\]
Therefore, possibly by choosing a smaller $\bar{\varepsilon}_{0}$ if necessary, for every $\varepsilon\in(0,\bar{\varepsilon}_{0})$ and every positive and radially symmetric minimizer $Q^{\varepsilon}$ of $K_{M}^{\varepsilon}$, we get
\[
\|Q^{\varepsilon}-Q^{0}\|_{H_{x}^{1}}+|\omega^{\varepsilon}-\omega^{0}|<\frac{\delta}{2}
\]
by using (i).
\end{ProofOf}

\section{Admissible form functions: proof of Pro\-po\-si\-tion \ref{non_perturbative_L+}}\label{large_class_sigma1}

The general strategy relies on the application of Proposition~\ref{perturbative_L+}; hence we have to
construct
a sequence of potentials $(\Sigma^{\varepsilon})_{\varepsilon>0}$, with the specific form $\Sigma^{\varepsilon}=\sigma_{1}^{\varepsilon}\star\sigma_{1}^{\varepsilon}$, which converges to $\Sigma^{0}$ in the sense of \eqref{cv_potential}.
This requires some care beyond the classical ``regularization and truncation'' approach.
A similar difficulty arises, but in a different manner, when justifying the asymptotic regime of the Vlasov-wave system~\eqref{kin}, \eqref{wave_kin2} towards the Vlasov-Poisson equation \cite{dBGV}.
The following simple examples are quite illuminating on the strategy.

\subsubsection*{Toy example 1} Let
$\chi:\mathbb R^3\to [0,1]$ be a $C^\infty_c$ function which satisfies
$\chi(x)= 1$ for $|x|\leq 1$ and $\chi(x)=0$ for $|x|\geq 2$.
Let
\[
\Sigma^{\varepsilon}(x)=\ds\frac{\chi(\varepsilon x)}{|x|}.
\]
The analysis of this kernel is simple: due to the scale invariance of $\sfrac{1}{|x|}$, we have
\[
\Sigma^{\varepsilon}(x)=\varepsilon\,\frac{\chi(\varepsilon x)}{|\varepsilon x|}=\varepsilon\Sigma^{1}(\varepsilon x).
\]
As a matter of fact, we have
\begin{enumeratei}
\item\label{enumc:i}
$H^{\Sigma^{\varepsilon}}(u)=\varepsilon^{3}H^{\Sigma^{1}}(u^{\varepsilon})$ where $u^{\varepsilon}(x)=\varepsilon^{-2}u(\varepsilon^{-1}x)$,
\item\label{enumc:ii}
$Q^{\varepsilon}$ is a minimizer of $K_{M}^{\Sigma^{\varepsilon}}\Leftrightarrow Q(x)=\varepsilon^{-2}Q^{\varepsilon}(\varepsilon^{-1}x)$ is a minimizer of $K_{\varepsilon^{-1}M}^{\Sigma^{1}}$,
\item\label{enumc:iii}
$K_{M}^{\Sigma^{\varepsilon}}=\varepsilon^{3}K_{\varepsilon^{-1}M}^{\Sigma^{1}}$,
\item\label{enumc:iv}
if $Q^{\varepsilon}$ is a minimizer of $K_{M}^{\Sigma^{\varepsilon}}$, then $\omega(\Sigma^{\varepsilon},Q^{\varepsilon})=\varepsilon^{2}\omega(\Sigma^{1},Q)$, where $Q(x)=\varepsilon^{-2}Q^{\varepsilon}(\varepsilon^{-1}x)$,
\item\label{enumc:v}
$\left\langle L_{+}(\Sigma^{\varepsilon},Q^{\varepsilon})f^{\varepsilon},f^{\varepsilon}\right\rangle_{L_{x}^{2}}=\varepsilon^{3}\left\langle L_{+}(\Sigma^{1},Q)f,f\right\rangle_{L_{x}^{2}}$, where $f(x)=\varepsilon^{-2}f^{\varepsilon}(\varepsilon^{-1}x)$ and still $Q(x)=\varepsilon^{-2}Q^{\varepsilon}(\varepsilon^{-1}x)$.
\end{enumeratei}
These relations provide several useful information. For example, since for any fixed $\varepsilon>0$, $\Sigma^{\varepsilon}$ lies in $L_{x}^{3/2}$, Lemma~\ref{basic_lemma} applies and justifies the existence of the mass threshold $M_{0}^{\Sigma^{\varepsilon}}$, which, in turn, can be expressed by means of $M_{0}^{\Sigma^{1}}$:
$M_{0}^{\Sigma^{\varepsilon}}=\varepsilon M_{0}^{\Sigma^{1}}\to\nobreak 0$. Furthermore, $\Sigma^{\varepsilon}$ converges to $\Sigma^{0}$ in the sense of \eqref{cv_potential},
and the
conclusions of Proposition~\ref{perturbative_L+} hold. Then, relation \eqref{enumc:v} allows us to extend the coercivity estimate to any radially symmetric minimizer of $K_{m}^{\Sigma^{1}}$ associated to a mass $m$ larger than $M/\bar{\varepsilon}_{0}$, as~illustrated by Figure~\ref{fig_st}. Indeed \eqref{enumc:ii}, \eqref{enumc:v} and Proposition~\ref{perturbative_L+}\eqref{perturbative_L+ii} yield
\begin{align*}
\bigl\langle L_{+}(&\Sigma^{1},Q)f,f\bigr\rangle_{L_{x}^{2}}=\varepsilon^{-3}\left\langle L_{+}(\Sigma^{\varepsilon},Q^{\varepsilon})f^{\varepsilon},f^{\varepsilon}\right\rangle_{L_{x}^{2}}\\
&\geq\varepsilon^{-3}\nu^{\varepsilon}\|f^{\varepsilon}\|_{H_{x}^{1}}^{2}-\frac{\varepsilon^{-3}}{\nu^{0}}\Bigl(\left|\langle f^{\varepsilon},Q^{\varepsilon}\rangle_{L_{x}^{2}}\right|^{2}+\sum\limits_{j=1}^{3}\left|\langle f^{\varepsilon},\partial_{x_{j}}Q^{\varepsilon}\rangle_{L_{x}^{2}}\right|^{2}\Bigr)\\
&=\nu^{\varepsilon}\|\nabla_{x}f\|_{L_{x}^{2}}^{2}+\varepsilon^{-2}\nu^{\varepsilon}\|f\|_{L_{x}^{2}}^{2}
-\frac{1}{\nu^{0}}\Bigl(\varepsilon^{-2}\left|\langle f,Q\rangle_{L_{x}^{2}}\right|^{2}+\varepsilon^{-1}\sum\limits_{j=1}^{3}\left|\langle f^{\varepsilon},\partial_{x_{j}}Q^{\varepsilon}\rangle_{L_{x}^{2}}\right|^{2}\Bigr),
\end{align*}
which implies the announced coercivity property.

This example can be compared to the case of the Yukawa potential seen as a perturbation of the Newtonian potential in \cite{jmaa}.

\begin{figure}[!ht]
\definecolor{qqqqff}{rgb}{0,0,1}
\begin{center}
\resizebox{0.5\linewidth}{!}{
\begin{tikzpicture}[line cap=round,line join=round,>=triangle 45
,x=3.8cm,y=3.8cm]
\draw[->,color=black] (-0.,0) -- (2,0);
\foreach \x in {,0.2,0.4,0.6,0.8,1,1.2,1.4,1.6,1.8}
\draw[shift={(\x,0)},color=black] (0pt,-2pt);
\draw[->,color=black] (0,-0.) -- (0,1.3);
\foreach \y in {,0.2,0.4,0.6,0.8,1,1.2,1.4}
\draw[shift={(0,\y)},color=black] (2pt,0pt);
\clip(-0.15,-0.15) rectangle (2,1.3);
\draw [dash pattern=on 3pt off 3pt] (0.4,1)-- (0.4,0.4);
\draw (0.4,0.4)-- (0.4,0);
\draw (0.8,1)-- (2.5,1);
\draw [->,dash pattern=on 6pt off 6pt] (0.4,0.4) -- (0.8,1);
\draw [->,dash pattern=on 6pt off 6pt] (0.4,0.3) -- (1.1,1);
\draw [->,dash pattern=on 6pt off 6pt] (0.4,0.2) -- (1.6,1);
\draw [->,dash pattern=on 6pt off 6pt] (0.4,0.1) -- (2.5,1);
\draw [dotted] (0.4,1)-- (0,1);
\draw [dotted] (0.8,1)-- (0.8,0);
\draw (-0.1,1.05) node[anchor=north west] {$1$};
\draw (-0.11,0.45) node[anchor=north west] {$\bar{\varepsilon}_{0}$};
\draw [dotted] (0.4,0.4)-- (0,0.4);
\draw (0.34,-0.02) node[anchor=north west] {$ M $};
\draw (-0.1,1.3) node[anchor=north west] {$ \varepsilon $};
\draw (1.77,-0.02) node[anchor=north west] {$ \text{Mass} $};
\draw (0.73,-0.02) node[anchor=north west] {\parbox{1.55 cm}{$ M/\bar{\varepsilon}_{0} $}};
\begin{scriptsize}
\fill [color=qqqqff] (0.4,1) circle (1.5pt);
\fill [color=qqqqff] (0.4,0.4) circle (1.5pt);
\fill [color=qqqqff] (0.8,1) circle (1.5pt);
\end{scriptsize}
\end{tikzpicture}}
\caption{Illustration of the strategy:
for the given mass $M$,
the stability of the ground states is proved for the potentials $\Sigma^\varepsilon$, with $0\leq \varepsilon<\bar \varepsilon_0$.
By rescaling, we can go back to the potentials $\Sigma^1$,
and ground states with a mass larger that $M/\bar \varepsilon_0$ are stable.\label{fig_st}
}
\end{center}
\end{figure}

\subsubsection*{Toy example 2}
Let $\alpha:\mathbb R^3\to [0,\infty)$ be a $C^\infty$ function such that $\int \alpha\ud x=1$. We~consider
\[
\Sigma^{\varepsilon}(x)=\varepsilon^{-3}\ds\int \ds\frac{\alpha(\varepsilon^{-1}y)}{|x-y|}\ud y.
\]
Now, we have the scaling relation: $\Sigma^{\varepsilon}(x)=\varepsilon^{-1}\Sigma^{1}(\varepsilon^{-1}x)$, where
\[
\Sigma^{1}(x)=\ds\int \ds\frac{\alpha(y)}{|x-y|}\ud y.
\]
We deduce that
\[
Q^{\varepsilon} \text{ is a minimizer of } K_{M}^{\Sigma^{\varepsilon}} \iff Q(x)=\varepsilon^{2}Q^{\varepsilon}(\varepsilon x) \text{ is a minimizer of } K_{\varepsilon M}^{\Sigma^{1}}.
\]
Reasoning as in the previous example, we obtain that, for $M$ sufficiently small, every positive and radially symmetric minimizer of $K_{M}^{\Sigma^{1}}$ satisfies the coercivity relation~\eqref{coercivity2}. In particular there is no mass threshold: $M_{0}^{\Sigma^{1}}=0$. Since $\Sigma^{1}\notin L_{x}^{3/2}$, this is not a contradiction with Lemma~\ref{basic_lemma}.

\subsubsection*{Toy example 3}
We go back to the case of the dimension $d=1$. In this case for any $\sigma_{1}$ satisfying \eqref{h2}--\eqref{h3} we consider the sequence of potential $(\Sigma^{\varepsilon})$ defined by
\[
\Sigma^{\varepsilon}(x)=\varepsilon^{-1}\Sigma(\varepsilon^{-1}x), \hspace{0.5cm} \Sigma=\sigma_{1}\star\sigma_{1}.
\]
Hence we obtain the equivalence
\[
Q^{\varepsilon} \text{ is a minimizer of } K_{M}^{\Sigma^{\varepsilon}} \iff Q(x)=\varepsilon Q^{\varepsilon}(\varepsilon x) \text{ is a minimizer of } K_{\varepsilon M}^{\Sigma^{1}}.
\]
Reasoning as above, we justify the existence of some $M^{*}>0$ such that for every $M\in(0,M^{*})$, every positive and even minimizer of $K_{M}^{\Sigma^{1}}$ satisfies the coercivity relation~\eqref{coercivity2}.

\subsubsection*{Main strategy} The toy examples 1 and 2 do not fit with our framework, where we are dealing with smooth and compactly supported potentials $\Sigma$. Then, in order to handle such a potential, the idea is (as usual) to combine the truncation and the regularization by setting (mind that $d=3$)
\[
\Sigma^{\varepsilon}(x)=\varepsilon^{-3}\chi(\varepsilon x)\ds\int\ds\frac{\alpha(\varepsilon^{-1}y)}{|x-y|}\ud y.
\]
However, the scaling for the truncation and for the regularization are not the same,
and the properties deduced from the scale
invariance of $\sfrac{1}{|x|}$ break down.
Instead, we consider a doubly indexed sequence of potentials
\[
\Sigma^{\lambda,\mu}(x)=\lambda^{-3}\chi(\mu x)\ds\int\ds\frac{\alpha(\lambda^{-1}y)}{|x-y|}\ud y
\]
with $\lambda,\mu>0$.
We also introduce
\[
\widetilde{\Sigma}^{\epsilon}(x)=\epsilon^{-3}\chi(x)\ds\int \frac{\alpha(\epsilon^{-1}y)}{|x-y|}\ud y.
\]
We have the scaling relation $\Sigma^{\lambda,\mu}(x)=\mu\widetilde{\Sigma}^{\lambda\mu}(\mu x)$ which leads to the following lemma.
\begin{lemma}\label{lemma_scaling}
The following assertions hold:
\begin{enumeratei}
\item
$H^{\Sigma^{\lambda,\mu}}(u)=\mu^{3}H^{\widetilde{\Sigma}^{\epsilon}}(u^{\mu})$, where $u^{\mu}(x)=\mu^{-2}u(\mu^{-1}x)$ and $\epsilon=\lambda\mu$,
\item
$Q^{\lambda,\mu}$ is a minimizer of $K_{M}^{\widetilde{\Sigma}^{\lambda,\mu}}\Leftrightarrow Q(x)=\mu^{-2}Q^{\lambda,\mu}(\mu^{-1}x)$ is a minimizer of $K_{\mu^{-1}M}^{\widetilde{\Sigma}^{\epsilon}}$ with $\epsilon=\lambda\mu$,
\item
$K_{M}^{\Sigma^{\lambda,\mu}}=\mu^{3}K_{\mu^{-1}M}^{\widetilde{\Sigma}^{\epsilon}}$ with $\epsilon=\lambda\mu$,
\item
if $Q^{\lambda,\mu}$ is a minimizer of $K_{M}^{\Sigma^{\lambda,\mu}}$, then $\omega(\Sigma^{\lambda,\mu},Q^{\lambda,\mu})=\mu^{2}\omega(\widetilde{\Sigma}^{\epsilon},Q)$, where $Q(x)=\mu^{-2}Q^{\lambda,\mu}(\mu^{-1}x)$ and $\epsilon=\lambda\mu$,
\item
$\left\langle L_{+}(\Sigma^{\lambda,\mu},Q^{\lambda,\mu})f^{\lambda,\mu},f^{\lambda,\mu}\right\rangle_{L_{x}^{2}}=\mu^{3}\bigl\langle L_{+}(\widetilde{\Sigma}^{\epsilon},Q)f,f\bigr\rangle_{L_{x}^{2}}$, where we have set $Q(x)=\mu^{-2}Q^{\lambda,\mu}(\mu^{-1}x)$, $f(x)=\mu^{-2}f^{\lambda,\mu}(\mu^{-1}x)$ and $\epsilon=\lambda\mu$.
\end{enumeratei}
\end{lemma}

Let us suppose for a while that the sequence $(\Sigma^{\lambda,\mu})_{\lambda, \mu >0}$ converges to $\Sigma^{0}$ in the sense of \eqref{cv_potential} as $\lambda$ and $\mu $ tend to 0. Then there exists $\lambda_{0}>0$ and $\mu_{0}>0$ such that for any $(\lambda,\mu)\in(0,\lambda_{0})\times(0,\mu_{0})$, the conclusions of Proposition~\ref{perturbative_L+} hold. Based on Lemma~\ref{lemma_scaling}, we
infer the following statement.

\skpt
\begin{proposition}\label{prop_scaling}
\begin{enumeratei}
\item\label{prop_scalingi}
For every $(\lambda,\mu)\in(0,\lambda_{0})\times(0,\mu_{0})$ and for every positive and radially symmetric minimizer Q of $K_{\mu^{-1}M}^{\widetilde{\Sigma}^{\epsilon}}$ with $\epsilon=\lambda\mu$, the operator $L_{+}(\widetilde{\Sigma}^{\epsilon},Q)$ satisfies Lemma~\ref{Lemma_L+}.

\item\label{prop_scalingii}
In particular, for $\epsilon\in(0,\lambda_{0}\mu_{0})$ fixed, applying (i) to any $(\lambda,\mu)\in(0,\lambda_{0})\times(0,\mu_{0})$ such that $\lambda\mu=\epsilon$
implies that for any $m\in(\mu_{0}^{-1}M,\lambda_{0}\epsilon^{-1}M)$ and any positive and radially symmetric minimizer $Q$ of $K_{m}^{\widetilde{\Sigma}^{\epsilon}}$, the operator $L_{+}(\widetilde{\Sigma}^{\epsilon},Q)$ satisfies Lemma~\ref{Lemma_L+}.
\end{enumeratei}
\end{proposition}

Item \eqref{prop_scalingii} implies, up to the fact that $\widetilde{\Sigma}^{\epsilon}$ can be cast under the form $\widetilde{\Sigma}^{\epsilon}=\widetilde{\sigma_{1}}^{\epsilon}\star\widetilde{\sigma_{1}}^{\epsilon}$, that the set of admissible form function $\mathscr{A}$ is non empty. Then, to conclude the proof it only remains to slightly adapt the previous construction in order to obtain a sequence $\Sigma^{\lambda,\mu}$ satisfying \eqref{h4}.
We proceed as follows.
Let $\alpha, \chi$ be two $C^\infty_c(\mathbb R^3)$, non negative, radially symmetric, compactly supported and non increasing functions, with $\chi(x)=1$ in a neighborhood of the origin.
Let us set
\[
\sigma_{1}^{\lambda,\mu}(x)=\lambda^{-3}\int_{\R^{3}}\alpha(\lambda^{-1}y)\,\frac{\chi(\mu[x-y])}{|x-y|^{2}}\ud y=\alpha^{\lambda}\star\Bigl(\frac{\chi^{\mu}}{|\cdot|^{2}}\Bigr)(x)
\quad\text{and}\quad
\Sigma^{\lambda,\mu}=\sigma_{1}^{\lambda,\mu}\star\sigma_{1}^{\lambda,\mu},
\]
where
\[
\alpha^{\lambda}(x)=\lambda^{-3}\alpha(\lambda^{-1}x) \hspace{0.25cm}\text{and}\hspace{0.25cm} \chi^{\mu}(x)=\chi(\mu x).
\]
Then each $\sigma_{1}^{\lambda,\mu}$ satisfies \eqref{h2}--\eqref{h3}. Moreover we can check that
\[
\sigma_{1}^{\lambda,\mu}(x)=\mu^{2}\widetilde{\sigma_{1}}^{\lambda\mu}(\mu x),\qquad
\Sigma^{\lambda,\mu}(x)=\mu\widetilde{\Sigma}^{\epsilon}(\mu x),
\]
where
\[
\widetilde{\sigma_{1}}^{\epsilon}(x)=\ds\int \alpha^{\epsilon}(x-y)\,\frac{\chi (y)}{| y |^{2}}\ud y,
\qquad
\widetilde{\Sigma}^{\epsilon}=\widetilde{\sigma_{1}}^{\epsilon}\star\widetilde{\sigma_{1}}^{\epsilon}.
\]
Then Lemma~\ref{lemma_scaling} applies to this new sequence as well and Proposition~\ref{prop_scaling} holds
provided we can show that it
converges to $\Sigma^{0}$ in the sense of \eqref{cv_potential}.
Such a form function appeared in \cite{dBGV}. The construction is based on the following two observations:
\[
\frac{1}{|\cdot|^{2}}\star\frac{1}{|\cdot|^{2}}(x)=\frac{C}{|x|}=C\,\Sigma^{0}(x),\quad\text{where } C=\int_{\R^{3}}\frac{\ud y}{|y|^2|e_1-y|^2}
\]
($e_1$ being the first vector of the canonical basis),
and
\[
\Sigma^{\lambda,\mu}=(\alpha^{\lambda}\star\alpha^{\lambda})\star\Bigl(\frac{\chi^{\mu}}{|\cdot|^{2}}\star\frac{\chi^{\mu}}{|\cdot|^{2}}\Bigr).
\]
Then, at least formally,
\begin{align*}
\alpha^{\lambda}\star\alpha^{\lambda}&\to\biggl(\int\alpha\star\alpha\ud x\biggr)\delta_{0}\quad\text{when }\lambda\to 0,\\
(\chi^{\mu}/|\cdot|^{2})\star(\chi^{\mu}/|\cdot|^{2})&\to(1/|\cdot|^{2})\star(1/|\cdot|^{2})=C\,\Sigma^{0}\quad\text{when }\mu\to 0,
\end{align*}
and we can expect that $\Sigma^{\lambda,\mu}$ looks like
$\Sigma^{0}$ when $\lambda,\mu\to 0$ provided $\int\alpha\ud x=1/\sqrt{C}$.
The intuition is confirmed by the following claim.

\begin{lemma}\label{lemma_cv_potential}
If $\int\alpha\ud x=1/\sqrt{C}$, then the sequence $(\Sigma^{\lambda,\mu})_{\lambda,\mu>0}$ converges to $\Sigma^{0}$ in the sense of \eqref{cv_potential} when $(\lambda,\mu)\to(0,0)$.
\end{lemma}

This approach allows us to construct a large class of admissible form functions, not necessarily close de $\Sigma^0$ in the sense of \eqref{cv_potential}, by using suitable rescalings that preserve the coercivity estimate as we did with the toy example 1.
Indeed, for any $\alpha$ and $\chi$ defined as before, if the form function $\sigma_{1}=\alpha\star(\chi/|\cdot|^{2})$ is not in $\mathscr{A}$ we know, at least up to rescaling $\alpha$ into $\alpha^{\epsilon}(x)=\epsilon^{-3}\alpha(\epsilon^{-1}x)$,
that the form functions $\widetilde{\sigma_{1}}^{\epsilon}=\alpha^{\epsilon}\star(\chi/|\cdot|^{2})$ belong to $\mathscr{A}$ provided $\epsilon$ is sufficiently small.
With the previous notation the non empty mass interval $I$ associated to the form function $\widetilde{\sigma_{1}}^{\epsilon}$ is given by $I=(\mu_{0}^{-1}M,\lambda_{0}\epsilon^{-1}M)$. It is also possible to rescale $\chi$ into $\chi^{\epsilon}(x)=\chi(\epsilon x)$ and obtain that form functions $\check{\sigma}_{1}^{\epsilon}=\alpha\star(\chi^{\epsilon}/|\cdot|^{2})$ equally belong to $\mathscr{A}$ provided $\epsilon$ is sufficiently small (this second example uses the scaling relation $\sigma_{1}^{\lambda,\mu}(x)=\lambda^{-2}\check{\sigma}_{1}^{\lambda\mu}(\lambda^{-1}x)$).
Moreover given an admissible function $\sigma_{1}$, we observe that $\sigma_{1}^{\lambda,\mu}(x)=\lambda\sigma_{1}(\mu x)$ is admissible too.
We obtain this way form functions with arbitrary support size and $L_{x}^{\infty}$-norm, which are non negative, non increasing, radially symmetric and concentrated around the origin. Such form functions are physically meaningful in the framework defined in \cite{BdB}.
Since they are simply derived by rescaling, we can check that the necessary coercivity estimate still holds, with constants that keep track of the rescaling, and they also provide stable ground states.

\begin{ProofOf}{Lemma~\ref{lemma_cv_potential}}
Let $0<R<\infty$ be fixed once for all. We decompose the difference $\Sigma^{\lambda,\mu}-\Sigma^{0}$ as follows
\begin{multline*}
\Sigma^{\lambda,\mu}(x)-\Sigma^{0}(x)=\left(\alpha^{\lambda}\star\alpha^{\lambda}\right)\star\Bigl(\frac{\chi^{\mu}}{|\cdot|^{2}}\star\frac{\chi^{\mu}}{|\cdot|^{2}}-\frac{1}{|\cdot|^{2}}\star\frac{1}{|\cdot|^{2}}\Bigr)(x)\\
+C\int\left(\alpha^{\lambda}\star\alpha^{\lambda}\right)(y)\left(\Sigma^{0}(x-y)-\Sigma^{0}(x)\right)\ud y=I_{1}(x)+I_{2}(x).
\end{multline*}
Bearing in mind that $\alpha^{\lambda}\star\alpha^{\lambda}(x)=\lambda^{-3}\alpha\star\alpha(\lambda^{-1}x)$,
we readily obtain
the convergence of $I_{2}\mathbf{1}_{|x|\leq R}$ to 0 in the $L_{x}^{3/2}$-norm. Moreover, since the support of $\alpha^{\lambda}\star\alpha^{\lambda}$ shrinks to~$\{0\}$ when $\lambda\to 0$ and since the function $x\mto 1/|x|$ is a Lipschitz function on every set of the form $\complement B(0,R)$ (with a Lipschitz constant $L(R)$ which blows up when $R\to 0$) we get
\[
\|I_{2}\mathbf{1}_{|x|>R}\|_{L_{x}^{\infty}}\lesssim\mathrm{meas}\left(\mathrm{supp}\left(\alpha^{\lambda}\star\alpha^{\lambda}\right)\right)\,\underset{\lambda\to 0}{\to}\,0.
\]
Next, for $y\in \mathrm{supp}(\alpha^\lambda\star\alpha^\lambda)$ with $\lambda$ sufficiently small, $|x|>R$ implies $|x-y|>R/2$; it follows that
\begin{align*}
\|I_{1}\mathbf{1}_{|x|>R}\|_{L_{x}^{\infty}}
&\leq\Bigl\|\Bigl(\frac{\chi^{\mu}}{|\cdot|^{2}}\star\frac{\chi^{\mu}}{|\cdot|^{2}}-\frac{1}{|\cdot|^{2}}\star\frac{1}{|\cdot|^{2}}\Bigr)\mathbf{1}_{|x|>R/2}\Bigr\|_{L_{x}^{\infty}}\\
&=\underset{|x|>R/2}{\sup}\,\left|\int\frac{\chi^{\mu}(x-y)\chi^{\mu}(y)-1}{|x-y|^{2}|y|^{2}}\ud y\right|\\
&\leq\underset{|x|>R/2}{\sup}\,\left|\int\frac{\chi^{\mu}(x-y)(\chi^{\mu}(y)-1)}{|x-y|^{2}|y|^{2}}\ud y\right|+\underset{|x|>R/2}{\sup}\,\biggl|\int\frac{\chi^{\mu}(z)-1}{|z|^{2}|x+z|^{2}}\ud z\biggr|.
\end{align*}
Since $0\leq\chi\leq 1$ and $\chi^{\mu}(x)=1$ when $|x|\leq \mu^{-1}$ this estimate yields
\[
\|I_{1}\mathbf{1}_{|x|>R}\|_{L_{x}^{\infty}}\leq 4\,\underset{|x|>R/2}{\sup}\,\int_{\complement B(0,\mu^{-1})}\frac{1}{|x-y|^{2}|y|^{2}}\ud y\,\underset{\mu\to 0}{\to}\,0.
\]
It remains to prove that $I_{1}\mathbf{1}_{|x|\leq R}$ converges to $0$ in $L_{x}^{3/2}$-norm as $\lambda,\mu\to 0$.
For $r\in(0,R)$ we split this quantity as follows
\[
\|I_{1}\mathbf{1}_{|x|\leq R}\|_{L_{x}^{3/2}}\leq\|I_{1}\mathbf{1}_{|x|\geq r}\|_{L_{x}^{3/2}}+\|I_{1}\mathbf{1}_{r<|x|\leq R}\|_{L_{x}^{3/2}}.
\]
We have
\[
\Bigl |(\alpha^{\lambda}\star\alpha^{\lambda})\star \Bigl(\frac{\chi^{\mu}}{|\cdot|^{2}}\star\frac{\chi^{\mu}}{|\cdot|^{2}}-\frac{1}{|\cdot|^{2}}\star\frac{1}{|\cdot|^{2}} \Bigr)\mathbf{1}_{|x|\leq r} \Bigr|\leq 2C\left(\alpha^{\lambda}\star\alpha^{\lambda}\right)\star\Sigma^{0}\mathbf{1}_{|x|\leq r}
\]
and we have already seen that $C(\alpha^{\lambda}\star\alpha^{\lambda})\star\Sigma^{0}\mathbf{1}_{|x|\leq r}$ converges to $\Sigma^{0}\mathbf{1}_{|x|\leq r}$ in the $L_{x}^{3/2}$-norm for any $0<r<\infty$.
Let $\eta>0$. We can choose $r=r(\eta)>0$ small enough and, next, find $\lambda(\eta)$ small enough so that for any $0<\lambda<\lambda(\eta)$, we get
\[
\|I_{1}\mathbf{1}_{|x|\leq r}\|_{L_{x}^{3/2}}\leq 2\|(C(\alpha^{\lambda}\star\alpha^{\lambda})\star\Sigma^{0}-\Sigma^{0})\mathbf{1}_{|x|\leq r}\|_{L_{x}^{3/2}}+2\|\Sigma^{0}\mathbf{1}_{|x|\leq r}\|_{L_{x}^{3/2}}
\leq \eta.
\]
Finally, the $L_{x}^{3/2}$-norm of $I_{1}\mathbf{1}_{r<|x|\leq R}$ can be estimated as we did for the $L_{x}^{\infty}$-norm of $I_{1}\mathbf{1}_{|x|>R}$.
Possibly at the price of taking $\lambda(\eta)$ smaller,
if $|x|>r$ we have $|x-y|>r/2$ for any $y\in \mathrm{supp}(a^\lambda\star a^\lambda)$. It follows that
\[
\|I_{1}\mathbf{1}_{r<|x|\leq R}\|_{L_{x}^{3/2}}\leq\mathrm{meas}\left(B(0,R)\right)^{2/3}\,\underset{r/2<|x|\leq R}{\sup}\,\int_{\complement B(0,\mu^{-1})}\frac{1}{|x-y|^{2}|y|^{2}}\ud y,
\]
which can be made $\leq \eta$ for $0<\mu<\mu(\eta)$, with $\mu(\eta)$ small enough.
This ends the proof.
\end{ProofOf}

\appendix
\section{Cauchy theory}\label{CauchyTheory}

From an energetic point of view, the natural functional spaces for the Cauchy theory of the Schrödinger-Wave equation are $C^{0}([0,T];H^{1}(\R_{x}^{d}))$ for the wave function $u$ and
\[
\mathcal{E}_{T}=C^{0}\left([0,T];L^{2}\left(\R_{x}^{d};\overbigdot{H}^{1}(\R_{z}^{n})\right)\right)\;\cap\;C^{1}\left([0,T];L^{2}\left(\R_{x}^{d};L^{2}(\R_{z}^{n})\right)\right)
\]
for the vibrational environment $\psi$.
We are going to prove the global existence of solutions to \eqref{schro}--\eqref{wave_schro}, with Cauchy data \eqref{CI_schro},
in these spaces, see Theorem~\ref{TC}.
Throughout this appendix, we work, without loss of generality, with $c=1$.

The proof of this theorem is quite classical: the most important part consists in applying Strichartz' estimates to the Schrödinger and the wave equation. In fact the main difficulty comes from the fact that Strichartz' estimates for \eqref{schro} lead to estimates of $u$ in $L_{t}^{q}L_{x}^{r}$ norms whereas Strichartz' estimates for \eqref{wave_schro} lead to estimates of $\psi$ in $L_{x}^{r}L_{t}^{q}L_{z}^{p}$ norms. In order to combine these two estimates
of different type,
we need to permute Lebesgue-norms in time and space. For that purpose we will use Hölder and Young inequalities (and the fact that $\sigma_{1}$ and $\sigma_{2}$ are in any $L^{p}$ space for $1\leq p\leq+\infty$) in order to work with $L_{t}^{q}L_{x}^{q}$ norms.

Let us introduce some notation that we will use until the end of this section. First we denote by $S$ the linear Schrödinger's group and by $(W,\overbigdot{W})$ the free wave group: for any $u_{0}\in L^{2}(\R_{x}^{d})$, $S(t)u_{0}$ is the unique solution at time $t$ of
\[
\left\{\begin{array}{l}
\ds i\partial_{t}u+\Delta_{x}u=0,\\
\ds u(0,x)=u_{0}(x),
\end{array}\right.
\]
and for any $(\psi_{0},\psi_{1})\in L^{2}(\R_{x}^{d};\overbigdot{H}^{1}(\R_{z}^{n}))\times L^{2}(\R_{x}^{d};L^{2}(\R_{z}^{n}))$, $\overbigdot{W}(t)\psi_{0}+W(t)\psi_{1}$ is the unique solution at time $t$ of
\[
\left\{\begin{array}{l}
\ds \partial_{tt}^{2}\psi-\Delta_{z}\psi=0,\\[3pt]
\ds (\psi(0,x,z),\partial_{t}\psi(0,x,z))=(\psi_{0}(x,z),\psi_{1}(x,z)).
\end{array}\right.
\]
With these notation we can now define (at least formally) the functions $\mathcal{L}$, $\mathcal{K}$ and~$\Phi$~by
\[
\left\{\begin{array}{l}
\ds \mathcal{L}(u,\psi):t\longmapsto S(t)u_{0}+\int_{0}^{t}S(t-s)\left[\left(\sigma_{1}\star_{x}\int\sigma_{2}\psi(s)\ud z\right)u(s)\right]\ud s,\\[10pt]
\ds\mathcal{K}(u,\psi):t\longmapsto \overbigdot{W}(t)\psi_{0}+W(t)\psi_{1}+\int_{0}^{t}W(t-s)\left[-\sigma_{2}\sigma_{1}\star_{x}|u(s)|^{2}\right]\ud s,\\
\Phi=(\mathcal{L},\mathcal{K}),
\end{array}\right.
\]
where $u_{0}\in H^{1}(\R_{x}^{d})$ and $(\psi_{0},\psi_{1})\in L^{2}(\R_{x}^{d};\overbigdot{H}^{1}(\R_{z}^{n}))\times L^{2}(\R_{x}^{d};L^{2}(\R_{z}^{n}))$ are now fixed until the end of this section. From here it is obvious that any fixed point $(u,\psi)$ of~$\Phi$ defines a solution of \eqref{schro}--\eqref{wave_schro} and $\eqref{CI_schro}$. In order to apply the Banach-Picard fixed point theorem we have to specify on which space we define the function $\Phi$. As already mentioned, since we wish to apply Strichartz estimates, we need that $\Phi$ is defined on a well adapted space for this approach. We introduce the following notations and spaces for that purpose. First let us define the Lebesgue exponent $p_{0}$ by
\begin{equation}
\label{p0}p_{0}=\frac{2n}{n-2}.
\end{equation}
Then, for any final time $T>0$ we introduce the following Banach spaces:
\[
X_{T}=L^{\infty}(0,T;H^{1}(\R_{x}^{d})),\quad Y_{T}=L^{2}(\R_{x}^{d};L^{\infty}(0,T;L^{p_{0}}(\R_{z}^{n})))\quad\text{and}\quad Z_{T}=X_{T}\times Y_{T},
\]
endowed with the norm $\|u,\psi\|_{Z_{T}}=\|u\|_{X_{T}}+\|\psi\|_{Y_{T}}$.

We introduce these spaces because $(\infty,2)$ is a \emph{Schrödinger-admissible} pair and $(\infty,p_{0})$ is a \emph{wave-admissible} pair for $n\geq 3$. Let us briefly recall what are the definition of Schrödinger and wave-admissible pairs and what are Strichartz' estimates (we~follow~\cite{MT} and the interested reader can find further information about Strichartz' estimates in \cite{Ginibre_Velo_Strichartz} and the references therein).

\skpt
\begin{defin}
\begin{enumeratei}
\item
We say that the exponent pair $(q,r)$ is \emph{Schrödinger-admissible} if $d\geq 1$, $q,r\geq 2$, $(q,r,d)\neq(2,\infty,2)$ and
\[
\frac{1}{q}+\frac{d}{2r}=\frac{d}{4}.
\]
\item
We say that the exponent pair $(q,p)$ is \emph{wave-admissible} if $n\geq 2$, $q,p\geq 2$, $(q,p,n)\neq(2,\infty,3)$ and
\[
\frac{1}{q}+\frac{n-1}{2p}\leq\frac{n-1}{4}.
\]
\end{enumeratei}
\end{defin}

From now on for any exponent $a\geq 1$, $a'$ will denote its conjugate exponent: $1/a+1/a'=1$.

\skpt
\begin{proposition}[Strichartz estimates]\label{strichartz}
\begin{enumeratei}
\item
Let $(q,r)$ and $(\bar{q},\bar{r})$ be Schrödinger-admissible pairs, $u_{0}\in L^{2}(\R_{x}^{d})$, $F\in L^{\bar{q}'}(0,T;L^{\bar{r}'}(\R_{x}^{d}))$ and let us denoted by $u$ the unique solution of $i\partial_{t}u+\Delta_{x}u=F$ with initial data $u_{0}$. Then there exists a constant $C>0$ independent of $T$ such that
\begin{equation}\label{schro_strichartz}
\|u\|_{L_{t}^{q}L_{x}^{r}}\leq C\bigl(\|u_{0}\|_{L_{x}^{2}}+\|F\|_{L_{t}^{\bar{q}'}L_{x}^{\bar{r}'}}\bigr).
\end{equation}

\item
Let $(q,p)$ and $(\bar{q},\bar{p})$ be wave-admissible pairs with $p,\bar{p}<+\infty$, $(\psi_{0},\psi_{1})\in\overbigdot{H}^{s}(\R_{z}^{n})\times\overbigdot{H}^{s-1}(\R_{z}^{n})$, $G\in L^{\bar{q}'}(0,T;L^{\bar{p}'}(\R_{z}^{n}))$ and let us denoted by $\psi$ the unique solution of $\partial_{tt}^{2}\psi-\Delta_{z}\psi=G$ with initial data $(\psi_{0},\psi_{1})$. Then, under the additional condition
\begin{equation}\label{ad_strichartz}
\frac{1}{q}+\frac{n}{p}=\frac{n}{2}-s=\frac{1}{\bar{q}'}+\frac{n}{\bar{p}'}-2,
\end{equation}
there exists a constant $K>0$ independent of $T$ such that
\begin{equation}\label{wave_strichartz}
\|\psi\|_{L_{t}^{q}L_{z}^{p}}+\|\psi\|_{L_{t}^{\infty}\overbigdot{H}_{z}^{s}}+\|\partial_{t}\psi\|_{L_{t}^{\infty}\overbigdot{H}_{z}^{s-1}}\leq K\bigl(\|\psi_{0}\|_{\overbigdot{H}_{z}^{s}}+\|\psi_{1}\|_{\overbigdot{H}_{z}^{s-1}}+\|G\|_{L_{t}^{\bar{q}'}L_{z}^{\bar{p}'}}\bigr).
\end{equation}
\end{enumeratei}
\end{proposition}

\begin{rmk}
We will apply \eqref{wave_strichartz} with the Sobolev regularity $s=1$. With this regularity the exponent pairs $(q,p)=(\infty,p_{0})$ and $(\infty,2)$ are wave-admissible and satisfies the additional condition \eqref{ad_strichartz}.
\end{rmk}

The following two Lemma justify that the application $\Phi$ is well defined on $Z_{T}$, sends $Z_{T}$ into itself and admits a fixed point on it.

\begin{lemma}\label{lemma_TC1}
There exists a constant $C>0$ independent of $T$ such that
\begin{subequations}
\begin{align}
\label{estimate_strichartz_1} &\ds\|\mathcal{L}(u,\psi)\|_{L_{t}^{\infty}L_{x}^{2}} \leq  C\bigl(\|u_{0}\|_{L_{x}^{2}}+|T|\|\psi\|_{Y_{T}}\|u\|_{L_{t}^{\infty}L_{x}^{2}}\bigr),\\
\label{estimate_strichartz_2} &\ds\|\nabla_{x}\mathcal{L}(u,\psi)\|_{L_{t}^{\infty}L_{x}^{2}} \leq  C\bigl(\|\nabla_{x}u_{0}\|_{L_{x}^{2}}+|T|\|\psi\|_{Y_{T}}\left[\|u\|_{L_{t}^{\infty}L_{x}^{2}}+\|\nabla_{x}u\|_{L_{t}^{\infty}L_{x}^{2}}\right]\bigr),\\
\label{estimate_strichartz_3}&
\|\mathcal{K}(u,\psi)\|_{Y_{T}}+\|\psi\|_{L_{x}^{2}L_{t}^{\infty}\overbigdot{H}_{z}^{1}}+\|\partial_{t}\psi\|_{L_{x}^{2}L_{t}^{\infty}L_{z}^{2}}\\
&\hspace*{4cm}\leq C\bigl(\|\psi_{0}\|_{L_{x}^{2}\overbigdot{H}_{z}^{1}}+\|\psi_{1}\|_{L_{x}^{2}L_{z}^{2}}+|T|\|u\|_{L_{t}^{\infty}L_{x}^{2}}^{2}\bigr),\notag
\end{align}
\end{subequations}
and
\begin{subequations}
\begin{align}
\label{estimate_strichartz_4} &\|\mathcal{L}(u,\psi)-\mathcal{L}(v,\varphi)\|_{L_{t}^{\infty}L_{x}^{2}}\\
&\hspace*{2cm}\leq  C\,|T|\left(\|\psi\|_{Y_{T}}\|u-v\|_{L_{t}^{\infty}L_{x}^{2}}+\|\psi-\varphi\|_{Y_{T}}\|v\|_{L_{t}^{\infty}L_{x}^{2}}\right),\notag\\
\label{estimate_strichartz_5}&
\|\nabla_{x}(\mathcal{L}(u,\psi)-\mathcal{L}(v,\varphi))\|_{L_{t}^{\infty}L_{x}^{2}} \leq  C\,|T|\Bigl(\|\psi\|_{Y_{T}}\Bigl[\|u-v\|_{L_{t}^{\infty}L_{x}^{2}}\\
&\hspace*{2cm}+\|\nabla_{x}(u-v)\|_{L_{t}^{\infty}L_{x}^{2}} \Bigr]+\|\psi-\varphi\|_{Y_{T}}\bigl[\|v\|_{L_{t}^{\infty}L_{x}^{2}}+\|\nabla_{x}v\|_{L_{t}^{\infty}L_{x}^{2}}\bigr]\Bigr)\notag
\\
\label{estimate_strichartz_6}
&\|\mathcal{K}(u,\psi)-\mathcal{K}(v,\varphi)\|_{Y_{T}} \leq  C|T|\left(\|u\|_{L_{t}^{\infty}L_{x}^{2}}+\|v\|_{L_{t}^{\infty}L_{x}^{2}}\right)\|u-v\|_{L_{t}^{\infty}L_{x}^{2}}.
\end{align}
\end{subequations}
\end{lemma}

\begin{lemma}\label{lemma_TC2}
There exists a universal constant $C_{1}>0$ such that for any final time $T>0$ small enough, $\Phi:B_{T}\to B_{T}$, where
\[
B_{T}=\left\{(u,\psi)\in Z_{T}\sep \|u,\psi\|_{Z_{T}}\leq C_{1}(\|u_{0}\|_{H_{x}^{1}}+\|\psi_{0}\|_{L_{z}^{2}\overbigdot{H}_{z}^{1}}+\|\psi_{1}\|_{L_{x}^{2}L_{z}^{2}})\right\}.
\]
Moreover, considering smaller $T$ if necessary, $\Phi$ is indeed a contraction on $B_{T}$.
\end{lemma}

We postpone the proof of Lemma~\ref{lemma_TC1} to the end of this section and we start by proving Lemma~\ref{lemma_TC2} and Theorem~\ref{TC}.

\begin{ProofOf}{Lemma~\ref{lemma_TC2}}
We can summarize the estimates \eqref{estimate_strichartz_1}--\eqref{estimate_strichartz_3} as follows:
\[
\|\Phi(u,\psi)\|_{Z_{T}}\leq C\Bigl[\|u_{0}\|_{H_{x}^{1}}+\|\psi_{0}\|_{L_{x}^{2}\overbigdot{H}_{z}^{1}}+\|\psi_{1}\|_{L_{x}^{2}L_{z}^{2}}+|T|\,\|u,\psi\|_{Z_{T}}^{2}\Bigr].
\]
Next, let $C_{1}=2C$; we thus obtain that for any $(u,\psi)\in B_{T}$,
\begin{multline*}
\|\Phi(u,\psi)\|_{Z_{T}}\leq C\left[1+C_{1}^{2}\,|T|\, \bigl(\|u_{0}\|_{H_{x}^{1}}+\|\psi_{0}\|_{L_{x}^{2}\overbigdot{H}_{z}^{1}}+\|\psi_{1}\|_{L_{x}^{2}L_{z}^{2}} \bigr)\right]\\
\times\bigl(\|u_{0}\|_{H_{x}^{1}}+\|\psi_{0}\|_{L_{x}^{2}\overbigdot{H}_{z}^{1}}+\|\psi_{1}\|_{L_{x}^{2}L_{z}^{2}}\bigr).
\end{multline*}
Since for $T$ small enough,
\[
C_{1}^{2}\,|T|\, \bigl(\|u_{0}\|_{H_{x}^{1}}+\|\psi_{0}\|_{L_{x}^{2}\overbigdot{H}_{z}^{1}}+\|\psi_{1}\|_{L_{x}^{2}L_{z}^{2}} \bigr)<1,
\]
we obtain that $\Phi$ sends $B_{T}$ into $B_{T}$ for $T$ small enough. As previously, we can recast \eqref{estimate_strichartz_4}--\eqref{estimate_strichartz_6} as follows:\vspace*{-3pt}
\[
\|\Phi(u,\psi)-\Phi(v,\phi)\|_{Z_{T}}\leq C\,|T|\,\left(\|(u,\psi)\|_{Z_{T}}+\|v,\phi\|_{Z_{T}}\right)\|(u,\psi)-(v,\phi)\|_{Z_{T}}.
\]
Therefore, for any $(u,\psi),(v,\phi)\in B_{T}$,\vspace*{-5pt}
\begin{multline*}
\|\Phi(u,\psi)-\Phi(v,\phi)\|_{Z_{T}}
\\
\leq 2C\,C_{1} \bigl(\|u_{0}\|_{H_{x}^{1}}+\|\psi_{0}\|_{L_{z}^{2}\overbigdot{H}_{z}^{1}}+\|\psi_{1}\|_{L_{x}^{2}L_{z}^{2}} \bigr)\,|T|\|(u,\psi)-(v,\phi)\|_{Z_{T}},
\end{multline*}
holds and $\Phi$ is a contraction as soon as $T$ is small enough.
\end{ProofOf}

\begin{ProofOf}{Theorem~\ref{TC}. Step 1: Local existence} For $T$ small enough $\Phi$ is a contraction on $B_{T}$, we thus know that \eqref{schro}--\eqref{wave_schro} has a solution in $Z_{T}$. Then it is clear that for any solution $(u,\psi)\in Z_{T}$ of \eqref{schro}--\eqref{wave_schro},
\begin{gather*}
u\in L^{\infty}(0,T;H^{1}(\R_{x}^{d})),\quad \psi\in L^{2}\bigl(\R_{x}^{d};L^{\infty}(0,T;\overbigdot{H}^{1}(\R_{z}^{n}))\bigr),\\
\tag*{and} \partial_{t}\psi\in L^{2}\bigl(\R_{x}^{d};L^{\infty}(0,T;L^{2}(\R_{z}^{n}))\bigr)
\end{gather*}
(for $\psi$ it comes from the Strichartz estimate \eqref{estimate_strichartz_3}). Moreover, using the fact that $(u,\psi)$ is a fixed point of $\Phi$ and the expressions of $\mathcal{L}$ and $\mathcal{K}$ in terms of $S$ and $(W,\overbigdot{W})$, one can prove that indeed $u\in C^{0}\left([0,T];H^{1}(\R_{x}^{d})\right)$, for almost every $x\in\R^{d}$,\vspace*{-3pt}
\begin{align*}
(t,z)&\mto \psi(t,x,z)\in C^{0}\bigl([0,T];\overbigdot{H}^{1}(\R_{z}^{n})\bigr)\\
\tag*{and}
(t,z)&\mto\partial_{t}\psi(t,x,z)\in C^{0}\bigl([0,T];L^{2}(\R_{z}^{n})\bigr).
\end{align*}
We finish the proof by applying the following lemma (proved at the end of this section) to $\psi$ and $\partial_{t}\psi$ in order to obtain that $\psi\in\mathcal{E}_{T}$.
\begin{lemma}\label{lemma_continuity}
If $f\in L_{x}^{2}L_{t}^{\infty}$ and for almost every $x\in\R^{d}$, $t\mto f(t,x)\in C^{0}([0,T])$, then $f\in C^{0}\left([0,T];L^{2}(\R_{x}^{d})\right)$.
\end{lemma}

\subsubsection*{Step 2: Uniqueness} The uniqueness in $B_{T}$ comes from the fixed point theorem and we can extend this uniqueness statement to the entire space $Z_{T}$. Then the uniqueness in $C_{t}^{0}H_{x}^{1}\times\mathcal{E}_{T}$ comes from the fact that any fixed point $(u,\psi)\in C_{t}^{0}H_{x}^{1}\times\mathcal{E}_{T}$ of $\Phi$ is also an element of $Z_{T}$ (thanks to the estimate \eqref{estimate_strichartz_3}, we get that $\psi$ is in $Y_{T}$).

\subsubsection*{Step 3: Global existence} Since the time $T$ in Lemma~\ref{lemma_TC2} depends only on universal constants and on\vspace*{-3pt}
\[
\|u_{0}\|_{H_{x}^{1}}+\|\psi_{0}\|_{L_{x}^{2}\overbigdot{H}_{z}^{1}}+\|\psi_{1}\|_{L_{x}^{2}L_{z}^{2}},
\]
the first two steps of this proof allow us to obtain the following proposition.
\begin{proposition}
Let $n\geq 3$. Then, for any $u_{0}\in H^{1}(\R_{x}^{d})$ and any $(\psi_{0},\psi_{1})\in L^{2}(\R_{x}^{d};\overbigdot{H}^{1}(\R_{z}^{n}))\times L^{2}(\R_{x}^{d};L^{2}(\R_{z}^{n}))$, there exists $T^{\star}>0$ such that for any $T$ with $0<T<T^{\star}$, the problem \eqref{schro}--\eqref{wave_schro} and \eqref{CI_schro} admits a unique solution \hbox{$(u,\psi)\in C^{0}\bigl([0,T];H^{1}(\R_{x}^{d})\bigr)\times\mathcal{E}_{T}$} on $[0,T]$. Moreover, if for some $0<T\leq T^{\star}$,
\[
\underset{t\nearrow T}{\limsup}\,\|u(t)\|_{H_{x}^{1}}+\|\psi(t)\|_{L_{x}^{2}\overbigdot{H}_{z}^{1}}+\|\partial_{t}\psi(t)\|_{L_{x}^{2}L_{z}^{2}}<+\infty,
\]
then, actually, $T<T^{\star}$.
\end{proposition}
Then in order to obtain the global existence we have to justify that the quantity
\[
\|u(t)\|_{H_{x}^{1}}+\|\psi(t)\|_{L_{x}^{2}\overbigdot{H}_{z}^{1}}+\|\partial_{t}\psi(t)\|_{L_{x}^{2}L_{z}^{2}}
\]
does not blow up in finite time. Thanks to the mass conservation of the wave function $u$ ($M=\|u(t)\|_{L_{x}^{2}}$ is constant in time) and thanks to \eqref{estimate_strichartz_3} we get
\[
\|u(t)\|_{H_{x}^{1}}+\|\psi(t)\|_{L_{x}^{2}\overbigdot{H}_{z}^{1}}+\|\partial_{t}\psi(t)\|_{L_{x}^{2}L_{z}^{2}}\!\lesssim\! M+\|\nabla_{x}u(t)\|_{L_{x}^{2}}+\|\psi_{0}\|_{L_{x}^{2}\overbigdot{H}_{z}^{1}}+\|\psi_{1}\|_{L_{x}^{2}L_{z}^{2}}+|t|M,
\]
and it only remains to control $\|\nabla_{x}u(t)\|_{L_{x}^{2}}$. For that purpose we use the energy conservation
\eqref{energy_Sch}
in order to obtain
\[
\frac{1}{2}\|\nabla_{x}u(t)\|_{L_{x}^{2}}+\int\left(\sigma_{1}\star\int\sigma_{2}\psi(t)\ud z\right)|u(t)|^{2}\ud x\leq \mathcal E_{\mathrm {Schr}}(t)= \mathcal E_{\mathrm {Schr}}(0).
\]
Then if $\|\nabla_{x}u(t)\|_{L_{x}^{2}}$ blows up in finite time, $|\int(\sigma_{1}\star\int\sigma_{2}\psi(t)\ud z)|u(t)|^{2}\ud x|$ has to blows up in finite time too. But
\begin{equation}\label{estimate_utile}
\begin{split}
\biggl\|\int(\sigma_{1}&\star\int\sigma_{2}\psi\ud z)|u|^{2}\ud x\biggr\|_{L_{t}^{\infty}}\leq M^{2}\,\left\|\sigma_{1}\star\int\sigma_{2}\psi\ud z\right\|_{L_{t}^{\infty}L_{x}^{\infty}}\\
&=M^{2}\,\left\|\sigma_{1}\star\int\sigma_{2}\psi\ud z\right\|_{L_{x}^{\infty}L_{t}^{\infty}}\leq M^{2}\|\sigma_{2}\|_{L_{z}^{p_{0}'}}\,\left\|\sigma_{1}\star\|\psi\|_{L_{z}^{p_{0}}}\right\|_{L_{x}^{\infty}L_{t}^{\infty}}\\
&\leq M^{2}\|\sigma_{2}\|_{L_{z}^{p_{0}'}}\bigl\|\sigma_{1}\star\|\psi\|_{L_{t}^{\infty}L_{z}^{p_{0}}}\bigr\|_{L_{x}^{\infty}}\leq M^{2}\|\sigma_{2}\|_{L_{z}^{p_{0}'}}\|\sigma_{1}\|_{L_{x}^{2}}\,\|\psi\|_{L_{x}^{2}L_{t}^{\infty}L_{z}^{p_{0}}},
\end{split}
\end{equation}
and eventually estimate \eqref{estimate_strichartz_3} tells us that $|\int(\sigma_{1}\star\int\sigma_{2}\psi(t)\ud z)|u(t)|^{2}\ud x|$ grows at most linearly in time.
\end{ProofOf}

\begin{rmk}
In fact the proof of the global existence gives us the additional information that the quantities
\[
\|\nabla_{x}u(t)\|_{L_{x}^{2}},\quad \|\psi(t)\|_{L_{x}^{2}\overbigdot{H}_{z}^{1}}+\|\partial_{t}\psi(t)\|_{L_{x}^{2}L_{z}^{2}},\quad\text{and}\quad
\biggl|\int\biggl(\sigma_{1}\star\int\sigma_{2}\psi(t)\ud z\biggr)|u(t)|^{2}\ud x\biggr|
\]
grow at most linearly in time.
\end{rmk}

We finish this section with the proofs of Lemma~\ref{lemma_TC1} and Lemma~\ref{lemma_continuity}.

\begin{ProofOf}{Lemma~\ref{lemma_TC1}. Estimate \eqref{estimate_strichartz_1}} We apply apply the Strichartz estimate \eqref{schro_strichartz} to $\mathcal{L}(u,\psi)$ with the Schrödinger-admissible pair $(\infty,2)$ on both side to obtain
\[
\|\mathcal{L}(u,\psi)\|_{L_{t}^{\infty}L_{x}^{2}}\lesssim\|u_{0}\|_{L_{x}^{2}}+\left\|\left(\sigma_{1}\star_{x}\int\sigma_{2}\psi\ud z\right)u\right\|_{L_{t}^{1}L_{x}^{2}}.
\]
Then, thanks to the following estimate
\begin{align*}
\left\|\left(\sigma_{1}\star_{x}\int\sigma_{2}\psi\ud z\right)u\right\|_{L_{t}^{1}L_{x}^{2}}&\leq |T|\left\|\left(\sigma_{1}\star_{x}\int\sigma_{2}\psi\ud z\right)u\right\|_{L_{t}^{\infty}L_{x}^{2}}\\
&\leq\left\|\sigma_{1}\star_{x}\int\sigma_{2}\psi\ud z\right\|_{L_{t}^{\infty}L_{x}^{\infty}}\,\|u\|_{L_{t}^{\infty}L_{x}^{2}},
\end{align*}
and thanks to \eqref{estimate_utile}, we eventually obtain
\[
\|\mathcal{L}(u,\psi)\|_{L_{t}^{\infty}L_{x}^{2}}\lesssim\|u_{0}\|_{L_{x}^{2}}+|T|\,\|\psi\|_{Y_{T}}\,\|u\|_{L_{t}^{\infty}L_{x}^{2}}.
\]

\subsubsection*{Estimate \eqref{estimate_strichartz_2}} Since\vspace*{-3pt}
\begin{multline*}
\ds\nabla_{x}\mathcal{L}(u,\psi)(t)=S(t)\nabla_{x}u_{0}\\[-3pt]
\ds+\int_{0}^{t}S(t-s)\left[\left(\nabla_{x}\sigma_{1}\star\int\sigma_{2}\psi(s)\ud z\right)u(s)+\left(\sigma_{1}\star\int\sigma_{2}\psi(s)\ud z\right)\nabla_{x}u(s)\right]\ud s,
\end{multline*}
we just apply the same estimates as before.

\subsubsection*{Estimate \eqref{estimate_strichartz_3}} We apply for almost every $x\in\R^{d}$ the Strichartz estimate \eqref{wave_strichartz} to $\mathcal{K}(u,\psi)(x)$ with the wave-admissible pair $(\infty,p_{0})$ on the left hand side and $(\infty,2)$ on the right hand side\vspace*{-3pt}
\begin{multline*}
\|\mathcal{K}(u,\psi)(x)\|_{L_{t}^{\infty}L_{z}^{p_{0}}}+\|\psi(x)\|_{L_{t}^{\infty}\overbigdot{H}_{z}^{1}}+\|\partial_{t}\psi(x)\|_{L_{t}^{\infty}L_{z}^{2}}\\
\lesssim\,\|\psi_{0}(x)\|_{\overbigdot{H}_{z}^{1}}+\|\psi_{1}(x)\|_{L_{z}^{2}}+\left\|\sigma_{2}\,\sigma_{1}\star|u|^{2}(x)\right\|_{L_{t}^{1}L_{z}^{2}}.
\end{multline*}
Then, since\vspace*{-3pt}
\[
\left\|\sigma_{2}\,\sigma_{1}\star|u|^{2}(x)\right\|_{L_{t}^{1}L_{z}^{2}}=\|\sigma_{2}\|_{L_{z}^{2}}\,\|\sigma_{1}\star|u|^{2}(x)\|_{L_{t}^{1}}\leq\|\sigma_{2}\|_{L_{z}^{2}}\,|\sigma_{1}|\star\|u\|_{L_{t}^{2}}^{2}(x),
\]
we can pass in $L_{x}^{2}$-norm to obtain\vspace*{-3pt}
\[
\left\|\sigma_{2}\,\sigma_{1}\star|u|^{2}\right\|_{L_{x}^{2}L_{t}^{1}L_{z}^{2}}\leq\|\sigma_{2}\|_{L_{z}^{2}}\,\bigl\||\sigma_{1}|\star\|u\|_{L_{t}^{2}}^{2}\bigr\|_{L_{x}^{2}}.
\]
Here, thanks to the Young inequality we have
\[
\bigl\||\sigma_{1}|\star\|u\|_{L_{t}^{2}}^{2}\bigr\|_{L_{x}^{2}}\leq \|\sigma_{1}\|_{L_{x}^{2}}\,\bigl\|\|u\|_{L_{t}^{2}}^{2}\bigr\|_{L_{x}^{1}}=\|\sigma_{1}\|_{L_{x}^{2}}\,\|u\|_{L_{t}^{2}L_{x}^{2}}^{2}\leq \|\sigma_{1}\|_{L_{x}^{2}}\,|T|\,\|u\|_{L_{t}^{\infty}L_{x}^{2}}^{2},
\]
and we eventually obtain\vspace*{-3pt}
\[
\|\mathcal{K}(u,\psi)\|_{L_{x}^{2}L_{t}^{\infty}L_{z}^{p_{0}}}+\|\psi\|_{L_{x}^{2}L_{t}^{\infty}\overbigdot{H}_{z}^{1}}+\|\partial_{t}\psi\|_{L_{x}^{2}L_{t}^{\infty}L_{z}^{2}}\!\lesssim\!\|\psi_{0}\|_{L_{x}^{2}\overbigdot{H}_{z}^{1}}+\|\psi_{1}\|_{L_{x}^{2}L_{z}^{2}}+|T|\,\|u\|_{L_{t}^{\infty}L_{x}^{2}}^{2}.
\]

\subsubsection*{Estimates \eqref{estimate_strichartz_4}, \eqref{estimate_strichartz_5} and \eqref{estimate_strichartz_6}} Since\vspace*{-5pt}
\begin{multline*}
\ds\mathcal{L}(u,\psi)(t)-\mathcal{L}(v,\varphi)(t)
= \ds\int_{0}^{t}S(t-s)\left[\left(\sigma_{1}\star_{x}\int\sigma_{2}\psi(s)\ud z\right)(u(s)-v(s))
\right.
\\
\left.+\left(\sigma_{1}\star_{x}\int\sigma_{2}(\psi(s)-\varphi(s))\ud z\right)v(s)\right]\ud s
\end{multline*}\vspace*{-5pt}
and
\begin{multline*}
\ds\mathcal{K}(u,\psi)(t)-\mathcal{K}(v,\varphi)(t)\\
= \ds\int_{0}^{t}W(t-s)\left[-\sigma_{2}\sigma_{1}\star_{x}\left([u(s)-v(s)]\bar{u}(s)+v(s)[\bar{u}(s)-\bar{v}(s)]\right)\right]\ud s,
\end{multline*}
we just follow closely the proof of \eqref{estimate_strichartz_1}, \eqref{estimate_strichartz_2} and \eqref{estimate_strichartz_3}.
\end{ProofOf}

\begin{ProofOf}{Lemma~\ref{lemma_continuity}}
Let us fix $\varepsilon>0$ and $t\in[0,T]$. We know that for all $x\in\R^{d}$ and for all $\eta>0$, there exists $\delta(\eta,t,x)\geq 0$ such that if $|t-s|\leq\delta(\eta,t,x)$, then $|f(t,x)-f(s,x)|\leq\eta$. Note that in fact $\delta(\eta,t,x)$ is positive
for almost every $x\in\R^{d}$. Moreover, since $f\in L_{x}^{2}L_{t}^{\infty}$ we now that
\[
\int_{\R^{d}}\mathbf{1}_{|x|\geq R}\|f(x)\|_{L_{t}^{\infty}}^{2}\ud x\,\underset{R\to \infty}{\to}\,0.
\]
Let $\delta>0$. Let us also introduce the following subset of $\R_{x}^{d}$
\[
B_{t,\delta}^{R,\eta}=\bigl\{x\in\R^{d}\sep |x|\leq R \text{ and } \delta(\eta,t,x)
\leq\delta\bigr\}.
\]
Note that $\mathrm{meas}(B_{t,\delta}^{R,\eta})\to 0$ when $\delta\to 0$. Then for all $R,\eta,\delta>0$ and for all $s$ such that $|t-s|\leq\delta$,
\begin{align*}
\|f(t)-f(s)\|_{L_{x}^{2}}&\leq\|\mathbf{1}_{|x|\geq R}(f(t)-f(s))\|_{L_{x}^{2}}+\|\mathbf{1}_{|x|\leq R}(f(t)-f(s))\|_{L_{x}^{2}}\\
&\leq 2\,\|\mathbf{1}_{|x|\geq R}f\|_{L_{x}^{2}L_{t}^{\infty}}+\eta\,\mathrm{meas}\left(B(0,R)\right)^{1/2}+2\,\mathrm{meas}(B_{t,\delta}^{R,\eta})\|f\|_{L_{x}^{2}L_{t}^{\infty}}.
\end{align*}
We can pick $R$ large enough to obtain
$
2\,\|\mathbf{1}_{|x|\geq R}f\|_{L_{x}^{2}L_{t}^{\infty}}\leq\sfrac{\varepsilon}{3}$, then we fix $\eta$ small enough to get
$\eta\,\mathrm{meas}\left(B(0,R)\right)^{1/2}\leq\sfrac{\varepsilon}{3}$, and we eventually fix $\delta$ small enough to get
\[
2\,\mathrm{meas}(B_{t,\delta}^{R,\eta})\|f\|_{L_{x}^{2}L_{t}^{\infty}}\leq\frac{\varepsilon}{3}.\qedhere
\]
\end{ProofOf}

\section{Semi-classical analysis}\label{SClassicalAnalysis}

In this section we rescale the Schrödinger-Wave system as follows
\begin{subequations}
\begin{alignat}{11}
\label{schroh}&\ds ih\,\partial_{t}u_{h}+\frac{h^{2}}{2}\Delta_{x}u_{h}=\left(\sigma_{1}\star_{x}\int\sigma_{2}\psi_{h}(t)\ud z\right)u_{h},&\hspace{1cm}t\in\R,\;x\in\R^{d},\\
\label{wave1_schroh}&\ds \partial_{t}\psi_{h}=\chi_{h},&\hspace{1cm}t\in\R,\;x\in\R^{d},\;z\in\R^{n},\\
\label{wave2_schroh}&\ds \partial_{t}\chi_{h}=c^{2}\Delta_{z}\psi_{h}-c^{2}\sigma_{2}(z)\left(\sigma_{1}\star_{x}|u_{h}(t)|^{2}\right)(x),&\hspace{1cm}t\in\R,\;x\in\R^{d},\;z\in\R^{n},
\end{alignat}
\end{subequations}
where $h>0$ denotes (a dimensionless version of) the Planck constant.
We wish to investigate the behavior of this rescaled system when $h\to 0$. This is expected to establish a connection between the classical and quantum models, see \cite{Lions_Paul_Wigner}. More precisely for every $h>0$ we consider the Wigner transform of $u_{h}$
\[
W_{h}(t,x,\xi)=\frac{1}{(2\pi)^{d}}\int_{\R^{d}}e^{-i\xi\cdot y}u_{h}(t,x+\frac{h}{2}\,y)\,\bar{u}_{h}(t,x-\frac{h}{2}\,y)\ud y
\]
and we address the question of the asymptotic behavior of $(W_{h},\psi_{h},\chi_{h})$ when $h$ goes to $0$. Our goal is to prove that $(W_{h},\psi_{h},\chi_{h})$ admits a limit and this limit is a solution of the Vlasov-wave system \eqref{kin}--\eqref{wave_kin}. For that purpose let us introduce some notations and assumptions.

We consider a sequence of initial data $(u_{0}^{h})_{h>0}\subset H_{x}^{1}$, $(\psi_{0}^{h})_{h>0}\subset L_{x}^{2}\overbigdot{H}_{z}^{1}$ and $(\chi_{0}^{h})_{h>0}\subset L_{x}^{2}L_{z}^{2}$ such that

\begin{enumerate}\renewcommand{\theenumi}{H\arabic{enumi}}\setcounter{enumi}{4}
\item\label{h5}
the quantities $\|u_{h}\|_{L_{x}^{2}}$ and
\begin{multline*}
\mathscr{E}_{0,+}^{h}=\frac{h^{2}}{2}\int_{\R^{d}}|\nabla_{x}u_{0}^{h}|^{2}\ud x+\int_{\R^{d}}\left(\sigma_{1}\star\int_{\R^{n}}\sigma_{2}\psi_{0}^{h}\ud z\right)_{+}|u_{0}^{h}|^{2}\ud x\\
+\frac{1}{2c^{2}}\iint_{\R^{d}\times\R^{n}}|\chi_{0}^{h}|^{2}\ud x\ud z+\frac{1}{2}\iint_{\R^{d}\times\R^{n}}|\nabla_{z}\psi_{0}^{h}|^{2}\ud x\ud z
\end{multline*}
are uniformly bounded with respect to $h$.
\end{enumerate}

\skpt
\begin{rmk}\label{rmk_h5}
\begin{enumeratei}
\item
Assumption \eqref{h5} guarantees us that the sequences $(\psi_{0}^{h})$ and $(\chi_{0}^{h})$ are uniformly bounded with respect to $h$ respectively in $L_{x}^{2}\overbigdot{H}_{z}^{1}$ and $L_{x}^{2}L_{z}^{2}$. Hence, there exist $\psi_{0}\in L_{x}^{2}\overbigdot{H}_{z}^{1}$ and $\chi_{0}\in L_{x}^{2}L_{z}^{2}$ such that sub-sequences still labeled $(\psi_{0}^{h})_{h>0}$ and $(\chi_{0}^{h})_{h>0}$ converge respectively to $\psi_{0}$ in $L_{x}^{2}\overbigdot{H}_{z}^{1}$-weakly and $\chi_{0}$ in $L_{x}^{2}L_{z}^{2}$-weakly.

\item
Moreover, since the rescaled Hamiltonian
\begin{multline*}
\mathscr{E}^{h}(t)=\frac{h^{2}}{2}\int_{\R^{d}}|\nabla_{x}u_{h}(t)|^{2}\ud x+\int_{\R^{d}}\left(\sigma_{1}\star\int_{\R^{n}}\sigma_{2}\psi_{h}(t)\ud z\right)|u_{h}(t)|^{2}\ud x\\
+\frac{1}{2c^{2}}\iint_{\R^{d}\times\R^{n}}|\chi_{h}(t)|^{2}\ud x\ud z+\frac{1}{2}\iint_{\R^{d}\times\R^{n}}|\nabla_{z}\psi_{h}(t)|^{2}\ud x\ud z
\end{multline*}
is preserved by the system \eqref{schroh}--\eqref{wave2_schroh}, we have
\bgroup
\multlinegap0pt
\begin{multline*}
0\leq\frac{h^{2}}{2}\int_{\R^{d}}|\nabla_{x}u_{h}(t)|^{2}\ud x+\frac{1}{2c^{2}}\iint_{\R^{d}\times\R^{n}}|\chi_{h}(t)|^{2}\ud x\ud z+\frac{1}{2}\iint_{\R^{d}\times\R^{n}}|\nabla_{z}\psi_{h}(t)|^{2}\ud x\ud z\\
= \mathscr{E}^{h}(0)-\int_{\R^{d}}\left(\sigma_{1}\star\int_{\R^{n}}\sigma_{2}\psi_{h}(t)\ud z\right)|u_{h}(t)|^{2}\ud x\\
\leq \mathscr{E}_{0,+}^{h}-\int_{\R^{d}}\left(\sigma_{1}\star\int_{\R^{n}}\sigma_{2}\psi_{h}(t)\ud z\right)|u_{h}(t)|^{2}\ud x.
\end{multline*}
\egroup
Then thanks to \eqref{estimate_utile} coupled with the mass conservation of the wave function $u_{h}$ and \eqref{estimate_strichartz_3} we have
\begin{multline*}
\left\|\int_{\R^{d}}\left(\sigma_{1}\star\int_{\R^{n}}\sigma_{2}\psi_{h}(t)\ud z\right)|u_{h}(t)|^{2}\ud x\right\|_{L_{t}^{\infty}}\\
\lesssim\bigl(\|\psi_{0}^{h}\|_{L_{x}^{2}\overbigdot{H}_{z}^{1}}+\|\chi_{0}^{h}\|_{L_{x}^{2}L_{z}^{2}}+|T|\|u_{0}^{h}\|_{L_{x}^{2}}\bigr)\|u_{0}^{h}\|_{L_{x}^{2}}^{2},
\end{multline*}
that means $h^{2}\|\nabla_{x}u_{h}(t)\|_{L_{x}^{2}}^{2}$, $\|\chi_{h}(t)\|_{L_{x}^{2}L_{z}^{2}}$ and $\|\psi_{h}(t)\|_{L_{x}^{2}\overbigdot{H}_{z}^{1}}$ are uniformly bounded with respect to $h$ and $t\in[0,T]$.
\end{enumeratei}
\end{rmk}

One can easily check that
the Wigner transform $W_{h}$ associated to a solution $u_h$ of \eqref{schroh} satisfies the following equation
\begin{equation}\label{eq_wigner}
\partial_{t}W_{h}+\xi\cdot\nabla_{x}W_{h}+K_{h}\star_{\xi}W_{h}=0,
\end{equation}
where
\begin{equation}\label{K_h}
K_{h}(t,x,\xi)=\frac{i}{(2\pi)^{d}}\int_{\R^{d}}e^{-i\xi\cdot y}\frac{1}{h}\Bigl(\Phi_{h}\Bigl(t,x+\frac{h}{2}\,y\Bigr)-\Phi_{h}\Bigl(t,x-\frac{h}{2}\,y\Bigr)\Bigr)\ud y.
\end{equation}
This follows by direct inspection when $u_{h}$ is a strong solution of \eqref{schroh}, which is the case if $u_{0}^{h}$ is regular enough;
dealing with weak solutions requires a step by regularization and approximation.

According to \cite{Lions_Paul_Wigner}, we introduce the separable Banach space
\[
\mathcal{A}=\bigl\{\varphi\in C^{0}(\R_{x}^{d}\times\R_{\xi}^{d})\sep \mathcal{F}_{\xi}\varphi(x,y)\in L^{1}\bigl(\R_{y}^{d};C^{0}(\R_{x}^{d})\bigr)\bigr\}
\]
equipped with the norm
\[
\|\varphi\|_{\mathcal{A}}=\|\mathcal{F}_{\xi}\varphi\|_{L_{y}^{1}C_{x}^{0}}=\int_{\R^{d}}\underset{x}{\sup}\,|\mathcal{F}_{\xi}\varphi(x,y)|\ud y,
\]
and notice that the space
\[
\mathcal{B}=\left\{\varphi\in\mathcal{S}\sep \mathcal{F}_{\xi}\varphi\in C_{c}^{\infty}(\R_{x}^{d}\times\R_{y}^{d})\right\}
\]
is dense in $\mathcal{A}$. We also denote by $\mathcal{M}=\mathcal{M}(\R_{x}^{d}\times\R_{\xi}^{d})$ the space of bounded measures on $\R_{x}^{d}\times\R_{\xi}^{d}$,
and $\mathcal{M}_+$ its positive cone.

\begin{theo}\label{SC}
Let \eqref{h1}--\eqref{h2} and \eqref{h5} be fulfilled. Up to a sub-sequence, the families $(W_{h})_{h>0}$, $(\psi_{h})_{h>0}$ and $(\chi_{h})_{h>0}$ converge respectively to $\mu\in C^{0}([0,T];\mathcal{M}-w\star)$, $\psi\in C^{0}([0,T];L_{x}^{2}\overbigdot{H}_{z}^{1}-w)$ and $\chi\in C^{0}([0,T];L_{x}^{2}L_{z}^{2}-w)$ respectively in the spaces $C^{0}([0,T];\mathcal{A}'-w\star)$, $C^{0}([0,T] ; L_{x}^{2}\overbigdot{H}_{z}^{1}-w)$ and $C^{0}([0,T];L_{x}^{2}L_{z}^{2}-w)$. Moreover $(\mu,\psi,\chi)$ is a solution of the Vlasov-wave system
\begin{subequations}
\begin{alignat}{11}
&\ds \partial_{t}\mu+\mathrm{div}_{x}(\xi\mu)-\mathrm{div}_{\xi}\left(\nabla_{x}\left[\sigma_{1}\star_{x}\int\sigma_{2}\psi(t)\ud z\right]\mu\right)=0,&\hspace{.4cm}\text{ in }\mathcal{D}'\bigl((0,T);\mathcal{B}'\bigr),\notag\\
&\ds \partial_{t}\psi=\chi,&\hspace{.4cm}\text{ in }\mathcal{D}'\bigl((0,T)\times\R_{x}^{d}\times\R_{z}^{n}\bigr),\notag\\
&\ds \partial_{t}\chi=c^{2}\Delta_{z}\psi-\sigma_{2}(z)\left(\sigma_{1}\star_{x}\int\ud\mu(\xi)\right)(x),&\hspace{.4cm}\text{ in }\mathcal{D}'\bigl((0,T)\times\R_{x}^{d}\times\R_{z}^{n}\bigr).\notag
\end{alignat}
\end{subequations}
\end{theo}

The proof follows closely the analysis of \cite{Lions_Paul_Wigner};
the main difference being that here we have to control also what happens as $h\to 0$ for the wave part of the system \eqref{schroh}--\eqref{wave2_schroh}.
Note that if the sequence of initial data is supposed to converge, then, by uniqueness of the solution of the limit equation \cite[Th.\,4]{dBGV}, the entire sequence
$(W_{h}, \psi_{h},\chi_{h})_{h>0}$ converges.

\begin{proof}[Proof. Step 1: Convergence of $(\psi_{h})_{h>0}$]
Thanks to Remark~\ref{rmk_h5} we already know that the sequence $(\psi_{h})_{h>0}$ is bounded in $L^{\infty}(0,T;L_{x}^{2}\overbigdot{H}_{z}^{1})$. Since any closed ball of $L_{x}^{2}\overbigdot{H}_{z}^{1}$ is metrizable and compact for the weak topology, we are going to apply the Ascoli-Arzelà theorem in order to justify that $(\psi_{h})_{h>0}$ admits a converging sub-sequence in $C_{t}^{0}(L_{x}^{2}\overbigdot{H}_{z}^{1}-w)$. For that purpose it only remains to show that $(\psi_{h})_{h>0}$ is equi-continuous in $C_{t}^{0}(L_{x}^{2}\overbigdot{H}_{z}^{1}-w)$. In fact, it is sufficient to prove that the family $\{t\mto \langle \psi_{h}(t),g\rangle_{L_{x}^{2}\overbigdot{H}_{z}^{1}}\}$ is equi-continuous for every $g$ in a dense countable subset of $L_{x}^{2}\overbigdot{H}_{z}^{1}$. Details on this argument can be found \eg in \cite[App.\,C]{PLL}.
For any $g\in C_{c}^{\infty}(\R_{x}^{d}\times\R_{z}^{n})$,
\[
\left|\frac{\ud}{\ud t}\left\langle\psi_{h}(t),g\right\rangle_{L_{x}^{2}\overbigdot{H}_{z}^{1}}\right|=\left|\iint_{\R^{d}\times\R^{n}}\hat{\chi}_{h}(t,k,\zeta)|\zeta|^{2}\overline{\hat{g}(k,\zeta)}\ud k\ud\zeta\right|\leq\|\chi_{h}(t)\|_{L_{x}^{2}L_{z}^{2}}\|g\|_{L_{x}^{2}H_{z}^{2}}
\]
is uniformly bounded in $h$ and $t\in[0,T]$ (see Remark~\ref{rmk_h5}) and the Ascoli-Arzelà theorem insures us that, up to a sub-sequence, $(\psi_{h})_{h>0}$ converges in \hbox{$C^{0}([0,T];L_{x}^{2}\overbigdot{H}_{z}^{1}-w)$} to $\psi\in C^{0}([0,T];L_{x}^{2}\overbigdot{H}_{z}^{1}-w)$.

\subsubsection*{Step 2: Convergence of $(\chi_{h})_{h>0}$} As in the previous step Remark~\ref{rmk_h5} insures us that the sequence $(\chi_{h})_{h>0}$ is bounded in $L^{\infty}(0,T;L_{x}^{2}L_{z}^{2})$. Moreover, for any $g\in C_{c}^{\infty}(\R_{x}^{d}\times\R_{z}^{n})$,
\begin{align*}
\left|\frac{\ud}{\ud t}\left\langle\chi_{h}(t),g\right\rangle_{L_{x}^{2}L_{z}^{2}}\right|
&\leq c^{2}\left|\iint_{\R^{d}\times\R^{n}}\nabla_{z}\psi_{h}(t)\cdot\nabla_{z}g\ud x\ud z\right|\\
&\hspace*{2cm}+c^{2}\left|\iint_{\R^{d}\times\R^{n}}\sigma_{2}(z)\sigma_{1}\star|u_{h}(t)|^{2}(x)\,g(x,z)\ud x\ud z\right|\\
&\leq\|\psi_{h}\|_{L_{x}^{2}\overbigdot{H}_{z}^{1}}\|g\|_{L_{x}^{2}H_{z}^{1}}+\|\sigma_{1}\|_{L_{x}^{2}}\|\sigma_{2}\|_{L_{z}^{2}}\|u_{h}(t)\|_{L_{x}^{2}}^{2}\|g\|_{L_{x}^{2}L_{z}^{2}}
\end{align*}
is uniformly bounded in $h$ and $t\in[0,T]$ (see Remark~\ref{rmk_h5}). Eventually the Ascoli-Arzelà theorem insures us that, up to a sub-sequence, $(\chi_{h})$ converges in $C^{0}([0,T];L_{x}^{2}L_{z}^{2}-w)$ to $\chi\in C^{0}([0,T];L_{x}^{2}L_{z}^{2}-w)$.

\subsubsection*{Step 3: Equation on $\psi$} Since $\chi_{h}$ converges to $\chi$ in $C^{0}([0,T];L_{x}^{2}L_{z}^{2}-w)$ we obtain directly that for any $g\in C_{c}^{\infty}(\R_{x}^{d}\times\R_{z}^{n})$,
\[
\frac{\ud}{\ud t}\left\langle\psi_{h}(t),g\right\rangle_{\mathcal{D}',\mathcal{D}}=\iint_{\R^{d}\times\R^{n}}\chi_{h}(t)\,g\ud x\ud z\,\underset{h\to 0}{\to}\,\left\langle\chi(t),g\right\rangle_{\mathcal{D}',\mathcal{D}},
\]
the convergence being uniform on $[0,T]$. Note that here, since the duality product on $L_{x}^{2}\overbigdot{H}_{z}^{1}$ is not compatible with the duality product in $\mathcal{D}'$, we have to say something in order to justify the following convergence
\[
\frac{\ud}{\ud t}\left\langle\psi_{h}(t),g\right\rangle_{\mathcal{D}',\mathcal{D}}\,\underset{h\to 0}{\to}\,\frac{\ud}{\ud t}\left\langle\psi(t),g\right\rangle_{\mathcal{D}',\mathcal{D}}\quad\text{in }\mathcal{D}'(0,T).
\]
Since for any $f\in C_{c}^{\infty}(0,T)$,\vspace*{-3pt}
\[
\Bigl\langle\frac{\ud}{\ud t}\langle\psi_{h},g\rangle_{\mathcal{D}'},f\Bigr\rangle_{\mathcal{D'}(0,T)}=-\int_{0}^{T}\langle\psi_{h}(t),g\rangle_{\mathcal{D}'}f'(t)\ud t,
\]
we have to justify the uniform convergence in time of $\langle\psi_{h}(t),g\rangle_{\mathcal{D}'}$ to $\langle\psi(t),g\rangle_{\mathcal{D}'}$. For any $g\in C_{c}^{\infty}(\R_{x}^{d}\times\R_{z}^{n})$, we have
\[
\langle\psi_{h}(t),g\rangle_{\mathcal{D}'}=\iint_{\R^{d}\times\R^{n}}|\zeta|\hat{\psi}_{h}(t,k,\zeta)\,|\zeta|\overline{\frac{\hat{g}(k,\zeta)}{|\zeta|^{2}}}\ud k\ud\zeta.
\]
The condition $n\geq 3$ implies that $\mathcal{F}^{-1}(\hat{g}(k,\zeta)/|\zeta|^{2})$ lies in $L_{x}^{2}\overbigdot{H}_{z}^{1}$, and the convergence of $\psi_{h}$ to $\psi$ in $C^{0}([0,T];L_{x}^{2}\overbigdot{H}_{z}^{1}-w)$ allows us to conclude. Eventually we have proved that $\partial_{t}\psi=\chi$ in $\mathcal{D}'$.

\subsubsection*{Step 4: Equation on $\chi$} Let us temporarily assume that
$|u_{h}(t)|^{2}$ converges to a certain $\rho\in C^{0}([0,T];\mathcal{M}-w\star)$ (see Step 7).
For any $g\in C_{c}^{\infty}(\R_{x}^{d}\times\R_{z}^{n})$, we have
\begin{multline}\label{eq_on_chi}
\frac{\ud}{\ud t}\left\langle\chi_{h}(t),g\right\rangle_{\mathcal{D}',\mathcal{D}}=-c^{2}\iint_{\R^{d}\times\R^{n}}\hspace*{-1mm}\nabla_{z}\psi_{h}(t)\cdot\nabla_{z}g\ud x\ud z\\
-c^{2}\iint_{\R^{d}\times\R^{n}}\hspace*{-1mm}\sigma_{2}\,\sigma_{1}\star|u_{h}(t)|^{2}\,g\ud x\ud z.
\end{multline}
The weak convergence of $(\psi_{h})_{h>0}$ insures us that
\[
-c^{2}\iint_{\R^{d}\times\R^{n}}\nabla_{z}\psi_{h}(t)\cdot\nabla_{z}g\ud x\ud z\,\underset{h\to 0}{\to}\,-c^{2}\iint_{\R^{d}\times\R^{n}}\nabla_{z}\psi(t)\cdot\nabla_{z}g\ud x\ud z
\]
and, if we rewrite the second term of the right hand side of \eqref{eq_on_chi} as follows
\[
c^{2}\iint_{\R^{d}\times\R^{n}}\sigma_{2}\,\sigma_{1}\star|u_{h}(t)|^{2}\,g\ud x\ud z=c^{2}\int_{\R^{d}}|u_{h}(t,y)|^{2}\left(\int_{\R^{n}}\sigma_{2}\,\sigma_{1}\star g(y)\ud z\right)\ud y,
\]
the weak convergence of $|u_{h}|^{2}$ leads to
\[
c^{2}\iint_{\R^{d}\times\R^{n}}\sigma_{2}\,\sigma_{1}\star|u_{h}(t)|^{2}\,g\ud x\ud z\,\underset{h\to 0}{\to}\,c^{2}\iint_{\R^{d}\times\R^{n}}\sigma_{2}\,\sigma_{1}\star\rho(t)\,g\ud x\ud z.
\]
These two convergences hold uniformly in time and we eventually obtain
\[
\partial_{t}\chi=c^{2}\Delta_{z}\psi-c^{2}\sigma_{2}\,\sigma_{1}\star\rho(t)\quad\text{in }\mathcal{D}'\bigl((0,T)\times\R_{x}^{d}\times\R_{z}^{n}\bigr).
\]

\subsubsection*{Step 5: Convergence of $(W_{h})_{h>0}$} We first prove that the sequence $(W_{h})_{h>0}$ is bounded in $L^{\infty}\left(0,T;\mathcal{A}'\right)$. Since
\begin{multline*}
\iint_{\R^{d}\times\R^{d}}W_{h}(t,x,\xi)\varphi(x,\xi)\ud x\ud \xi\\
=\frac{1}{(2\pi)^{d}}\iint_{\R^{d}\times\R^{d}}u_{h}(t,x+\frac{h}{2}\,y)\bar{u}_{h}(t,x-\frac{h}{2}\,y)\mathcal{F}_{\xi}\varphi(x,y)\ud x\ud y,
\end{multline*}
we obtain directly
\begin{multline*}
\left|\iint_{\R^{d}\times\R^{d}}W_{h}(t,x,\xi)\varphi(x,\xi)\ud x\ud \xi\right|\\
\leq\frac{1}{(2\pi)^{d}}\left(\underset{y}{\sup}\,\int_{\R^{d}}\left|u_{h}(t,x+\frac{h}{2}\,y)\bar{u}_{h}(t,x-\frac{h}{2}\,y)\right|\ud x\right)\left(\underset{x}{\sup}\,\int_{\R^{d}}|\mathcal{F}_{\xi}\varphi(x,y)|\ud y\right)\\
\leq\frac{1}{(2\pi)^{d}}\|u_{h}(t)\|_{L_{x}^{2}}^{2}\|\varphi\|_{\mathcal{A}},
\end{multline*}
which insures us that
\[
\|W_{h}(t)\|_{\mathcal{A}'}\leq\frac{1}{(2\pi)^{d}}\|u_{h}(t)\|_{L_{x}^{2}}^{2}
\]
is bounded with respect to $h$ and $t$. Since any closed ball of $\mathcal{A}'$ is metrizable and compact for the weak-$\star$ topology, we will apply again the Ascoli-Arzelà theorem in order to justify that $(W_{h})_{h>0}$ admits a converging sub-sequence in $C_{t}^{0}(\mathcal{A}'-w\star)$. For that purpose we will prove that for any $\varphi\in\mathcal{B}$, the functions $t\mto \langle W_{h}(t),\varphi\rangle_{\mathcal{A}',\mathcal{A}}$ are equi-continuous. Direct computations yield
\begin{multline}\label{eq_on_Wh}
\frac{\ud}{\ud t}\left\langle W_{h}(t),\varphi\right\rangle_{\mathcal{A}',\mathcal{A}}=-\iint_{\R^{d}\times\R^{d}}W_{h}(t,x,\xi)\,\xi\cdot\nabla_{x}\varphi(x,\xi)\ud x\ud\xi\\
+\iint_{\R^{d}\times\R^{d}}W_{h}(t,x,\eta)\left(\int_{\R^{d}}K_{h}(t,x,\xi-\eta)\varphi(x,\xi)\ud\xi\right)\ud x\ud\eta,
\end{multline}
with
\begin{align*}
L_{h}(t,x,\eta)&:=\int_{\R^{d}}K_{h}(t,x,\xi-\eta)\varphi(x,\xi)\ud\xi\\
&\hphantom{:}=\frac{i}{(2\pi)^{d}}\int_{\R^{d}}e^{i\eta\cdot y}\frac{1}{h}\Bigl(\Phi_{h}\Bigl(t,x+\frac{h}{2}\,y\Bigr)-\Phi_{h}\Bigl(t,x-\frac{h}{2}\,y\Bigr)\Bigr)\mathcal{F}_{\xi}\varphi(x,y)\ud y
\end{align*}
and
\[
\mathcal{F}_{\eta}L_{h}(t,x,y)=\frac{i}{h}\Bigl(\Phi_{h}\Bigl(t,x+\frac{h}{2}\,y\Bigr)-\Phi_{h}\Bigl(t,x-\frac{h}{2}\,y\Bigr)\Bigr)\mathcal{F}_{\xi}\varphi(x,y).
\]
From \eqref{eq_on_Wh} we get for any $\varphi\in\mathcal{B}$,
\[
\Bigl|\frac{\ud}{\ud t}\left\langle W_{h}(t),\varphi\right\rangle_{\mathcal{A}',\mathcal{A}}\Bigr|\leq\|W_{h}(t)\|_{\mathcal{A}'}\left(\|\xi\cdot\nabla_{x}\varphi\|_{\mathcal{A}}+\|L_{h}(t)\|_{\mathcal{A}}\right)
\]
and it only remains to prove that $\mathcal{F}_{\eta}L_{h}(t)$ is bounded in $L_{y}^{1}C_{x}^{0}$, uniformly with respect to $t\in[0,T]$ and $h$. Since $\Phi_{h}=\sigma_{1}\star\int\sigma_{2}\psi_{h}\ud z$,
\[
\frac{1}{h}\Bigl(\Phi_{h}\Bigl(t,x+\frac{h}{2}\,y\Bigr)-\Phi_{h}\Bigl(t,x-\frac{h}{2}\,y\Bigr)\Bigr)=\frac{y}{h}\cdot\int_{-\sfrac{h}{2}}^{\sfrac{h}{2}}\nabla\sigma_{1}\star\left(\int_{\R^{n}}\sigma_{2}\psi_{h}(t)\ud z\right)(x+sy)\ud s
\]
and we can estimate $\mathcal{F}_{\eta}L_{h}(t)$ as follows
\begin{align*}
\|\mathcal{F}_{\eta}L_{h}(t)\|_{L_{y}^{1}C_{x}^{0}}&\leq\|y\mathcal{F}_{\xi}\varphi\|_{L_{y}^{1}C_{x}^{0}}\biggl\|\frac{1}{h}\int_{-\sfrac{h}{2}}^{\sfrac{h}{2}}\nabla\sigma_{1}\star\left(\int_{\R^{n}}\sigma_{2}\psi_{h}(t)\ud z\right)(x+sy)\ud s\biggr\|_{L_{x,y}^{\infty}}\\
&\leq\|y\mathcal{F}_{\xi}\varphi\|_{L_{y}^{1}C_{x}^{0}}\biggl\|\nabla\sigma_{1}\star\left(\int_{\R^{n}}\sigma_{2}\psi_{h}(t)\ud z\right)\biggr\|_{L_{x}^{\infty}}.
\end{align*}
The following estimate coupled with \eqref{estimate_strichartz_3} and Remark~\ref{rmk_h5} allows us to conclude
\[
\left\|\nabla\sigma_{1}\star\left(\int_{\R^{n}}\sigma_{2}\psi_{h}(t)\ud z\right)\right\|_{L_{x}^{\infty}}\leq\|\nabla\sigma_{1}\|_{L_{x}^{2}}\|\sigma_{2}\|_{L_{z}^{p_{0}'}}\|\psi_{h}\|_{L_{x}^{2}L_{t}^{\infty}L_{z}^{p_{0}}}.
\]

\subsubsection*{Step 6: Equation on $\mu$} For any $\varphi\in\mathcal{B}$, we have
\[
\frac{\ud}{\ud t}\langle W_{h}(t),\varphi\rangle_{\mathcal{B}',\mathcal{B}}=-\langle W_{h}(t),\xi\cdot\nabla_{x}\varphi\rangle_{\mathcal{B}',\mathcal{B}}+\langle W_{h}(t),L_{h}(t)\rangle_{\mathcal{B}',\mathcal{B}}.
\]
The weak convergence of $(W_{h})_{h>0}$ allows us to obtain
\[
\frac{\ud}{\ud t}\langle W_{h}(t),\varphi\rangle_{\mathcal{B}',\mathcal{B}}\,\underset{h\to 0}{\to}\,\frac{\ud}{\ud t}\langle \mu(t),\varphi\rangle_{\mathcal{B}',\mathcal{B}} \hspace{0.5cm}\text{ in }\mathcal{D}'(0,T),
\]
and
\[
\langle W_{h}(t),\xi\cdot\nabla_{x}\varphi\rangle_{\mathcal{B}',\mathcal{B}}\,\underset{h\to 0}{\to}\,\langle \mu(t),\xi\cdot\nabla_{x}\varphi\rangle_{\mathcal{B}',\mathcal{B}} \hspace{0.5cm}\text{ uniformly in time }(t\in[0,T]),
\]
and it only remains to prove that $L_{h}(t)$ converges strongly in $\mathcal{A}$ (uniformly with respect to $t\in[0,T]$) to $\nabla_{x}\left(\sigma_{1}\star\int\sigma_{2}\psi(t)\ud z\right)\cdot\nabla_{\xi}\varphi$, which is equivalent to prove the strong convergence of $\mathcal{F}_{\xi}L_{h}(t)$ to $iy\cdot(\nabla\sigma_{1}\star\int\sigma_{2}\psi(t)\ud z)\mathcal{F}_{\xi}\varphi$ in $L_{y}^{1}C_{x}^{0}$. For that purpose we decompose the difference of these two terms as follows
\bgroup
\multlinegap0pt
\begin{multline*}
\mathcal{F}_{\xi}L_{h}(t,x,y)-iy\cdot\left(\int_{\R^{d}}\nabla\sigma_{1}(x-\bar{x})\left[\int\sigma_{2}(z)\psi(t,\bar{x},z)\ud z\right]\ud\bar{x}\right)\,\mathcal{F}_{\xi}\varphi(x,y)\\
=iy\cdot\left(\int_{\R^{d}}\nabla\sigma_{1}(x-\bar{x})\left[\int\sigma_{2}(z)(\psi(t,\bar{x},z)-\psi_{h}(t,\bar{x},z))\ud z\right]\ud\bar{x}\right)\,\mathcal{F}_{\xi}\varphi(x,y)\\
+iy\cdot\biggl(\int_{\R^{d}}\frac{1}{h}\biggl[\int_{-\sfrac{h}{2}}^{\sfrac{h}{2}}\nabla\sigma_{1}(x-\bar{x})-\nabla\sigma_{1}(x+sy-\bar{x})\ud s\biggr]\biggl[\int\sigma_{2}(z)\psi_{h}(t,\bar{x},z)\ud z\biggr]\ud\bar{x}\biggr)\\
\shoveright{\times \mathcal{F}_{\xi}\varphi(x,y)}\\
=\mathrm{I}(t,x,y)+\mathrm{II}(t,x,y).
\end{multline*}
\egroup
We estimate the first term as follows (where the support of $\mathcal{F}_{\xi}\varphi$ is supposed to be included in the compact $K_{1}\times K_{2}$)
\[
\|\mathrm{I}(t)\|_{L_{y}^{1}C_{x}^{0}}\leq\|y\mathcal{F}_{\xi}\varphi\|_{L_{y}^{1}C_{x}^{0}}\;\underset{x\in K_{1}}{\sup}\,\left|\nabla\sigma_{1}\star\left(\sigma_{2}(\psi(t)-\psi_{h}(t))\right)(x)\right|
\]
and the weak convergence of $(\psi_{h})_{h>0}$ insures us that for every $x\in K_1$
\begin{multline*}
\nabla\sigma_{1}\star\left(\sigma_{2}(\psi(t)-\psi_{h}(t))\right)(x)\\
=\iint_{\R^{d}\times\R^{d}}|\zeta|\,\nabla\sigma_{1}(x-\bar{x})\frac{\hat{\sigma}_{2}(\zeta)}{|\zeta|^{2}}\;|\zeta|\,\overline{\left(\hat{\psi}(t,\bar{x},\zeta)-\hat{\psi}_{h}(t,\bar{x},\zeta)\right)}
\ud \bar x \ud\zeta\underset{h\to 0}{\to}\,0.
\end{multline*}
This convergence is not a priori uniform in $x\in K_1$. Nevertheless, we can combine the fact that $\psi(t)-\psi_{h}(t)$ is uniformly bounded with respect to $t$ and $h$ in $L_{x}^{2}\overbigdot{H}_{z}^{1}$, $K_1$ is compact and the application
\[
x\in\mathbb R^d \longmapsto\,\left((\bar{x},z)\mapsto\nabla\sigma_{1}(x-\bar{x})\mathcal{F}_{\zeta}^{-1}(\hat{\sigma}_{2}(\zeta)/|\zeta|^{2})(z)\right)\,\in L_{x}^{2}\overbigdot{H}_{z}^{1}
\]
is continuous,
to prove that the convergence is indeed uniform in $x$. For the second term, the estimate\enlargethispage{-\baselineskip}
\begin{align*}
\|\mathrm{II}(t)\|_{L_{y}^{1}C_{x}^{0}}&\leq \|y\mathcal{F}_{\xi}\varphi\|_{L_{y}^{1}C_{x}^{0}}\,\|\sigma_{2}\|_{L_{z}^{p_{0}'}}\,\|\psi_{h}\|_{L_{x}^{2}L_{t}^{\infty}L_{z}^{p_{0}}}\\
&\hspace*{1.5cm}\times\underset{\substack{x\in K_{1} \\ y\in K_{2}}}{\sup}\biggl(\int_{\R^{d}}\frac{1}{h^{2}}\biggl|\int_{-\sfrac{h}{2}}^{\sfrac{h}{2}}\nabla\sigma_{1}(x-\bar{x})-\nabla\sigma_{1}(x+sy-\bar{x})\ud s\biggr|^{2}\biggr)^{1/2}\\
&=\|y\mathcal{F}_{\xi}\varphi\|_{L_{y}^{1}C_{x}^{0}}\,\|\sigma_{2}\|_{L_{z}^{p_{0}'}}\,\|\psi_{h}\|_{L_{x}^{2}L_{t}^{\infty}L_{z}^{p_{0}}}\\
&\hspace*{1.5cm}\times \underset{y\in K_{2}}{\sup}\biggl(\int_{\R^{d}}\frac{1}{h^{2}}\biggl|\int_{-\sfrac{h}{2}}^{\sfrac{h}{2}}\nabla\sigma_{1}(x)-\nabla\sigma_{1}(x+sy)\ud s\biggr|^{2}\ud x\biggr)^{1/2}
\end{align*}
coupled with the regularity and the compactness of the support of $\nabla\sigma_{1}$ and the uniform boundedness with respect to $h$ of $\|\psi_{h}\|_{L_{x}^{2}L_{t}^{\infty}L_{z}^{p_{0}}}$, allows us to conclude that $\|\mathrm{II}(t)\|_{L_{y}^{1}C_{x}^{0}}\to 0$ when $h\to 0$.

\subsubsection*{Step 7: Final details} To conclude the proof it remains to justify that in fact the limit $\mu$ of the sequence $(W_{h})_{h>0}$ defines an element of $C^{0}([0,T],\mathcal{M}_+-w\star)$ and that the sequence $(|u_{h}|^{2})_{h>0}$ converges in $C^{0}([0,T],\mathcal{M}(\R_{x}^{d})-w\star)$ to $\rho=\int\ud\mu(\xi)$. The first point comes from the study of the Husimi transform of $u_{h}$:
\[
\widetilde{W}_{h}(t)=W_{h}(t)\star\frac{e^{-(|x|^{2}+|\xi|^{2})/h}}{(\pi h)^{d}}.
\]
One can prove that, for every time $t\in[0,T]$, $\widetilde{W}_{h}(t)$ is non negative and the \hbox{sequence} $(\widetilde{W}_{h}(t))_{h>0}$ is bounded in $L_{x}^{1}L_{\xi}^{1}$. This allows us to conclude that, up to a sub-sequence, $\widetilde{W}_{h}(t)$ converges weakly in the sense of measures to a certain $\tilde{\mu}(t)\in\mathcal{M}_+$ and it is then possible to prove that indeed $\mu(t)=\tilde{\mu}(t)$. We refer the reader to \cite[\S III]{Lions_Paul_Wigner} for details. However it is not possible yet to conclude that $\mu$ is an element of $C^{0}([0,T],\mathcal{M}-w\star)$. In~the previous argument each sub-sequence depends on $t$ (then it is not possible to apply a diagonal argument) and we have no information about the time continuity. The missing step can be obtained by slightly modifying the compactness argument in Step~5, in order to obtain the compactness of the sequence $(\widetilde{W}_{h})_{h>0}$ in \hbox{$C^{0}([0,T],\mathcal{M}-w\star)$}, and conclude that, up to a sub-sequence, $(\widetilde{W}_{h})_{h>0}$ converges in \hbox{$C^{0}([0,T],\mathcal{M}-w\star)$} to $\tilde{\mu}\in C^{0}([0,T],\mathcal{M}-w\star)$. We eventually obtain that $\mu=\tilde{\mu}\in C^{0}([0,T],\mathcal{M}-w\star)$.

Finally, we make use of the results in the \cite[\S III ]{Lions_Paul_Wigner} which tell us that
if the sequence $(h^{-d}|\hat{u}_{h}(t,h^{-1}\xi)|^{2})_{h>0}$ is tightly relatively compact, then $(|u_{h}(t)|^{2})$ converges weakly in the sense of measures to $\rho(t)=\int\ud\tilde{\mu}(t,\xi)=\int\ud\mu(t,\xi)$.
Moreover, we already know that $(\widetilde{W}_{h})_{h>0}$ converges in $C^{0}([0,T],\mathcal{M}-w\star)$ to $\tilde{\mu}$, so that if $(h^{-d}|\hat{u}_{h}(t,h^{-1}\xi)|^{2})_{h>0}$ is tightly relatively compact, \emph{uniformly in time},
then the proof \cite[Th.\,III.1(3)]{Lions_Paul_Wigner} can be revisited in order to obtain that $(|u_{h}|^{2})_{h>0}$ converges in $C^{0}([0,T],\mathcal{M}(\R^{d})-w\star)$ to $\rho=\int\ud\tilde{\mu}(\xi)=\int\ud\mu(\xi)\in C^{0}([0,T],\mathcal{M}(\R^{d})-w\star)$.

Let us conclude the proof by proving that the sequence $(h^{-d}|\hat{u}_{h}(t,h^{-1}\xi|^{2})_{h>0}$ is tightly relatively compact uniformly in time, which can be cast as
\[
\underset{t\geq 0}{\sup}\;\underset{h>0}{\sup}\;\frac{1}{h^{d}}\int_{|\xi|\geq R}\left|\hat{u}_{h}(t,h^{-1}\xi)\right|^{2}\ud\xi\;\underset{R\to \infty}{\to}\;0.
\]
Remark~\ref{rmk_h5} insures the existence of a constant $C>0$, independent of $h>0$ and $t\in[0,T]$, such that $h^{2}\|\nabla_{x}u_{h}(t)\|_{L_{x}^{2}}^{2}\leq C$. Then a direct computation shows that
\begin{multline*}
h^{2}\int_{\R^{d}}|\nabla_{x}u_{h}(t,x)|^{2}\ud x=h^{2}\int_{\R^{d}}|\xi|^{2}\left|\hat{u}_{h}(t,\xi)\right|^{2}\ud\xi\\
=\frac{1}{h^{d}}\int_{\R^{d}}|\xi|^{2}\left|\hat{u}_{h}(t,h^{-1}\xi)\right|^{2}\ud\xi\geq\frac{1}{h^{d}}\int_{|\xi|\geq R}R^{2}\left|\hat{u}_{h}(t,h^{-1}\xi)\right|^{2}\ud\xi,
\end{multline*}
and we eventually obtain
\[
\underset{t\geq 0}{\sup}\;\underset{h>0}{\sup}\;\frac{1}{h^{d}}\int_{|\xi|\geq R}\left|\hat{u}_{h}(t,h^{-1}\xi)\right|^{2}\ud\xi\leq \frac{C}{R^{2}}.\qedhere
\]
\end{proof}

\section{Scaling of the model}\label{adim}

Let us provide some details about the scaling of the equations
explaining the meaning of the parameters we are dealing with in this paper. Working with physical units, the system reads
\begin{subequations}
\begin{align}
\label{schro_niv0}
\Bigl(i\hbar \partial_{t}u+\frac{\hbar ^2}{2m}\Delta_{x}u\Bigr)(t,x)&=\left(\int_{\mathbb R^d\times\mathbb R^n}
\sigma_{1}(x-y) \sigma_{2}(z)\psi(t,y,z)\ud y\ud z\right)u(t,x),\\
\label{wave_niv0}(\partial_{tt}^{2}\psi-\varkappa^{2}\Delta_{z}\psi)(t,x,z)&=-\sigma_{2}(z)\biggl(\int_{\mathbb R^d} \sigma_{1}(x-y)|u(t,y)|^{2}
\ud y\biggr).
\end{align}
\end{subequations}
In order to bring out the parameters of the model we need to introduce relevant units of observations: $\textup M$, $\textup L$ and $\textup T$
stand for these mass, length and time units, respectively.
Here $\hbar $ is the Planck constant; its homogeneity is $\sfrac{\textup{Mass}\times \textup{Length}^2}{ \textup{Time}}$ (and its value is $1,055 \times10^{-34}\ J s$); $m$ is the inertial mass of the considered particle and $\psi$ is the wave function, such that $\int_\Omega |\psi|^2\ud x$ is
the probability of presence of the quantum particle in $\Omega\subset \mathbb R^d$ (dimensionless).
It helps the intuition to think of the $z$ directions as homogeneous to a length, but in fact this is not necessarily the case: we denote by $\Psi$ and $\mathrm Z$ the (unspecified) units for $\psi$ and the $z_j$'s. Accordingly, $\varkappa$ has the same dimension as $\sfrac{\mathrm Z}{\mathrm T}$.
Finally, we remark that what really matters to make the coupling consistent is the unit of the product of the form functions $\sigma_1\sigma_2$, that will be denoted~$\Sigma$. It turns out that $\Sigma $ has the same homogeneity
as
$\sfrac{\hbar }{ \mathrm T \Psi\mathrm L^d\mathrm Z^n}$
from \eqref{schro_niv0}
and
as $\sfrac{ \Psi}{\mathrm L^d \mathrm T^2}$
from \eqref{wave_niv0}, both are thus measured with the same units.

\pagebreak[2]
Besides, system \eqref{schro_niv0}--\eqref{wave_niv0} preserves the total energy $\mathscr E_0$ which is the given by the sum of
\begin{itemize}
\item $\ds\frac{\hbar ^2}{2m}\int_{\mathbb R^d}|\nabla_x u |^2 \ud x$, which has the homogeneity of $\sfrac{\hbar ^2}{m \mathrm L^2}$,
\item $\ds\frac{1}{2}\iint_{\mathbb R^d\times \mathbb R^n}\big(|\partial_t \psi |^2 + \varkappa^2|\nabla_z\psi|^2\big) \ud z\ud x$
which has the homogeneity of $\sfrac{\Psi^2 \mathrm L^d\mathrm Z^n}{ \mathrm T^2}$,
\item $\ds\iint_{\mathbb R^d\times \mathbb R^n}
|u|^2 \sigma_2\sigma_1\star \psi \ud z\ud x$
which has the homogeneity of $\Sigma \Psi \mathrm L^d\mathrm Z^n$.
\end{itemize}
All these quantities thus scale like $\sfrac{ \mathrm M \mathrm L^2}{\ \mathrm T^2}$.
From now on, we
assume that the total amount of energy $\mathscr E_0$ is given, so that, given the time and length units $\mathrm T, \mathrm L$, it defines
the mass unit
\[
\mathrm M= \ds\frac{\mathrm T^2}{\mathrm L^2\mathscr E_0}
\]
which has to be compared to the particle's mass $m$.

Based on this discussion, we introduce dimensionless quantities defined as follows
\begin{gather*}
u'(t',x')=u(t'\mathrm T, x'\mathrm L) \ \sqrt{\mathrm L^{d}\ds\frac{m}{\mathrm M}},
\qquad
\psi'(t',x',z')=\ds\frac{1}{\Psi} \psi(t'\mathrm T, x'\mathrm L,z'\mathrm Z),
\\
\sigma_1'(x')\sigma_2(z')=\ds\frac1\Sigma\ \sigma_1(x'\mathrm L)\sigma_2(z'\mathrm Z).
\end{gather*}
Note that with these definitions, we have
\[
\int_{\mathbb R^d} |u'(t',x')|^2\ud x'=\ds\frac{m}{\mathrm M}.
\]
To reduce the amount of notation, we drop the primes and we rewrite the system in the following dimensionless form
\begin{subequations}
\begin{alignat}{11}
\label{schro_res}&\ds \Bigl(i\partial_{t}u+\frac{\hbar \mathrm T}{m\mathrm L^2}\frac12\Delta_{x}u\Bigr)(t,x)
\\\nonumber
&\hspace*{2cm}=
\ds\frac{\Psi \Sigma \mathrm L^d\mathrm Z^n \mathrm T}{\hbar}
\left(\ds\int_{\mathbb R^d\times\mathbb R^n}
\sigma_{1}(x-y) \sigma_{2}(z)\psi(t,y,z)\ud y\ud z\right)u(t,x),\\
\label{wave_res}&\Bigl(\partial_{tt}^{2}\psi-\frac{\varkappa^{2}\mathrm T^2}{\mathrm Z^2}
\Delta_{z}\psi\Bigr)(t,x,z)=-
\ds\frac{ \Sigma\mathrm T^2}{\Psi}\ds\frac {\mathrm M}{m}
\sigma_{2}(z)\biggl(\ds\int_{\mathbb R^d} \sigma_{1}(x-y)|u(t,y)|^{2}
\ud y\biggr).
\end{alignat}
\end{subequations}
The energy conservation takes the form
\begin{multline*}
\frac{\hbar ^2}{m\mathrm L^2}\ds\frac{\mathrm M}{m}\frac12\ds\int_{\mathbb R^d}|\nabla_x u |^2 \ud x
+
\frac{\Psi^2\mathrm L^d\mathrm Z^n}{\mathrm T^2}\ds\frac12\iint_{\mathbb R^d\times \mathbb R^n}
\Big(|\partial_t \psi|^2 +\ds\frac{\varkappa^2\mathrm T^2}{\mathrm Z^2} |\nabla_z\psi|^2
\Big)
\ud z\ud x
\\
+
\Sigma \Psi \mathrm L^d\mathrm Z^n\ds\frac {\mathrm M}m
\iint_{\mathbb R^d\times \mathbb R^n}
|u|^2 \sigma_2\sigma_1\star \psi \ud z\ud x
=\mathscr E_0=\frac{\mathrm M\mathrm L^2}{\mathrm T^2}.
\end{multline*}
The coefficients of this system are dimensionless; we obtain \eqref{schro}--\eqref{wave_schro} by assuming
\[
\frac{\hbar \mathrm T}{m\mathrm L^2}=1= \ds\frac{\Psi \Sigma \mathrm L^d\mathrm Z^n \mathrm T}{\hbar}
,\qquad
c^2=\frac{\varkappa^{2}\mathrm T^2}{\mathrm Z^2}=\ds\frac{ \Sigma\mathrm T^2}{\Psi}\ds\frac {\mathrm M}{m}.
\]
By using the identities
\[
\Bigl(\ds\frac{\hbar\mathrm T}{m\mathrm L^2}\Bigr)^2=1,\quad
\ds\frac{ \Sigma \Psi \mathrm L^d\mathrm Z^n \mathrm T^2}{\hbar}\times \ds\frac {\hbar\mathrm T}{m\mathrm L^2}
=1,\quad
\ds\frac{\Sigma\Psi\mathrm L^d\mathrm Z^n\mathrm T}{\hbar }
\times\ds\frac{\Psi}{\Sigma \mathrm T^2}\times\ds\frac m{\mathrm M} \times \ds\frac{\hbar \mathrm T}{m\mathrm L^2}
=\ds\frac{1}{c^2},
\]
we check that the energy
recasts as announced in \eqref{energy_Sch}.
Of course, they are multiple ways to rephrase
this setting. As a matter of fact, the choice of units and the parameters $c$ and $\sfrac{m}{\mathrm M}$ can be
reinterpreted by means of ratio of the energies involved in the model
\begin{itemize}
\item quantum energy to total energy $\sfrac{\hbar^2}{m\mathrm L^2}
=\Psfrac{m}{\mathrm M}\mathscr E_0$,
\item vibrational energy to total energy $\sfrac{\Psi^2\mathrm L^d\mathrm Z^n}{\mathrm T^2} =
\Psfrac{1}{c^2}\mathscr E_0$,
\item coupling energy to the total energy $\Sigma \Psi \mathrm L^d\mathrm Z^n=\Psfrac{m}{\mathrm M}\mathscr E_0$.
\end{itemize}
